\documentclass[11pt,a4paper]{article}

\usepackage[T1]{fontenc}
\usepackage[utf8]{inputenc}
\usepackage[british]{babel}
\usepackage{lmodern}
\usepackage{microtype}
\usepackage[a4paper,margin=1in]{geometry}
\usepackage{amsmath,amssymb,amsthm,mathtools}
\usepackage{enumitem}
\usepackage{booktabs,tabularx,array}
\usepackage{xcolor}
\usepackage{hyperref}
\hypersetup{
  colorlinks=true,
  linkcolor=blue!60!black,
  citecolor=blue!60!black,
  urlcolor=blue!60!black,
  hypertexnames=false,
  pdftitle={A Preference for Traceable Agency: An Axiomatization of Expected Utility plus Mutual Information},
  pdfauthor={Daniele Condorelli},
  pdfsubject={A behavioural characterisation of expected utility plus mutual information}
}
\setlength{\emergencystretch}{2em}
\interfootnotelinepenalty=10000

\newtheorem{theorem}{Theorem}
\newtheorem{proposition}[theorem]{Proposition}
\newtheorem{lemma}[theorem]{Lemma}
\newtheorem{corollary}[theorem]{Corollary}
\theoremstyle{definition}
\newtheorem{definition}[theorem]{Definition}
\newtheorem{remark}[theorem]{Remark}

\newcommand{\D}{\Delta}
\newcommand{\E}{\mathbb E}
\newcommand{\I}{I}
\newcommand{\Sh}{H}
\newcommand{\supp}{\operatorname{supp}}
\newcommand{\Id}{\operatorname{Id}}
\newcommand{\Rnum}{\mathbb R}
\newcommand{\R}{\mathbb R}
\newcommand{\KL}{D_{\mathrm{KL}}}
\newcommand{\1}{\mathbf 1}

\title{A Preference for Traceable Agency\\[0.35em]
\large An Axiomatization of Expected Utility plus Mutual Information\thanks{Email:
\href{mailto:d.condorelli@gmail.com}{d.condorelli@gmail.com}.  I thank the LSE Theory Group for hospitality and stimulating discussion during the summer of 2025, when this project was first conceived. I thank Peter Hammond and Agustin Troccoli Moretti for comments. Joe Basford suggested a way to simplify Axiom~(A8). 

\smallskip
All parts of this manuscript have been written with assistance from a multi-model LLM workflow.  The
question, the model, the axioms, and the form of the representation are the work of human hands.  The author
personally contributed very little to the writing of the proofs beyond advice and encouragement.  A
machine-checked Lean~4 proof of Theorem~\ref{thm:main} is available at
\url{https://github.com/dcond/TraceableAgency}: the verified declaration is
\href{https://github.com/dcond/TraceableAgency/blob/main/TraceableAgency/Theorem1/Proof.lean}{\texttt{trace\_tempered\_choice\_v10\_theorem1}},
and the \href{https://github.com/dcond/TraceableAgency/blob/main/Certificate/FORMALIZATION_CERTIFICATE.md}{formalisation certificate}
maps the axioms and theorem to their formal statements.  All claims, errors, and omissions remain the sole responsibility of the author.}}
\author{Daniele Condorelli\\University of Warwick}
\date{\today}

\begin{document}
\maketitle

\begin{abstract}
\noindent
An individual may value her actions both for their payoffs and for what their observable consequences reveal
about which one she took.  A channel maps actions to distributions over consequences; in each channel she ranks
lotteries over actions.  Two axioms are distinctive: she never prefers her action or its consequence garbled, and
a sure-thing principle holds with interim records as events.  Together with standard ones, they characterise
expected payoff utility plus a positive, channel-independent weight on the mutual information between action and
consequence, her trace.  The consequences of a predetermined act reveal nothing new, so randomising can be
strictly optimal.
\end{abstract}

\medskip
\noindent\textit{Keywords:} traceable agency; mutual information; Shannon entropy; preference for
randomisation; intrinsic preference for information; non-expected utility.\\
\textit{JEL codes:} D81, D83, D91.

\clearpage

% ===============================================================
\section{Introduction}
% ===============================================================


\begin{quote}
\emph{Shower upon him every earthly blessing \ldots\ He would even risk his cakes and would deliberately desire the
most fatal rubbish, the most uneconomical absurdity \ldots\ simply in order to prove to himself---as though that were
so necessary---that men still are men and not the keys of a piano.\ \ldots\ If you say that all this, too, can be
calculated and tabulated \ldots\ then man would purposely go mad in order to be rid of reason and gain his point!}

\hfill --- Fyodor Dostoevsky, \emph{Notes from the Underground}, Part I, Chapter VIII \cite{Dostoevsky1864}
\end{quote}


Dostoevsky's Underground Man faces a choice problem.  The cakes are material payoff, and he will risk them so
that his conduct can attest to being his own, not the mechanical implication of a table.  The proof he seeks
is self-defeating.  An act fixed in advance cannot provide it, and deliberate rebellion can be
“calculated and tabulated.”  So he would “purposely go mad in order to be rid of reason”---not because
madness is what he values, but because it is his last imagined way of preventing the table from closing before he
acts.  This paper reads madness as making oneself unpredictable, and takes the evidentiary idea literally.  In
Dostoevsky's example conduct is observed directly; we consider the general case, in which an act is observed
only through its visible consequences.

For visible consequences to carry information about an act, two conditions are needed.  There must be uncertainty
beforehand about which act will occur, and the consequence must discriminate among the
possible acts.  Either condition can fail on its own.  A lottery over acts leaves no trace when every act generates
the same distribution of visible consequences; a consequence that perfectly reveals the act conveys
nothing new when the act was already known.  We call the intrinsic preference for the two together a preference for \emph{traceable agency}.  The paper axiomatises it jointly with ordinary
preferences over payoff risk.

Let $A$ be a finite set of actions, $O$ a finite set of payoff-bearing
outcomes, $R$ a finite set of explicit visible records, and $K:A\to\Delta(O\times R)$ the channel relating actions to
visible consequences.  An action lottery $q\in\Delta(A)$ gives the distribution of the realised action.  The primitive is a family
$\{\succeq_K\}_K$, with one binary relation $\succeq_K$ on $\Delta(A)$ for each fixed channel $K$.
The statement $q\succeq_Kp$ means that $q$ is weakly preferred to $p$ in $K$.  Eight axioms are
imposed on this family.  To compare lotteries governed
by different channels, the channels are placed side by side as blocks of one larger environment; duplicating a
channel or adding an unused block does not alter a comparison.  Two processing axioms say that the agent never gains when her record is garbled or her action variable is processed noisily, and a benchmark comparison rules out trace indifference. A recordwise sure-thing principle for compounded records supplies
the cardinal structure.  The taste prices how much the consequence reveals, not what it reveals.

A family $\{\succeq_K\}_K$ satisfies the eight axioms if and
only if there are a nonconstant payoff utility $u:O\to\mathbb R$ and $\lambda>0$ such that, for every finite joint
channel $K$ and every $q,p\in\Delta(A)$,
\[
 q\succeq_Kp
 \quad\Longleftrightarrow\quad
 \E_{qK}[u(O)]+\lambda I_{qK}(A;O,R)
 \ge
 \E_{pK}[u(O)]+\lambda I_{pK}(A;O,R),
\]
where $I_{qK}(A;O,R)$ is the mutual information between the action and its visible consequence.  We call such
families \emph{trace-tempered}.  The utility $u$ is unique up to positive affine transformations; after fixing that
scale and the logarithm base, $\lambda$ is measured in payoff units per unit of information, and the same pair prices comparisons across channels.  The sign is set by three axioms: reversing them characterises $\lambda<0$, a preference for statistical
deniability, and replacing them by indifference clauses yields expected utility, $\lambda=0$. 

While there is no observer in the model, preferences admit an instrumental reading: an audience rewarding the agent in proportion to what the visible consequence lets it learn about the act induces the criterion. Rather than a taste, $\lambda$ becomes a price of attribution, positive for a politician backed when acts move outcomes, negative if what is learned is used against~her.


The information term has a simple interpretation.  Observing $(O,R)$ induces a posterior belief about the realised
action; traceability is the expected reduction in that uncertainty, measured by Shannon entropy $\Sh$.  If visible
consequences are independent of actions, the value is zero.  If the action is known in advance because
$q$ is degenerate, the value is also zero.  If the visible-consequence supports of the actions in the support of
$q$ are pairwise disjoint, the value is $\Sh(q)$---Dostoevsky's case, in which conduct is observed directly.\footnote{If the
visible-consequence supports of all channel rows are pairwise disjoint, the unique optimal lottery has
multinomial-logit weights.  Overlap can make independence of irrelevant alternatives fail; see
\S\ref{sec:randomisation}.}  Thus the criterion is
not a taste for randomisation as such.  Randomisation matters only because it creates uncertainty about the act for the
visible consequence to resolve.

The taste is for the coupling of act and consequence.  A
four-action channel turns this into a direct test.  Every action pays the same amount and writes one of two
symbols on the record.  Action $a_1$ always writes the first symbol and $a_2$ the second; $a_3$ and $a_4$ each
write either symbol with equal probability.  The choice is between a fair coin over $\{a_1,a_2\}$ and a fair
coin over $\{a_3,a_4\}$.  Both coins are fair and induce the same joint distribution of payoff
and record.  Under the first coin the record identifies the realised action; under the second it reveals
nothing.  Any criterion that reads only the distribution of consequences is indifferent between the coins; so is
a taste for randomisation as such.  Traceable agency predicts a strict preference for the first.  Because the
record affects neither payoffs nor any later decision, the preference cannot be
instrumental.\footnote{A second treatment completes the design: repeat
the choice in a channel that writes a fair record independently of the realised action.  Every representation
then assigns both flips the same value, whatever $u$ and $\lambda$, so a strict preference there rejects the
model.  Section~\ref{sec:implications} develops the test and what it excludes.}  The comparison identifies a
value of act--consequence dependence; by itself it does not distinguish trace from control or identify the
psychological source of that value.\footnote{The information term admits a control interpretation. Read backward: it measures trace. Read forward, it measures how strongly the action controls the distribution of visible consequences. We return to this in Section~\ref{sec:interpretation}.}


This comparison has not been run.  The nearest evidence comes from two experiments.  Bartling, Fehr, and Herz elicit each principal's point of indifference between retaining and delegating the right to choose a project and the effort devoted to it.  At that point, the delegation lottery has a certainty equivalent 16.7\% higher on average than the lottery produced by the principal's own choices, when both are valued outside the delegation setting \cite{BartlingFehrHerz2014}.  Principals pay to keep the payoff lottery dependent on their own choices, which a trace reading accommodates only if that choice is itself uncertain. The same evidence may instead reflect autonomy.  Mistry and Liljeholm, and later Norton and Liljeholm, ask subjects to select a room containing two slot machines and then gamble in that room.  Subjects are more likely to choose a high-divergence room over a zero-divergence room when they choose between its machines themselves than when a computer chooses for them, holding current expected payoffs and outcome entropies fixed \cite{MistryLiljeholm2016,NortonLiljeholm2020}.  With equal weight on the two machines, the divergence between their token distributions is the mutual information between the machine and the token.  This also admits a trace interpretation: in self-play, the token is evidence about a voluntarily selected machine.  But token values could change after the room was selected, allowing subjects in self-play to adapt their machine choices.  Thus trace is confounded with the value of deciding in the first design and with that of adapting in the second.  The test proposed above removes both confounds.

Section~\ref{sec:model} sets out the domain and the eight axioms.  Section~\ref{sec:theorem} states the representation theorem,
with a proof sketch, and derives uniqueness, the sign trichotomy, and the
independence of the axioms.
Section~\ref{sec:implications} presents a two-treatment test that can reject the model; derives deliberate randomisation, with multinomial logit as the boundary case when channel-row supports are pairwise disjoint; and measures $\lambda$ from choices.
Section~\ref{sec:literature} reviews the literature, in particular rational inattention, where mutual information is a cost of attention rather than an object of
taste, and the axiomatic characterisations of information indices.
Section~\ref{sec:interpretation} discusses interpretation and scope.  The proofs are in the appendices:
Appendix~\ref{app:main-proof} gives a proof of Theorem~1, and Appendix~\ref{app:auxiliary-results} contains the remaining results.


% ===============================================================
\section{The model}\label{sec:model}
% ===============================================================

Fix a finite set $O$ of \emph{payoff-bearing outcomes}, with $|O|\ge2$.  Let $A$ be a nonempty finite set of
\emph{actions} and $R$ a nonempty finite set of \emph{explicit visible records}.  A joint channel is a stochastic
kernel
\[
  K:A\longrightarrow\D(O\times R).
\]
For each finite joint channel $K$, the primitive behavioural datum is a binary relation
\[
 \succeq_K\quad\text{on}\quad\D(A).
\]
The interpretation is that $q$ is a lottery over actions used by the agent and $q\succeq_Kp$ means that, in the fixed
environment $K$, $q$ is evaluated at least as highly as $p$.  Within that environment, the lottery is the
procedure being evaluated and the channel is fixed.\footnote{Why not a single
preference over joint laws on $A\times O\times R$?  Every joint law factors as $q(a)K(o,r\mid a)$ for some pair,
so the domain is rich enough.  But in any one choice situation the channel is fixed: the agent moves the
marginal on $A$, never the coupling.  A single preference would rank couplings, and no choice reveals those
rankings.} 

\bigskip
Axioms are imposed on the family $\{\succeq_K\}_K$.
The first two axioms are standard, and Axiom~\textup{(A3)} is the usual nondegeneracy condition on material
outcomes.  Axioms~\textup{(A4)}, \textup{(A6)}, and \textup{(A7)} carry the preference for trace.
Axiom~\textup{(A5)} puts the separate within-channel rankings on a common scale, while
Axiom~\textup{(A8)} imposes consistency under compounding; neither fixes the sign of the trace term.

\begin{description}[leftmargin=*,style=sameline]

\item[(A1) Weak order.]
For every finite joint channel $K$, $\succeq_K$ is complete and transitive on $\D(A)$.

\item[(A2) Continuity.]
For each fixed triple of alphabets $(A,O,R)$, the set
\[
 \{(K,q,p):q\succeq_Kp\}
 \subseteq \D(O\times R)^A\times\D(A)^2
\]
is closed.


\item[(A3) Material relevance.]
There are distinct outcomes $o^+,o^-\in O$ and a two-action channel $K$ with a one-point record set in which $a^+$ delivers
$o^+$ for sure and $a^-$ delivers $o^-$ for sure.  For this channel,
\[
 \delta_{a^+}\succ_K\delta_{a^-}.
\]
The agent strictly ranks at least two outcomes when each is delivered for sure and the two actions leave the same record.

\item[(A4) Trace relevance.]
There is an outcome $o_\ast\in O$ and a four-action channel $K$ in which every action delivers $o_\ast$ for
sure.  Actions $a_1$ and $a_2$ leave distinct records $r_1$ and $r_2$, respectively, while each of $a_3$ and
$a_4$ leaves either record with equal probability.  For this channel,
\[
 \frac12(\delta_{a_1}+\delta_{a_2})
 \succ_K
 \frac12(\delta_{a_3}+\delta_{a_4}).
\]
Both are fair lotteries and produce the same distribution of visible consequences. The axiom requires a strict
preference for the first, whose record identifies the realised action, over the second, whose record does not.

\end{description}

\paragraph{Comparison environments.}
Some axioms compare lotteries carried by different channels.  Each relation $\succeq_K$ ranks lotteries within
one channel, so such comparisons need a domain in which two channels appear in a single preference statement.  A
larger channel that contains both as labelled blocks provides it.

For
$K:A\to\D(O\times R)$ and $L:B\to\D(O\times R_L)$, let $K\sqcup L$ be the block-diagonal joint channel---the
\emph{block environment}---with action set
$(A\times\{0\})\cup(B\times\{1\})$, payoff set $O$, and record set
$(R\times\{0\})\cup(R_L\times\{1\})$: each block operates its own channel, and all cross-block probabilities
are zero.  For $q\in\D(A)$ and $p\in\D(B)$,
let $q^0$ and $p^1$ denote their
block-supported copies, defined on the tagged action set by
\[
\begin{aligned}
 q^0(a,0)&:=q(a), & q^0(b,1)&:=0,\\
 p^1(a,0)&:=0,    & p^1(b,1)&:=p(b),
\end{aligned}
\qquad(a\in A,\ b\in B).
\]

Under \(q^0\), the visible consequence is generated by \(K\) after the action is drawn according to \(q\); under \(p^1\), it is generated by \(L\) after the action is drawn according to \(p\). So a choice between them is
also a choice between the two channels, made inside a single environment.
Write
\[
 (q,K)\succeq(p,L)
 \quad\Longleftrightarrow\quad
 q^0\succeq_{K\sqcup L}p^1.
\]
Each instance of this shorthand is a single comparison inside the two-block environment $K\sqcup L$, which is
built from the two channels; $\succ$ and $\sim$ denote the asymmetric and symmetric parts of this shorthand
relation.  The pair $(q,K)$
has no meaning in isolation.\footnote{Formally, $(K\sqcup L)(o,(r,0)\mid(a,0))=K(o,r\mid a)$ and
$(K\sqcup L)(o,(r_L,1)\mid(b,1))=L(o,r_L\mid b)$, with every other entry zero.  In the binary case, with
$A=\{a_1,a_2\}$ and $B=\{b_1,b_2\}$, the block environment has the four tagged actions $(a_1,0)$, $(a_2,0)$,
$(b_1,1)$, $(b_2,1)$; a lottery $q=(q(a_1),q(a_2))$ becomes $q^0=(q(a_1),q(a_2);0,0)$ and $p=(p(b_1),p(b_2))$
becomes $p^1=(0,0;p(b_1),p(b_2))$, the semicolon separating the blocks.  Then $(q,K)\succeq(p,L)$ abbreviates
$q^0\succeq_{K\sqcup L}p^1$: one ordinary comparison of two lotteries in this one four-action environment.}  The construction extends to any finite family
$\{K_i:A_i\to\D(O\times R_i)\}_{i\in I}$: write $\bigsqcup_{i\in I}K_i$ for the analogous block environment and
$q_i^i$ for the lottery that assigns $q_i(a)$ to $(a,i)$ and zero to every other block.

\begin{description}[leftmargin=*,style=sameline]

\item[(A5) Block-comparison coherence.]
\emph{Duplication:} for every $K:A\to\D(O\times R)$ and all $q,p\in\D(A)$,
\[
 q\succeq_Kp
 \quad\Longleftrightarrow\quad
 (q,K)\succeq(p,K).
\]
\emph{Irrelevant blocks:} for every finite block environment $\bigsqcup_{k\in I}K_k$, every ordered pair of
distinct blocks $i,j\in I$, and all $q_i\in\D(A_i)$ and $q_j\in\D(A_j)$,
\[
 q_i^i\succeq_{\bigsqcup_{k\in I}K_k}q_j^j
 \quad\Longleftrightarrow\quad
 (q_i,K_i)\succeq(q_j,K_j).
\]

Comparing $q$ and $p$ across two tagged copies of $K$ is just comparing them in $K$; adding blocks not involved
in a comparison leaves it unchanged.

\end{description}

\paragraph{Processing.}
The next two axioms use two ways of transforming an environment.  A \emph{record processor} is a stochastic
kernel $T:O\times R\to\D(O\times R')$ that garbles the record without changing the payoff:
$T(o',r'\mid o,r)=0$ whenever $o'\neq o$.  The garbled channel is simply the matrix product
$KT:A\to\D(O\times R')$,
\begin{equation}
 (KT)(o',r'\mid a)
 :=\sum_{o,r}K(o,r\mid a)T(o',r'\mid o,r).
 \label{eq:record-processing}
\end{equation}

An \emph{action processor} is a stochastic kernel $S:A\to\D(B)$ that replaces the action variable $A$ by a
processed action variable $B$.  Given $q\in\D(A)$, it induces the lottery $qS\in\D(B)$,
\[
 (qS)(b):=\sum_aq(a)S(b\mid a).
\]
A \emph{Bayesian pushforward} of $K$ by $S$ at $q$ is any channel $\widehat K:B\to\D(O\times R)$
satisfying\footnote{If $(qS)(b)=0$, the
equation leaves that row unrestricted; every statement below applies to every consistent completion, making
unreached processed actions irrelevant.}
\begin{equation}
 (qS)(b)\,\widehat K(o,r\mid b)
 =\sum_a q(a)S(b\mid a)K(o,r\mid a)
 \quad\text{for every }(b,o,r).
 \label{eq:action-processing}
\end{equation}
In the processed environment, the alternatives are the noisy names $b$, while the visible consequences remain
$(O,R)$.  Choosing $qS$ under $\widehat K$ replaces dependence on the original action with dependence on its
noisy name, without changing the marginal distribution of visible consequences.

\begin{description}[leftmargin=*,style=sameline]

\item[(A6) Record data processing.]
For every $K:A\to\D(O\times R)$, every record processor $T$, and every $q\in\D(A)$,
\[
 (q,K)\succeq(q,KT).
\]
The agent never gains when her record is garbled.

\item[(A7) Action data processing.]
For every $K:A\to\D(O\times R)$, every action processor $S:A\to\D(B)$, every $q\in\D(A)$, and every Bayesian
pushforward $\widehat K$ of $K$ by $S$ at $q$,
\[
 (q,K)\succeq(qS,\widehat K).
\]
Nor does she gain when the action variable is processed, holding the distribution of visible consequences fixed. Special cases include swapping action labels, adding or deleting unused actions, splitting an action into identical copies, merging actions, and noisy recoding. 
\end{description}

\paragraph{Compounding.}
For $k\ge1$, let $P:A\to\D(\{1,\ldots,k\})$ be a first-stage record channel and let
$(K_1,\ldots,K_k)$ be an ordered list of continuation channels $K_i:A\to\D(O\times R)$.  Their
\emph{compound environment} is
\begin{equation}
 \bigl(P\triangleright(K_1,\ldots,K_k)\bigr)(o,(i,r)\mid a)
 :=P(i\mid a)K_i(o,r\mid a),
 \qquad i=1,\ldots,k.
 \label{eq:compound}
\end{equation}


For expositional purposes, consider the following two-stage description of the compound environment.\footnote{The stage structure merely factorises the compound channel, which is a one-shot experiment. Preferences over the timing of uncertainty resolution are therefore abstracted from.}
  First, the action $a$ generates only
an interim record $i$ through $P$; no payoff is delivered at this stage.  Record $i$ selects the continuation
channel $K_i$, through which the same action generates the payoff $o$ and a further record $r$.  The agent does not
act again.  The interim record remains visible, so the final record is $(i,r)$.  Under $q$ and $P$, record $i$ is
\emph{reached} if $\sum_aq(a)P(i\mid a)>0$.  For every reached record, let $q_i(a):=q(a)P(i\mid a)/\sum_{a'}q(a')P(i\mid a')$ denote the posterior distribution
of the action.
\begin{description}[leftmargin=*,style=sameline]

\item[(A8) Recordwise sure-thing principle.]
For every $q\in\D(A)$, every binary first-stage record channel $P:A\to\D(\{1,2\})$ for which record $1$ is
reached, and every three continuation channels $K,L,M:A\to\D(O\times R)$,
\begin{equation}
 (q_1,K)\succeq(q_1,L)
 \quad\Longleftrightarrow\quad
 \bigl(q,P\triangleright(K,M)\bigr)
 \succeq
 \bigl(q,P\triangleright(L,M)\bigr).
 \label{eq:recordwise-sure-thing}
\end{equation}

Fix $q$ and a binary first-stage record channel $P$, and suppose record $1$ is reached, inducing posterior $q_1$.
The two compound environments share the continuation $M$ after record $2$ and differ only after record $1$, where
one continues with $K$ and the other with $L$.  The axiom says that the compounds are ranked exactly as the lottery
$q_1$ in channel $K$ and the same lottery in channel $L$ are ranked in their block environment.  The common
continuation on the other branch cannot affect the comparison.\footnote{This is a version of Savage's sure-thing principle
with the event given by an interim record \cite{Savage1954}.  Unlike the exogenous event in Savage's formulation,
reaching the record can be informative here: $P(1\mid a)$ may vary with the action.  This is why the conditional
comparison is evaluated at $q_1$.  If record $1$ is never reached, the axiom is silent, as on Savage's null
events.  When $P$ is an action-independent public coin, $q_1=q$, and
Axiom~\textup{(A8)} reduces to ordinary mixture independence with the coin retained in the record.}

\end{description}

% ===============================================================
\section{The representation theorem}\label{sec:theorem}
% ===============================================================

For $q\in\D(A)$ and $K:A\to\D(O\times R)$, let
\[
 p_{qK}(o,r):=\sum_a q(a)K(o,r\mid a)
\]
be the induced distribution of visible consequences, and define
\begin{equation}
 \I_{qK}(A;O,R)
 :=\sum_{a,o,r}q(a)K(o,r\mid a)
 \log\frac{K(o,r\mid a)}{p_{qK}(o,r)}.
 \label{eq:mi}
\end{equation}
Summands with $q(a)K(o,r\mid a)=0$ are set to zero, the logarithm has any fixed base greater than one, and expectations below are
taken under the joint distribution $q(a)K(o,r\mid a)$.

\begin{theorem}\label{thm:main}
For a family $\{\succeq_K\}_K$ as described above, the following are equivalent.
\begin{enumerate}[label=\textup{(\roman*)},leftmargin=*]
\item The family satisfies Axioms~\textup{(A1)--(A8)}.
\item There are a nonconstant function $u:O\to\mathbb R$ and a number $\lambda>0$ such that, for every nonempty
finite $A,R$, every $K:A\to\D(O\times R)$, and every $q,p\in\D(A)$,
\begin{equation}
 q\succeq_Kp
 \quad\Longleftrightarrow\quad
 V(q,K)\ge V(p,K),
 \qquad
 V(q,K):=\E_{qK}[u(O)]+\lambda\I_{qK}(A;O,R).
 \label{eq:representation}
\end{equation}
\end{enumerate}
Moreover, under these equivalent conditions, block-supported cross-channel comparisons are represented on the same
scale: for every finite block environment $\bigsqcup_{k\in I}K_k$, every two distinct blocks $i,j\in I$, and every
$q_i\in\D(A_i)$ and $q_j\in\D(A_j)$,
\[
 q_i^i\succeq_{\bigsqcup_{k\in I}K_k}q_j^j
 \quad\Longleftrightarrow\quad
 V(q_i,K_i)\ge V(q_j,K_j).
\]
\end{theorem}

We call a family satisfying \textup{(A1)--(A8)} \emph{trace-tempered}, and likewise the criterion that
represents it.  Its central restriction is that payoff and trace enter separately.  The agent may value good
outcomes and may value evidence about her action, but she cannot place an extra value on a good payoff being the
revealing consequence.  Thus, if under the same lottery two channels give the same expected utility from payoffs
and induce the same distribution of posteriors over actions, the agent is indifferent between them, regardless of
how payoffs and records are paired (Corollary~\ref{cor:matching}, Appendix~\ref{app:representation-consequences}).  This
separation is derived, not assumed.

\subsection{Proof sketch}
That the representation satisfies the axioms is a direct verification.  The converse has four steps.  It first establishes the common affine
machinery, then restricts it to the constant-payoff domain and proves a pure-trace lemma, returns to general
payoffs to align the material and trace scales, and finally assembles arbitrary channels branch by branch.

\emph{Step 1: reduction to terminal laws.}  Fix a full-support $q$.  If two channels induce the same joint
distribution of the realised payoff and the posterior over actions, each is a record processing of the other.
Block coherence and record processing therefore make them indifferent.  Action processing makes action
relabellings immaterial and, for an arbitrary $q$, permits null actions to be deleted.  Thus $(q,K)$ is identified
by its \emph{terminal law}, the joint distribution of the realised payoff and posterior
(Lemma~\ref{lf:lem:terminal-affine}\textup{(a)} and
Lemma~\ref{lf:lem:bookkeeping}\textup{(a)--(b)}).

\emph{Step 2: cardinalisation.}  Let a public coin select between two channels.  Because the coin is independent
of the action, the posterior in either branch remains $q$.  The recordwise sure-thing principle therefore gives
mixture independence over terminal laws.  Together with the weak order and continuity, the mixture-space theorem
gives an affine representation on each full-support fibre (Lemma~\ref{lf:lem:terminal-affine}\textup{(b)}).

\emph{Step 3: payoff and trace.}  On singleton action sets there is no trace, so the affine representation reduces
to expected utility.  Material relevance makes its index $u$ nonconstant, while action processing and block
coherence make the same payoff scale apply throughout.  At the outcome $o_\ast$ in the channel from Axiom~\textup{(A4)},
the payoff term is constant.  The constant-payoff phase of Appendix~\ref{app:main-proof} shows that the remaining order is represented by
mutual information.  The finite-branch form of Axiom~\textup{(A8)} supplies the ordinal continuation comparisons;
fibrewise affinity and scale coherence
convert them into an exact chain rule.  Together with action relabelling and support transport, that rule gives
Faddeev's grouping property, while Axiom~\textup{(A2)} supplies continuity and block coherence aligns comparisons
across priors.  Hence the pure-trace order is a scalar multiple of Shannon entropy \cite{Faddeev1956}.
Record data processing makes the scalar nonnegative, trace relevance makes it strictly positive, and material
calibration determines one coefficient $\lambda$ relative to $u$
(Lemmas~\ref{lf:lem:material-scale} and~\ref{lf:lem:trace-scale}).

\emph{Step 4: assembly.}  The finite-branch sure-thing result makes insertion on a reached branch an order
embedding.  For a fixed background, the marked terminal law is affine in the inserted continuation; the affine
fibre representation therefore makes this embedding affine.
Calibration on the constant-$o_\ast$ face and the mutual-information chain rule identify its coefficient as the
branch probability; crossed public mixtures show that this coefficient is independent of the other branches.
Begin with a compound that pays $o_\ast$
after every first-stage record.  If record $y$ occurs with probability $m_y$, changing only the payoff after $y$
from $o_\ast$ to $o$ changes value by
$
 m_y\bigl[u(o)-u(o_\ast)\bigr].
$
This baseline compound has value
\[
 u(o_\ast)+\lambda\I_q(A;Y).
\]
Summing across branches gives the representation when the payoff is determined by the first-stage record.  For
an arbitrary channel, factor it through the interim record $(o,r)$ and a deterministic continuation whose payoff
is $o$.  The resulting compound channel and
the original channel are record processings of one another and hence are indifferent.  The representation
follows (Lemmas~\ref{lf:lem:branch-weight} and~\ref{lf:lem:deterministic-assembly}).

\subsection{Uniqueness}

A real-valued index $W$ on all finite lottery--channel pairs is a \emph{common-scale representation} if, for every channel
$K:A\to\D(O\times R)$ and every $q,p\in\D(A)$,
\[
 q\succeq_Kp
 \quad\Longleftrightarrow\quad
 W(q,K)\ge W(p,K),
\]
and, for every finite block environment $\bigsqcup_{k\in I}K_k$, every two distinct blocks $i,j\in I$, and all
$q_i\in\D(A_i)$ and $q_j\in\D(A_j)$,
\[
 q_i^i\succeq_{\bigsqcup_{k\in I}K_k}q_j^j
 \quad\Longleftrightarrow\quad
 W(q_i,K_i)\ge W(q_j,K_j).
\]

A common-scale representation $W$
is \emph{branch-additive} if, for every $k\ge1$, every finite $A$, every $q\in\D(A)$, every $P:A\to\D(\{1,\ldots,k\})$, and every
two ordered lists of continuation channels $(K_1,\ldots,K_k)$ and $(L_1,\ldots,L_k)$, with
$K_i,L_i:A\to\D(O\times R)$, writing
$m_i:=\sum_aq(a)P(i\mid a)$,
\begin{equation}
 W\bigl(q,P\triangleright(K_1,\ldots,K_k)\bigr)
 -W\bigl(q,P\triangleright(L_1,\ldots,L_k)\bigr)
 =\sum_{i:m_i>0}m_i
 \bigl[W(q_i,K_i)-W(q_i,L_i)\bigr].
 \label{eq:branch-additive-index}
\end{equation}
In words, a continuation change worth $x$ after record $i$ changes ex-ante value by $m_i x$, and changes on
different branches add. Axiom~\textup{(A8)} is ordinal: it says that a common continuation on the other branch cannot affect a conditional
comparison.  By Lemma~\ref{pt:lem:finite-branch-extension}\textup{(a)}, this is equivalent---given the other structural
axioms---to the requirement that branchwise improvements cannot be reversed by finite compounding.  Branch
additivity is the corresponding requirement on numerical differences.  

\begin{proposition}[Uniqueness]\label{prop:uniqueness}
Suppose \textup{(A1)--(A8)} hold, and fix a representation $V$, with parameters $(u,\lambda)$, as in
Theorem~\ref{thm:main}.  For any real-valued index $W$ on all finite lottery--channel~pairs:
\begin{enumerate}[label=\textup{(\roman*)},leftmargin=*]
\item $W$ is a common-scale representation if and only if $W=\varphi\circ V$ for a single strictly increasing
$\varphi:[\underline u,\infty)\to\mathbb R$, where $\underline u:=\min_{o\in O}u(o)$.
\item A common-scale representation $W$ is branch-additive if and only if
$W=\alpha V+\beta$ for some $\alpha>0$ and $\beta\in\mathbb R$.
\end{enumerate}
\end{proposition}

The proof is in Appendix~\ref{app:representation-consequences}.

Every trace-tempered representation is a common-scale representation and, by the expected-utility and
mutual-information chain rules, branch-additive.  Part~\textup{(ii)} therefore implies that if $(u,\lambda)$ and
$(\widetilde u,\widetilde\lambda)$ are two such representations, then
\[
 \widetilde u=\alpha u+\beta,
 \qquad
 \widetilde\lambda=\alpha\lambda
\]
for some $\alpha>0$ and $\beta\in\mathbb R$.  Once two non-indifferent payoff outcomes are assigned numerical
utilities, $\lambda$ is unique for the chosen logarithm base.  In the constant-payoff phase of
Appendix~\ref{app:main-proof}, multiplying the pure-trace index by a positive constant has no behavioural content.  Here
the payoff normalisation fixes that scale, and $\lambda$ measures the trade-off between payoff and trace.

\subsection{Sign of the trace weight}

Axioms~\textup{(A4)},
\textup{(A6)}, and \textup{(A7)} orient the trace component.  Write $(\mathrm{A}4^-)$,
$(\mathrm{A}6^-)$, and $(\mathrm{A}7^-)$ for the clauses obtained by reversing their displayed comparisons,
and write $(\mathrm{A}4^0)$, $(\mathrm{A}6^0)$, and $(\mathrm{A}7^0)$ for the corresponding indifference clauses.

\begin{corollary}[Sign trichotomy]\label{cor:signduality}
Suppose Axioms~\textup{(A1)--(A3)}, \textup{(A5)}, and
\textup{(A8)} hold.  Then:
\begin{enumerate}[label=\textup{(\roman*)},leftmargin=*]
\item under \textup{(A4)}, \textup{(A6)}, and \textup{(A7)}, the family has the representation in
Theorem~\ref{thm:main}, with $\lambda>0$;
\item under $(\mathrm{A}4^-)$, $(\mathrm{A}6^-)$, and $(\mathrm{A}7^-)$, it has the same representation with
$\lambda<0$;
\item under $(\mathrm{A}4^0)$, $(\mathrm{A}6^0)$, and $(\mathrm{A}7^0)$, it is represented by expected utility,
equivalently by the same formula with $\lambda=0$.
\end{enumerate}

\end{corollary}

The proof is in Appendix~\ref{app:representation-consequences}.  The directions of the three orientation clauses cannot be chosen independently.  For example, under the
corollary's background axioms and any orientation of Axiom~\textup{(A7)}, Axiom~\textup{(A4)} is inconsistent
with $(\mathrm{A}6^-)$.  Restrict the two lotteries in Axiom~\textup{(A4)} to their supports and relabel the
actions.  The second channel is then a total garbling of the first.  Axiom~$(\mathrm{A}6^-)$ weakly prefers the
second, whereas Axiom~\textup{(A4)} strictly prefers the first.


\subsection{Independence of the axioms}

\begin{proposition}[Independence of the axioms]\label{prop:independence}
For each $i\in\{1,\ldots,8\}$ there is a family satisfying the other seven axioms and violating
Axiom~\textup{(A$i$)}.
\end{proposition}

The eight families are constructed in Appendix~\ref{app:axiom-independence}.  Three are of independent interest.  A conditional-entropy penalty satisfies every axiom except record data processing and values the removal of
extraneous record noise.  A Gini posterior-reduction criterion satisfies every axiom except action data processing.
The satiation criterion $\E[u(O)]+\min\{I(A;O,R),c\}$, for $c>0$, satisfies every axiom except the recordwise
sure-thing principle and ceases to value trace once $c$ units of information are secured.

% ===============================================================
\section{Choice implications}\label{sec:implications}
% ===============================================================

This section works out what an agent with the representation of Theorem~\ref{thm:main} does: which tests
falsify the model directly (\S\ref{sec:tracetest}), when she deliberately randomises (\S\ref{sec:randomisation}),
how $\lambda$ can be measured from choices (\S\ref{sec:measurement}), and how the criterion works in a
standard contracting model (\S\ref{sec:application}).
Throughout, after fixing the affine normalisation of $u$ and the logarithm base, $\lambda$ is measured in payoff units per unit of information.  Proofs for this
section, together with full statements of the supporting results that the text presents in summary form, are
collected in Appendix~\ref{app:choice-proofs}.

\subsection{The trace test}\label{sec:tracetest}

The criterion $V$ reads the joint distribution of action and visible consequence, rather than the consequence
distribution or action lottery alone.  The following two-treatment design separates these objects.

\begin{proposition}[The trace test]\label{prop:signature}
Four actions pay the same outcome $o_0$.  Two of them leave distinct marks: $a_1$ always produces record $r_1$,
and $a_2$ always produces the distinct record $r_2$.  The other two are indistinguishable: $a_3$ and $a_4$ each
produce $r_1$ or $r_2$ with equal probability.  Call this channel $K$, and let $q'$ be the fair coin over
$\{a_1,a_2\}$ and $q''$ the fair coin over $\{a_3,a_4\}$.
\begin{enumerate}[label=\textup{(\roman*)},leftmargin=*]
\item The two coins induce the same distribution of visible consequences, $p_{q'K}=p_{q''K}$; yet
\[
 V(q',K)-V(q'',K)=\lambda\log2>0,
 \qquad\text{so}\qquad
 q'\succ_Kq''.
\]
\item Let $K^\circ$ be the channel in which every action pays $o_0$ and produces $r_1$ or $r_2$ with equal
probability.  Then $q'\sim_{K^\circ}q''$.
\end{enumerate}
\end{proposition}


Part \textup{(i)} rules out any criterion that depends only on the distribution of visible consequences.
The two coins induce the same distribution over payoff--record pairs, so any such criterion is indifferent between them. However, a theory that also reads the lottery can
match it.  Existing theories of deliberate randomisation are of this kind: in the hedging model
of Cerreia-Vioglio, Dillenberger, Ortoleva, and Riella \cite{CerreiaVioglioDillenbergerOrtolevaRiella2019}, and
in the preference for randomisation documented by Agranov and Ortoleva \cite{AgranovOrtoleva2017}, the value of
mixing depends only on the induced distribution and the lottery.  Call an index $W$ \emph{channel-blind} if it
reads only these two objects: there is a function $F$, fixed across channels, such that
\[
 W(q,K)=F\bigl(p_{qK},q\bigr).
\]

Part \textup{(ii)}  rules these out.  Nothing a channel-blind index reads changes between the treatments: in
all four lottery--channel combinations the distribution of visible consequences is uniform on
$\{(o_0,r_1),(o_0,r_2)\}$, and the lotteries are the same two lotteries.  A channel-blind index therefore
separates the coins in both treatments or in neither, so none reproduces the choice pattern of
Proposition~\ref{prop:signature}.%
\footnote{If $W(q,K)=F(p_{qK},q)$ is channel-blind and $F(p,\cdot)$ is invariant under permutations of action labels---an entropy bonus, say---the first
treatment is enough: $p_{q'K}=p_{q''K}$ and $q''$ is a permutation of $q'$, so $W(q',K)=W(q'',K)$.  The second treatment is needed for label-sensitive indices.}

The first treatment also identifies the orientation of taste: it selects the revealing coin when $\lambda>0$,
the noisy coin when $\lambda<0$, and indifference when $\lambda=0$.  A strict preference between the coins in the second treatment rejects all three cases of Corollary~\ref{cor:signduality}, and with them the model.

Overall, the choice pattern in the trace test detects a taste for action--consequence dependence; by itself it does not identify mutual information among all dependence indices.

\subsection{Deliberate randomisation}\label{sec:randomisation}


Consider a decision maker free to choose lotteries among two actions, and let $a$ pay weakly more.  Playing $a$ for sure maximises payoff and leaves no
trace; mixing in $b$ costs payoff but gives the record something to reveal.  Write $K_a:=K(\cdot\mid a)$ for the distribution of visible
consequences under action $a$, $\bar u(a):=\E_{K_a}[u(O)]$ for its expected utility, and
$\Delta:=\bar u(a)-\bar u(b)\ge0$ for the payoff gap.  Each unit of probability moved from $a$ to $b$ costs
$\Delta$.  What it buys is evidence.  A consequence is evidence about the act that produced it through a
likelihood ratio: how much more probable it is given the realised act than given only the mixture she plays.
With all probability on $a$, consequences are distributed exactly as $K_a$: every draw has ratio one and
carries no evidence.  Now move a marginal unit of probability to $b$.  The mixture's consequences are still,
to first order, distributed as $K_a$; but when $b$ realises, its consequence is drawn from $K_b$, so the draw
carries $\log\bigl(K_b(o,r)/K_a(o,r)\bigr)$.  Its expected value over $b$'s draws is, by definition, the
divergence $\KL(K_b\,\|\,K_a)$.\footnote{For distributions $P$ and $Q$,
$\KL(P\,\|\,Q):=\sum_{o,r}P(o,r)\log\bigl(P(o,r)/Q(o,r)\bigr)$, with zero-mass summands set to zero.  The
divergence is $+\infty$ if $P$ assigns positive probability to a consequence and $Q$ assigns zero.}  Because
information is concave in the lottery, whether to mix at all is settled at the margin (Proposition~\ref{prop:binary},
Appendix~\ref{app:randomisation-proofs}):
\[
 \delta_a\text{ is optimal}
 \quad\Longleftrightarrow\quad
 \Delta\ge\lambda\,\KL(K_b\,\|\,K_a).
\]
When the payoff gap is below $\lambda\KL(K_b\,\|\,K_a)$, the optimal lottery is unique and mixes the two
actions, with weakly more probability on $b$ as $\lambda$ rises---strictly when $\Delta>0$---and less as $\Delta$ rises.


With more than two actions, each action is compared with the distribution of visible consequences generated by
the lottery as a whole, $p_{qK}$. An action is credited with its expected utility and with how far its visible consequences depart from those generated on average.  At an optimum this total is the same for every action in use; an unused action cannot
offer more.

\begin{proposition}[Optimality condition]\label{prop:optimal-choice}
Let $\lambda>0$.  A lottery $q^*$ is optimal if and only if there is a constant $c$ such that
\[
 \bar u(a)+\lambda\KL\!\left(K_a\,\middle\|\,p_{q^*K}\right)=c
 \quad\text{for }a\in\supp(q^*),
 \qquad
 \le c\quad\text{for }a\notin\supp(q^*).
\]
\end{proposition}


The condition has an immediate implication: every visible consequence that can arise under some action
occurs with positive probability under any optimal lottery (Proposition~\ref{cor:support-pure},
Appendix~\ref{app:randomisation-proofs}).  Hence, if each action can produce some consequence that no other action can produce, every optimal lottery uses every action. At the other extreme, suppose every action produces the same distribution of payoffs and records. The consequence then reveals nothing about the action, and all actions have the same expected utility. Hence the agent is indifferent among all lotteries.\footnote{With $\lambda<0$, every payoff-maximising pure action is optimal.  More generally, $q$ is optimal in a fixed channel if and only if it is supported on payoff-maximising actions and $\I_{qK}(A;O,R)=0$.}


Now suppose that each action is fully identified by its visible consequences.  The next result shows that behaviour is then multinomial logit and the preference for trace becomes a preference for randomisation. 

\begin{corollary}[Multinomial logit]\label{rem:logit}
If the row supports of $K$ are pairwise disjoint, then $\I_{qK}(A;O,R)=\Sh(q)$.  Therefore, the unique optimum is
\[
 q^*(a)
 =\frac{\exp\bigl(\bar u(a)/\lambda\bigr)}
 {\sum_{a'}\exp\bigl(\bar u(a')/\lambda\bigr)},
\]
with $1/\lambda$ as the payoff sensitivity when information is measured in natural logarithms.\footnote{The same formula arises in rational inattention, where, in contrast, $q$ is fixed and $K$ is chosen \cite{MatejkaMcKay2015}.}
\end{corollary}

Here, Debreu's red-bus/blue-bus conclusion depends on what is recorded.  If blue is an exact copy of red, only their combined probability is pinned down. Any split of the red bus passengers between the two buses is optimal (Proposition~\ref{prop:optimal-marginal}, Appendix~\ref{app:randomisation-proofs}).  If colour is recorded, however, red and blue are distinct traces.  With equal expected utilities, the corollary assigns probability $1/3$ to each alternative.  The model is invariant to duplicate names, not to distinct recorded traces.


With overlapping supports, the odds between two actions can depend on the rest of the menu.  The optimality
condition sets each action's probability against the average consequence distribution, and an added action
shifts that average toward its own consequences, reducing the measured distinctiveness of the actions it
resembles more than of the others; the odds then move.  Pairwise-disjoint supports therefore imply independence
of irrelevant alternatives, while overlap permits, but does not require, its failure.


\subsection{Measuring the trace weight}\label{sec:measurement}

Under the stated normalisation, the coefficient $\lambda$ is the price of one unit of information in expected-utility units.  Any indifference
between two environments with different amounts of trace reveals it: writing $E_{qK}:=\E_{qK}[u(O)]$, if
$(q,K)\sim(p,L)$ and $\I_{qK}\ne\I_{pL}$, the representation rearranges to
\[
 \frac{E_{pL}-E_{qK}}{\I_{qK}-\I_{pL}}=\lambda,
\]
and $\lambda$ is the only number consistent with every such indifference
(Proposition~\ref{prop:single-price}, Appendix~\ref{app:measurement-proofs}).

Plot expected payoff utility against information: indifference curves are parallel lines of slope $-\lambda$.  The slope does
not depend on the prior, the payoff level, the action alphabet, or the amount of information already present.
This constancy is the observable restriction that Axiom~\textup{(A8)} adds to the rest of the system.  A taste for trace that satiates satisfies every axiom except~\textup{(A8)} (Proposition~\ref{prop:independence}) and has a slope that falls to zero beyond the satiation point.


The identity presupposes the utility function; the design in the next proposition does not, because its only
payoffs are the two normalised outcomes and a baseline that cancels.  At a constant payoff, revelation of the
realised action is worth $\lambda\Sh(q)$ more than silence, so compensating the silent side until the agent is
indifferent measures her willingness to pay for trace.

\begin{proposition}[Eliciting the trace weight]\label{prop:elicit}
Normalise $u(o^+)=1$ and $u(o^-)=0$.  Fix $q\in\D(A)$ with $\Sh(q)>0$, a constant payoff, and two channels, one
whose record reveals the realised action and one whose record is silent.  Dilute both by the same public side
branch, reached with probability $\varepsilon\in(0,1)$: on the revealing side the branch pays $o^-$ for sure; on
the silent side it pays $o^+$ with probability $\beta$ and $o^-$ otherwise.  If
$\lambda\Sh(q)\le\varepsilon/(1-\varepsilon)$, a unique $\beta\in[0,1]$ produces indifference, and
\[
 \lambda=\frac{\varepsilon}{1-\varepsilon}\,\frac{\beta}{\Sh(q)};
\]
otherwise the revealing side is strictly preferred at every $\beta$.
\end{proposition}

The common baseline payoff's utility cancels, so it need not be known; the information gap $\Sh(q)$ is set by
the design, so no information quantity need be estimated.  The proof, the general compensation identity, and
the experimental controls are in Appendix~\ref{app:measurement-proofs} (Proposition~\ref{prop:wtp} and
Remark~\ref{rem:elicitation}).

The design also tests the model.  Run it twice with the same $\varepsilon$, at $q=(1/2,1/2)$ and at
$q=(1/4,3/4)$.  Two runs must return the same $\lambda$. Hence, the two indifference
reports must satisfy
\[
 \frac{\beta_1}{\beta_2}
 =\frac{\Sh(1/2,1/2)}{\Sh(1/4,3/4)}
 =\frac{\log 2}{\frac14\log4+\frac34\log\frac43}
 \approx1.23.
\]
 Any other ratio is not consistent with 
the model.  This number also separates Shannon entropy from the Gini-reduction rival
$1-\sum_aq(a)^2$, which instead commits to $4/3$.


% ===============================================================
\subsection{An application: moral hazard with a trace-tempered agent}\label{sec:application}
Take the simplest version of Holmstr\"om's moral-hazard problem \cite{Holmstrom1979}: a principal cannot observe
the agent's effort and offers a wage contingent on observable output.  The risk-neutral agent's effort is binary,
$A\in\{0,1\}$; high effort costs $c>0$ and low effort is costless.  Output $X\in\{0,1\}$ succeeds with
probability $p_a$ under effort $a$, where $0<p_0<p_1<1$, and the contract pays $W=w_0+\Delta wX$, with
$\Delta w\ge0$.  Write
$K_a$ for the Bernoulli row with parameter $p_a$, let $q:=\Pr(A=1)$, and attach the expected effort cost
$qc$ to the chosen lottery.\footnote{Writing cost at the level of the lottery prevents its realisation
from becoming a trace of effort.  The cost can nevertheless be kept inside the paper's finite consequence
domain.  For each candidate $q$, append to both action rows a payoff coordinate
$Z_q\in\{0,c\}$ with $\Pr(Z_q=c)=q$, and take $u(w,z)=w-z$.  Then $\E Z_q=qc$ and, because $Z_q$ is
independent of all production variables while $W$ is a function of $X$,
$\I(A;W,X,Z_q)=\I(A;X)$.  A plan $q$ is represented by the binary lottery $(1-q,q)$ in its augmented channel; any two
such plans are compared by placing their channels in separate blocks and applying the block-supported
common-scale clause of Theorem~\ref{thm:main}.}  

For a given contract $(w_0, \Delta w)$,  the agent chooses $q\in[0,1]$ to maximise
\[
 V(q)=w_0+\bigl[p_0+q(p_1-p_0)\bigr]\Delta w-qc
      +\lambda\,\I_{qK}(A;X).
\]
Pure high effort is optimal if and only if
\[
 \Delta w\ge \frac{c+\lambda\,\KL(K_0\,\|\,K_1)}{p_1-p_0}
\]
(Proposition~\ref{prop:binary}).  Thus the performance bonus must cover both effort cost and the trace forgone
at a pure action.  At a flat wage, the unique optimum is interior if and only if
\[
 c<\lambda\,\KL(K_1\,\|\,K_0),
\]
so trace alone can sustain some effort without an incentive bonus.

Now consider the team setting from \cite{Holmstrom1982}. Suppose $n$ agents choose efforts and the output is $X\in\{0,1\}$  with success probability $\varepsilon < g(a_1,\ldots,a_n) < 1-\varepsilon$ for some $\varepsilon>0$. Fix the
strategies of agent $i$'s teammates.  If $i$ works, the signal is positive with probability
$\bar g_1:=\E[g(1,A_{-i})]$; if she shirks, with probability $\bar g_0:=\E[g(0,A_{-i})]$.  

The agent's problem is the one described above with $p_a$ replaced by $\bar g_a$, where $0<\bar g_0<\bar g_1<1$, and both conditions apply as
written.  The minimum bonus that sustains pure high effort is now
\[
 \Delta w^*=\frac{c+\lambda\,\KL(\bar K_0\,\|\,\bar K_1)}{\bar g_1-\bar g_0},
\]
where $\bar K_a$ is the Bernoulli row with parameter $\bar g_a$.

In technologies in which one member's effect vanishes as the team grows, as when success depends on the
fraction of members working, success is almost as likely whether she works or shirks, and a draw of $X$
says almost nothing about her effort.
The gap $\bar g_1-\bar g_0$ then closes.  The effort share of the bonus grows as it does in the standard
model; the divergence is second order in the gap, so the trace share vanishes: trace becomes second order
relative to pay.  The exception is an essential task, whose effect does not shrink with the team.  

Suppose a signal $S$ of her work is visible---posted within the team, say---but not verifiable, so wages
remain contingent on $X$ alone.  Recording it beside $X$, at $q$ held fixed, raises her value by
$\lambda\,\I_q(A_i;S\mid X)$ (Proposition~\ref{prop:garbling-price}, Appendix~\ref{app:record-value-proofs}).
In the standard model such a signal is worthless; here, if posting it costs nothing, the principal can lower
the base wage by the same amount: non-contractible measurement substitutes for pay.


% ===============================================================
\section{Related literature}\label{sec:literature}
% ===============================================================

Grant, Kajii, and Polak \cite{GrantKajiiPolak1998} characterise an intrinsic preference for information
independent of its use in planning; Caplin and Leahy \cite{CaplinLeahy2001} model anticipatory feelings generated
by beliefs, while Ely, Frankel, and Kamenica \cite{ElyFrankelKamenica2015} study preferences over paths of beliefs
and the resulting value of suspense and surprise.  These papers show that information need not be instrumental,
but the relevant beliefs concern uncertainty faced by the agent, whereas the posterior here concerns her
realised action.  The axioms here also restrict its value to expected entropy reduction and fix its rate of
substitution with material utility.

Cerreia-Vioglio, Dillenberger, Ortoleva, and Riella
\cite{CerreiaVioglioDillenbergerOrtolevaRiella2019} study deliberate randomisation and characterise a model in
which mixing hedges uncertainty about how to evaluate lotteries; Agranov and Ortoleva
\cite{AgranovOrtoleva2017} document stochastic choice under announced repetition and evidence that much of it is
deliberate.  Trace-tempered choice can also make a nondegenerate lottery optimal. The test in Proposition~\ref{prop:signature} tells the two models apart.

Corrao, Fudenberg, Mudekereza, and Levine \cite{CorraoFudenbergMudekerezaLevine2026} represent
concave and smooth utilities over lotteries as expected utility plus suspense, the expected error of an
adversary's best forecast of the outcome.  Within a fixed channel the present criterion has this form: the
adversary forecasts the visible consequence, and the suspense is $\lambda\,\I_{qK}(A;O,R)$.  Making the
consequence hard to forecast for someone who does not know the action and making the action easy to infer from
the consequence are the same thing.  Their error function is a primitive.  Here the axioms determine it, in every channel, with the same $\lambda$.

In signaling models, actions convey information about hidden characteristics.  Andreoni and Bernheim
\cite{AndreoniBernheim2009} let chance override the dictator's choice in a dictator game with known probability, and fair behaviour
falls: the observer can no longer tell a selfish choice from chance.  In B\'enabou and Tirole
\cite{BenabouTirole2006}, the agent values the audience's belief that her motives are good; a paid good deed
proves less, so rewards can crowd out good conduct.  In their earlier paper \cite{BenabouTirole2002} the
audience is internal: a future self whose confidence sustains effort.  Bernheim \cite{Bernheim1994} shows how
status concerns can induce agents with heterogeneous preferences to choose the same action, making behaviour less
informative about predispositions---an instrumental counterpart to the deniability case in
Corollary~\ref{cor:signduality}.  In all these models the audience infers a hidden characteristic from the
action, and the response depends on what is inferred: a better type earns more esteem.  Here no response enters
utility, and inference runs the other way, from consequence to action.  The one audience the representation can
absorb pays for how much is learned, not for what (\S\ref{sec:interpretation}).

The proof also relates to axiomatic characterisations of entropy.  Shannon \cite{Shannon1948} characterised
entropy from continuity, monotonicity on uniform distributions, and a grouping recursion; Faddeev
\cite{Faddeev1956} reduced these to symmetry, continuity of the binary entropy, and the recursion itself.  The
proof uses the strong-additivity form of Faddeev's theorem given by Baez, Fritz, and Leinster
\cite{BaezFritzLeinster2011}.
In those characterisations, a numerical information functional is primitive and structural conditions are
imposed directly on it. Here the information chain rule and Faddeev's recursion are derived from axioms on an ordinal family of weak orders over action lotteries. The theorem also calibrates the information term against material utility.

Blackwell \cite{Blackwell1953} compares experiments by their value in all downstream decision problems; Gilboa
and Lehrer \cite{GilboaLehrer1991} characterise which functions on partitions can arise as values of information, and Azrieli and Lehrer \cite{AzrieliLehrer2008} extend the analysis to stochastic
information structures.  Frankel and Kamenica \cite{FrankelKamenica2019} characterise the measures of
information and uncertainty that arise as the value of news.  Cabrales, Gossner, and
Serrano \cite{CabralesGossnerSerrano2013} obtain entropy reduction from the decisions of ruin-averse investors.
For Blackwell's decision-maker, information is an input: she can use it or ignore it, so more
information cannot hurt.  Here it is an output: the record is evidence about the agent, and no
decision follows it.
Uniform preference between two record technologies identifies the more-capable
order of K\"orner and Marton \cite{KornerMarton1977}, not Blackwell garbling
(Proposition~\ref{prop:blackwell}, Appendix~\ref{app:record-value-proofs}).

Mutual information is also the cost of attention in rational inattention \cite{Sims2003,MatejkaMcKay2015}. The difference is in what is fixed and what is chosen. There the distribution of the state is fixed and the agent chooses how her action depends on it. The cost is the information the action carries about the state.  Here the channel from actions to consequences is fixed and the agent chooses the distribution of her action. The benefit is the information the consequence carries about
the action. The data that provide identification differ as
well.  Rational inattention is identified from state-dependent stochastic choice; Denti \cite{Denti2022} gives
testable conditions under which such data are rationalised by posterior-separable costs.  Here the primitive
is a ranking of action lotteries in a fixed channel, and the trace test compares two lotteries that induce the
same distribution of consequences.  When the rows of the channel have disjoint supports, both models produce
choice probabilities of the multinomial-logit form, so choice frequencies alone do not separate them.
Rankings of lotteries that are not chosen and comparisons across channels do.  Caplin, Dean, and Leahy
\cite{CaplinDeanLeahy2022} show that, among uniformly posterior-separable costs, Invariance under Compression
selects Shannon entropy: merging states that pay the same in every action does not change behaviour.  The
corresponding condition here is indifference to undetectable distinctions, which follows from action data
processing: merging two actions with identical rows does not change the ranking
(Lemma~\ref{pt:lem:transport}).  The Gini criterion in Appendix~\ref{app:axiom-independence} satisfies every
axiom except action data processing and fails it at such a merge.

Pomatto, Strack, and Tamuz \cite{PomattoStrackTamuz2023} ask what an experiment costs.  On a cardinal,
prior-free cost function they impose continuity, monotonicity in Blackwell's order, additivity over
independent experiments, and constant marginal cost: an experiment run with probability $\alpha$ costs
$\alpha$ times as much.  Then the only answer is a weighted sum of Kullback--Leibler divergences
between the signal distributions in different states.
Their cost functional and the trace value studied here do not nest.
Full revelation has trace value $\Sh(q)$, finite and prior-dependent, while every nonzero cost in their class
makes it infinitely costly.  The information-cost literature prices
learning about the world; the present paper prices being learned about.


% ===============================================================
\section{Interpretation and scope}\label{sec:interpretation}
% ===============================================================

\paragraph{Evidence, not authorship or control.}
The trace term prices evidence of agency: not authorship, and not control in its everyday sense.  A pure action
fixes the act in advance, although its consequence may remain stochastic; it has
zero trace value even when the consequence names the act perfectly, because an act settled in advance leaves no
uncertainty for the consequence to resolve.  A coin the agent flips herself can do better.  The model does not
say that surrendering a choice to chance creates more agency; it says that an unsettled act can leave more
evidence than a predetermined one.  The primitive ranks lotteries she implements herself; a lottery assigned to her
from outside is not in the domain.  This is the Underground Man's problem: the proof he wants cannot come from an act
fixed in advance.  A reading of control or authorship survives if what is valued is its display rather than its
possession; but the display of which act was taken is exactly what the trace term measures, so the surviving reading
is the trace reading.

\paragraph{Relation to empowerment.}
The channel maps actions into distributions over consequences, and the information term can be read from either
end.  Read backward, from consequence to action, it measures trace: how much the visible consequence reveals
about the realised action.  Read forward, from action to consequence, it measures how strongly the action
shapes the distribution of visible consequences.  Mutual information is symmetric, so the two readings share
one index.  No choice among lotteries separates them.

The forward reading has an exact counterpart.  Suppose $\E[u(O)\mid a]=\bar u$ for every action.  Since
$\lambda>0$, maximising $V$ over $q$ gives
\[
 \max_{q\in\Delta(A)}V(q,K)=\bar u+\lambda C(K),
 \qquad
 C(K):=\max_{q\in\Delta(A)}I_{qK}(A;O,R),
\]
the Shannon capacity of the channel.  When the visible consequence is read as the agent's sensor state, $C(K)$
is the one-step empowerment index of Klyubin, Polani, and Nehaniv \cite{KlyubinPolaniNehaniv2005}; see also
\cite{SalgeGlackinPolani2014}.  Under this supposition, the trace-tempered agent solves exactly the problem
that defines one-step empowerment: her optimal lotteries are the capacity-achieving action distributions, and
the size of $\lambda$ plays no role in the ranking---the price is visible only against payoffs.  Empowerment
summarises the channel by the optimised value $C(K)$; the representation recovers the ranking beneath it,
$q\mapsto I_{qK}(A;O,R)$, from preferences and prices it against material utility.

\paragraph{What the axioms preclude.}
The main exclusion is selective legibility.  Mutual information responds to an act's
probability and to the distinctiveness of its consequences, not to what the act is.  An agent who wants her generous
acts recognised and her mean acts hidden may therefore prefer a channel to its garbling in one case and the garbling to its source
in another.  Such sign-reversing preferences violate the global data-processing requirement in Axiom~\textup{(A6)} and fall
outside all three cases of the sign trichotomy.  Such preferences are not irrational; they are tastes for
something other than trace.  Wanting particular acts known points to an audience, outside the agent or within
her, for whom what is revealed matters.  The taste axiomatised here does not depend on one: it prices how much
the consequence reveals, not what.  A model of selective legibility would have to supply the audience, making
the value of evidence depend on its content and not only on the amount of uncertainty it removes.

\paragraph{The payoff--record distinction.}
Although the information term treats $O$ and $R$ jointly, the axioms cannot.  Axiom~\textup{(A6)} permits the
record to be garbled while the payoff is held fixed; this separates a change in evidence from a change in material
consequences.  Axiom~\textup{(A8)} permits records to arrive in stages while the payoff is delivered once.  Its
recordwise sure-thing restriction makes branchwise comparisons commensurable and supplies the cardinal structure used to
place payoff utility and information on a common scale.  Thus either $O$ or $R$ may reveal the action, but only
$O$ enters material utility.  The record is also assumed to have no utility beyond what it reveals: a record
that is itself pleasant or shameful belongs in $O$.  Moving the shame into $O$ is a change of description, not
of model: the same actions produce the same visible consequences with the same probabilities, so the
information term is unchanged.  Shame is a payoff; it cannot be counted again as negative trace.

\paragraph{One shot, with no audience.}
The model describes one episode.  It specifies neither an audience nor any response to the beliefs induced by
the record; the posterior is a statistic of the implemented lottery and the channel, not anyone's belief.  The
stages in Axiom~\textup{(A8)} factor a single one-shot interaction: no action follows the record, and no relationship
extends beyond the episode.  Now take the instrumental reading of the introduction literally: an observer
watches the visible consequence and pays the agent for what the consequence reveals about the action.  Let the
payment be proportional to the fall in the observer's uncertainty.  In expectation the payment is
$\lambda\,\I_{qK}(A;O,R)$ when the observer's prior is the implemented lottery $q$, and only that prior delivers
the match for every lottery and channel: the payment rewards what the consequence adds for someone who already
knows the plan.  An audience that correctly anticipates the plan, as it would in
equilibrium, pays $\lambda$ per unit of information.  The price is positive for the politician backed when
outcomes trace to her acts, and negative for the agent whose record will be used against her.  The audience
cares about how much is revealed, not which act it was, so selective legibility remains excluded.  The audience
is not a primitive: the representation predicts no change when an audience that pays nothing is added while the action
lottery and channel are held fixed.  If choices change, the model is missing the value of being observed by
others.

\clearpage
\appendix

\part*{Proofs}
\addcontentsline{toc}{part}{Proofs}

\section{Proof of the representation theorem}
\label{app:main-proof}

This appendix gives a proof of Theorem~\ref{thm:main}.  Its
dependency spine is simple.  We first derive the affine and behavioural tools
used throughout.  We then restrict attention to constant payoffs, align the
comparison scales across priors and support faces, and identify the resulting
trace value as mutual information.  Returning to the full payoff--record
domain, we put material utility and trace on one scale, prove the exact
reached-branch identity, and assemble an arbitrary channel by a finite
telescope and sequentialisation.


\subsection{Shared affine preliminaries}


The results of this subsection are bookkeeping: only
their statements are used later, and the reader may refer back as needed.


\begin{lemma}[Blocks, supports, and dummies]\label{lf:lem:bookkeeping}
Suppose Axioms \textup{(A1)} and \textup{(A5)--(A7)} hold, and recall that
$(q,K)\succeq(p,L)$ denotes the paper's canonical two-block comparison.
\begin{enumerate}[label=\textup{(\alph*)},leftmargin=*]
\item The canonical two-block comparison of experiments is complete and
  transitive.  It is unchanged by invertible relabellings of actions or
  records, and by replacement of either experiment with an indifferent one.
\item For every finite experiment $(q,K)$, write $S:=\supp q$, $q_S$ for $q$
  viewed on $S$, and $K_S$ for the restriction of $K$ to the rows in $S$.
  Then $(q_S,K_S)\sim(q,K)$.
\item For every finite experiment $(q,K)$, every nonempty finite $B$, and
  every $s\in\D(B)$, write $(q\otimes s)(a,b):=q(a)s(b)$ and
  $K^\uparrow(o,r\mid a,b):=K(o,r\mid a)$.  Then
  \begin{equation}
  (q\otimes s,K^\uparrow)\sim(q,K).
  \label{lf:eq:dummy-ordinal}
  \end{equation}
\end{enumerate}
\end{lemma}

\begin{proof}
Place any finite collection of experiments in one finite block environment.
Axiom~\textup{(A5)} removes unused blocks, and Axiom~\textup{(A1)} supplies
completeness and transitivity there.  Invertible action and record relabellings
are neutral by applying Axioms~\textup{(A6)} and~\textup{(A7)} in both
directions.  In particular, swapping both tags identifies the two orientations
of a two-block comparison.  This proves part \textup{(a)}; a three-block
environment gives transitivity and a four-block environment gives replacement.

For part \textup{(b)}, process an action in $S$ by its inclusion in $A$.  In the
other direction, output $a\in S$ when the realised action is $a$, and choose an
arbitrary element of $S$ off the support.  Both processors satisfy
\eqref{eq:action-processing}; rows attached to null processed actions may be
completed arbitrarily.  Axiom~\textup{(A7)} gives the two comparisons.

For part \textup{(c)}, first project $(a,b)$ to $a$.  Conversely, given $a$,
draw $b$ independently from $s$ and output $(a,b)$.  Again the two joint-law
identities in \eqref{eq:action-processing} are exact, so Axiom~\textup{(A7)}
applies in both directions.  Part \textup{(a)} completes both arguments.
\end{proof}


Axiom~\textup{(A8)} was stated in its weakest form, for a binary first
stage.  The proof uses its finite-branch form; the next lemma records the
equivalence and makes explicit where the other axioms enter it.

\begin{lemma}[Finite-branch extension]\label{pt:lem:finite-branch-extension}
Suppose Axioms~\textup{(A1)} and \textup{(A5)--(A7)} hold.
\begin{enumerate}[label=\textup{(\alph*)},leftmargin=*]
\item Axiom~\textup{(A8)} is equivalent to the following property.  Given
$k\ge1$, a first-stage record channel
$P:A\to\D(\{1,\ldots,k\})$, a prior $q\in\D(A)$, and two ordered lists of
continuation channels $(K_1,\ldots,K_k)$ and $(L_1,\ldots,L_k)$, with
$K_i,L_i:A\to\D(O\times R)$, define for every reached $i$
\[
 m_i:=\sum_aq(a)P(i\mid a),
 \qquad q_i(a):=\frac{q(a)P(i\mid a)}{m_i}.
\]
If every reached $i$ satisfies
\[
 (q_i,K_i)\succeq(q_i,L_i),
\]
then
\begin{equation}
 \bigl(q,P\triangleright(K_1,\ldots,K_k)\bigr)
 \succeq
 \bigl(q,P\triangleright(L_1,\ldots,L_k)\bigr),
 \label{pt:eq:finite-branch-extension}
\end{equation}
strictly if one reached branch comparison is strict.
\item Suppose in addition that Axiom~\textup{(A8)} holds.  Fix $q\in\D(A)$, a first-stage record channel
$P:A\to\D(\{1,\ldots,k\})$, a reached record $i^*$ with posterior $q_{i^*}$, and continuation channels
$K_1,\ldots,K_k,L_1,\ldots,L_k:A\to\D(O\times R)$ such that $K_i=L_i$ for every $i\ne i^*$.  Then
\[
 (q_{i^*},K_{i^*})\succeq(q_{i^*},L_{i^*})
 \quad\Longleftrightarrow\quad
 \bigl(q,P\triangleright(K_1,\ldots,K_k)\bigr)
 \succeq
 \bigl(q,P\triangleright(L_1,\ldots,L_k)\bigr),
\]
and the same equivalence holds for strict preference.
\end{enumerate}
\end{lemma}

\begin{proof}
Assume first Axiom~\textup{(A8)}.  Put
\[
 m_i:=\sum_aq(a)P(i\mid a),
 \qquad
 I^+:=\{i\in\{1,\ldots,k\}:m_i>0\}.
\]
If $m_i=0$, then $P(i\mid a)=0$ for every $a\in\supp q$.  Consequently, changing the continuation after such a
record changes the compound channel only on action rows outside $\supp q$.  Apply Axiom~\textup{(A7)} in both
directions using the identity action processor: its Bayesian pushforwards agree on supported rows and allow
arbitrary completions elsewhere.  Continuations following unreached records are therefore immaterial.  We may
hold them fixed and enumerate the reached records as $I^+=\{i_1,\ldots,i_n\}$.

If $n=1$, the first-stage record is constant on $\supp q$.  Deleting that
constant label and adjoining it again are payoff-preserving record processings
in both directions on every reached row; the off-support rows are immaterial
by the preceding paragraph.  Axiom~\textup{(A6)} therefore identifies the
compound comparison with the sole branch comparison.  Suppose henceforth that
$n\ge2$.

For $j=0,\ldots,n$, define hybrid continuation channels by
\[
 H_i^j:=
 \begin{cases}
 K_i,&i\in\{i_1,\ldots,i_j\},\\
 L_i,&i\in I^+\setminus\{i_1,\ldots,i_j\},
 \end{cases}
\]
and use a common arbitrary continuation on $\{1,\ldots,k\}\setminus I^+$.  Write
$E_j:=P\triangleright(H_1^j,\ldots,H_k^j)$.  Fix $j\ge1$ and collapse the first-stage record to the binary channel
\[
 B_j(1\mid a):=P(i_j\mid a),
 \qquad
 B_j(2\mid a):=\sum_{i\ne i_j}P(i\mid a).
\]
Both binary records are reached.  Put all binary continuations on the common record set
$\{1,\ldots,k\}\times R$.  Pad the
two continuations that may differ by
\[
 \widetilde K_j(o,(i,r)\mid a):=\1_{\{i=i_j\}}K_{i_j}(o,r\mid a),
 \qquad
 \widetilde L_j(o,(i,r)\mid a):=\1_{\{i=i_j\}}L_{i_j}(o,r\mid a).
\]
On record $2$, collect all unchanged branches in the continuation channel
\begin{equation}
 M_j(o,(i,r)\mid a)
 :=\frac{P(i\mid a)}{B_j(2\mid a)}H_i^j(o,r\mid a),
 \qquad i\ne i_j,
 \label{pt:eq:complement-continuation}
\end{equation}
with zero probability on records $(i_j,r)$, whenever $B_j(2\mid a)>0$, and complete the other rows arbitrarily.
Since $H_i^j=H_i^{j-1}$ for $i\ne i_j$, this continuation is common to the two comparisons.
Binary record $1$ has probability
$m_{i_j}$ and carries the conditional action lottery $q_{i_j}$; record $2$ carries the complementary conditional
lottery, and their mixture is $q$.  On the active record support, deleting or
adjoining the constant label $i_j$ gives payoff-preserving record processings
in both directions.  Axiom~\textup{(A6)} therefore identifies the comparison
of $\widetilde K_j$ and $\widetilde L_j$ with that of $K_{i_j}$ and $L_{i_j}$.
Hence Axiom~\textup{(A8)} gives
\[
 \bigl(q,B_j\triangleright(\widetilde K_j,M_j)\bigr)
 \succeq
 \bigl(q,B_j\triangleright(\widetilde L_j,M_j)\bigr),
\]
strictly whenever the comparison at $q_{i_j}$ is strict.

On the active record support, the nested records are in bijection with the
flat records through
\[
 (1,(i_j,r))\longmapsto(i_j,r),
 \qquad
 (2,(i,r))\longmapsto(i,r)\quad(i\ne i_j).
\]
Extend these maps arbitrarily over records of zero probability.  Together
with the preceding removal of off-support rows, they give payoff-preserving
record processings in both directions.  Axiom~\textup{(A6)} therefore
makes the two binary compounds indifferent to $E_j$ and $E_{j-1}$, respectively.  Thus
$(q,E_j)\succeq(q,E_{j-1})$, strictly at any strict replacement.  Place the finite collection
$E_0,\ldots,E_n$ in one block environment.  Axiom~\textup{(A5)} transports each pairwise comparison to that
environment, and transitivity in Axiom~\textup{(A1)} yields $(q,E_n)\succeq(q,E_0)$, strictly if at least one
step is strict.  Restoring the immaterial null branches proves \eqref{pt:eq:finite-branch-extension}.

Conversely, assume the displayed finite-branch property, including its strict clause.  In the binary environment
of Axiom~\textup{(A8)}, reflexivity and duplication coherence make the common continuation weakly comparable to
itself, so the finite-branch property gives the forward implication in
\eqref{eq:recordwise-sure-thing}.  If the aggregate comparison in that equation held while
$(q_1,K)\not\succeq(q_1,L)$, completeness would give $(q_1,L)\succ(q_1,K)$.  Applying the strict
finite-branch property to the reversed lists would give the strict reverse aggregate comparison, a contradiction.  This
proves the converse implication and hence Axiom~\textup{(A8)}.  This
establishes part~\textup{(a)}.

For part~\textup{(b)}, the forward implication follows from part~\textup{(a)}, using duplication coherence on
the common branches.  If the aggregate comparison held but the branch comparison failed, completeness would make
the reversed branch comparison strict; the strict clause of part~\textup{(a)} would then give a strict
reverse aggregate comparison.  The strict equivalence follows by applying the weak equivalence in both directions.
\end{proof}

In the appendices, references to Axiom~\textup{(A8)} with more than two first-stage records mean the
finite-branch property of Lemma~\ref{pt:lem:finite-branch-extension}\textup{(a)}; when only one branch
changes, we cite part~\textup{(b)}.

\subsection{Terminal laws and affine fibres}

For channels $E,G$ at a common prior $q$, write $E\succeq_q G$ for
$(q,E)\succeq(q,G)$.  Fix for the moment a full-support prior $q\in\D(A)$.  For a joint channel $K$,
write
\[
p_K(o,r):=\sum_a q(a)K(o,r\mid a)\;(=p_{qK}(o,r)),
\qquad
\pi^K_{or}(a):=\frac{q(a)K(o,r\mid a)}{p_K(o,r)}
\]
when $p_K(o,r)>0$, and define its \emph{marked terminal law}
\begin{equation}
\eta_K:=\sum_{(o,r):p_K(o,r)>0}
p_K(o,r)\,\delta_{(o,\pi^K_{or})},
\label{lf:eq:marked-law}
\end{equation}
where $\delta_x$ denotes the point mass at $x$.
It records the terminal payoff and the posterior about the realised action.
Write $I_q(E)$ for the mutual information between the action and the complete
visible output of $E$ under $q$.  When the output is a first-stage record $Y$,
we also write this quantity as $I_q(A;Y)$.

\begin{lemma}[Terminal laws and affine fibres]\label{lf:lem:terminal-affine}
Under Axioms \textup{(A1)}, \textup{(A2)}, and
\textup{(A5)--(A8)}, the following statements hold.
\begin{enumerate}[label=\textup{(\alph*)},leftmargin=*]
\item If $q$ has full support and $K,L$ are finite experiments at $q$ with
  $\eta_K=\eta_L$, then
  $(q,K)\sim(q,L)$.
\item At every full-support $q$, the quotient of finite experiments by their
  marked terminal laws has a representation $W_q$ of $\succeq_q$, affine in
  the marked terminal law.
\end{enumerate}
\end{lemma}

\begin{proof}
We give the finite Blackwell argument behind part \textup{(a)}.  For a payoff
$o$ and posterior $\pi$, let
\[
M_K(o,\pi):=
\sum_{\substack{r:\,p_K(o,r)>0\\ \pi^K_{or}=\pi}} p_K(o,r).
\]
Equality of marked laws gives $M_K(o,\pi)=M_L(o,\pi)=:M(o,\pi)$.
For $p_K(o,r)>0$, draw a record $s$ of $L$ according to
\begin{equation}
T(s\mid o,r):=
\1_{\{p_L(o,s)>0,\,\pi^L_{os}=\pi^K_{or}\}}
\frac{p_L(o,s)}{M(o,\pi^K_{or})}.
\label{lf:eq:matching-processor}
\end{equation}
Choose an arbitrary probability row when $p_K(o,r)=0$.  This is a
payoff-preserving record processor.  Bayes' formula gives
\[
q(a)K(o,r\mid a)=p_K(o,r)\pi^K_{or}(a).
\]
For a reached target record $s$, with posterior $\pi=\pi^L_{os}$, it follows
that
\[
\sum_r K(o,r\mid a)T(s\mid o,r)
=\frac{p_L(o,s)}{M(o,\pi)}
  \frac{\pi(a)}{q(a)}M(o,\pi)
=L(o,s\mid a).
\]
If $s$ is unreached, full support makes both sides zero.  Thus postprocessing
$K$ by $T$ gives $L$.  The symmetric construction gives a
processor from $L$ to $K$.  Axiom~\textup{(A6)} and
Lemma~\ref{lf:lem:bookkeeping}\textup{(a)} give indifference.

Boundary priors are handled later by first restricting to their positive
support using Lemma~\ref{lf:lem:bookkeeping}\textup{(b)}.

For part \textup{(b)}, take the quotient of finite experiments at $q$ by
equality of \eqref{lf:eq:marked-law}.  If $E$ and $G$ are two such experiments,
their public mixture $tE\oplus(1-t)G$ draws a coin independent of the action,
selecting $E$ with probability $t$ and $G$ otherwise, runs the selected
experiment, and retains the coin in the record.  Its marked law is
\begin{equation}
\eta_{tE\oplus(1-t)G}=t\eta_E+(1-t)\eta_G.
\label{lf:eq:public-mixture}
\end{equation}
Axiom~\textup{(A8)} gives, for $0<t<1$ and any third experiment $L$,
\begin{equation}
E\succeq_q L
\quad\Longleftrightarrow\quad
tE\oplus(1-t)G\succeq_q tL\oplus(1-t)G.
\label{lf:eq:mixture-independence}
\end{equation}
This is the action-independent special case of the axiom: the public coin is the binary first-stage record and
$G$ is the common continuation.  If record alphabets differ, put them in a common disjoint union before invoking
the axiom; part \textup{(a)} identifies the result with
\eqref{lf:eq:public-mixture}.

It remains only to check the Archimedean step.  A closed segment between two
experiments, and any fixed comparator, can be realised on one fixed union of
their finite record alphabets.  Its channel entries vary continuously with
$t\in[0,1]$.
Axiom~\textup{(A2)} therefore makes both contour sets along the segment closed.
Completeness says that they cover $[0,1]$; when a target lies between the two
endpoints, the sets contain opposite endpoints and connectedness makes them
intersect.  Thus every such target is indifferent to a point of the segment.
The mixture-space theorem \cite{HersteinMilnor1953} now supplies an affine
representative $W_q$.  This
quotient is closed under every mixture and embedding used below.
\end{proof}

\subsection{Fixed-channel and environment relevance}

The relevance axioms in the main text are stated inside fixed channels.  The
proofs below use their equivalent comparison-environment forms.

For $o\in O$, write $K_o:\{\ast\}\to\D(O\times\{\ast\})$ for the channel that
delivers $o$ for sure.  On a finite action set $A$, write
$K^{\mathrm{id}}_o:A\to\D(O\times A)$ for the channel that delivers $o$ and
reports the action, and $K^0_o:A\to\D(O\times\{\ast\})$ for the channel that
delivers $o$ and no record.

\begin{lemma}[Fixed-channel and environment relevance]\label{lem:relevance-bridge}
Given Axioms~\textup{(A1)}, \textup{(A5)}, \textup{(A6)},
\textup{(A7)}, and \textup{(A8)}, the following equivalences hold.
\begin{enumerate}[label=\textup{(\alph*)},leftmargin=*]
\item Axiom~\textup{(A3)} is equivalent to the existence of
  $o^+,o^-\in O$ such that
  \[
    (\delta_\ast,K_{o^+})\succ(\delta_\ast,K_{o^-}).
  \]
\item Axiom~\textup{(A4)} is equivalent to the following condition: there is
  an outcome $o_\ast\in O$ such that, for every finite $A$ with $|A|\ge2$
  and every full-support $q\in\D(A)$,
  \[
    (q,K^{\mathrm{id}}_{o_\ast})\succ(q,K^0_{o_\ast}).
  \]
\end{enumerate}
\end{lemma}

\begin{proof}
For part \textup{(a)}, action processing in both directions identifies each
pure lottery in Axiom~\textup{(A3)} with the corresponding singleton
experiment; the unrestricted row at the unreached action supplies the reverse
pushforward.  Replacement transports the strict comparison both ways.

For part \textup{(b)}, restrict the witnessing lotteries to their supports.
Two-way action processing identifies them with fair binary full revelation
and a channel with two identical rows; two-way record processing identifies
the latter with silence.  Thus Axiom~\textup{(A4)} is equivalent to
\begin{equation}
  ((\tfrac12,\tfrac12),K^{\mathrm{id}}_{o_\ast})
  \succ
  ((\tfrac12,\tfrac12),K^0_{o_\ast})
  \label{eq:bridge-fair-binary}
\end{equation}
on a binary action set.

For any full-support binary $p$, choose
$0<\varepsilon\le2\min_i p_i$ and a binary record channel $P$ with
$P(y\mid i)=\varepsilon/(2p_i)$ for a record $y$.  It has mass $\varepsilon$ and fair
posterior.  Reveal only after $y$ in one continuation plan and remain silent
everywhere in the other.  Equation~\eqref{eq:bridge-fair-binary} and
Lemma~\ref{pt:lem:finite-branch-extension}\textup{(b)} rank the compounds strictly.  Full
revelation garbles to the first and the second garbles to silence, so
$(p,K^{\mathrm{id}}_{o_\ast})\succ(p,K^0_{o_\ast})$.

For general full-support $q$ on $|A|\ge2$, report a two-cell partition.
Two-way action processing identifies this with binary full revelation and
identifies silence with binary silence.  Hence the cell report is strictly
better than silence; full revelation weakly dominates the cell report by
record processing.  The converse takes the fair binary prior and reverses
the same support identifications.
\end{proof}

\subsection{The pure-trace core}
\label{pt:sec:environment}\label{lf:app:pure-trace}

Whenever Axiom~\textup{(A4)} is imposed below, fix one of its witnessing
outcomes and denote it by $o_\ast$.  Every channel between nonempty finite sets
$P:A\to\D(X)$ has the constant-payoff lift
\begin{equation}
K_P(o,x\mid a):=\1_{\{o=o_\ast\}}P(x\mid a).
\label{pt:eq:constant-payoff-lift}
\end{equation}
Throughout this constant-payoff phase, comparisons of pure record channels mean comparisons
of these lifts in the model already defined in the main text:
\[
q\succeq_Pq'\quad:\Longleftrightarrow\quad q\succeq_{K_P}q'.
\]
Likewise, $(q,P)\succeq(r,Q)$ denotes the main-text two-block comparison of
the lifts, and $P\sqcup Q$, processing, continuation lists, public mixtures
$tP\oplus(1-t)Q$, and $P_1\triangleright(Q_1,\ldots,Q_k)$ denote the lifts of
the corresponding earlier constructions.
The lift commutes with them up to tagged or singleton-record bijections,
neutral by Axiom~\textup{(A6)} in both directions.  These are only
abbreviations, not a new primitive relation or behavioural condition.

Two pieces of notation will be used repeatedly.  If
$S:A\to\Delta(B)$ is an action processor, let $S^qP:B\to\Delta(X)$ denote any
Bayesian pushforward satisfying
\[
 (qS)(b)(S^qP)(x\mid b)
 =\sum_a q(a)S(b\mid a)P(x\mid a).
\]
Its rows at zero-mass reports are arbitrary, exactly as in
\eqref{eq:action-processing}.  For channels $P:A\to\Delta(X)$ and
$Q:B\to\Delta(Y)$, their independent product is
\[
 (P\otimes Q)((x,y)\mid(a,b)):=P(x\mid a)Q(y\mid b).
\]

For $q\in\D(A)$ and $P:A\to\D(X)$, write
\[
m_{qP}(x):=\sum_a q(a)P(x\mid a),
\qquad
r_x^{qP}(a):=\frac{q(a)P(x\mid a)}{m_{qP}(x)}
\quad\text{when }m_{qP}(x)>0,
\]
and define the finite posterior law
\[
\mu_{qP}:=
\sum_{x:m_{qP}(x)>0}m_{qP}(x)\,\delta_{r_x^{qP}}.
\]
In subscripts, a comma separates the prior from the channel whenever either
is compound or decorated.
Let $\Id_A$ be full revelation and $U_A$ the one-record uninformative
channel.  Call a prior, and the fibre at it, \emph{nondegenerate} if its
support contains at least two actions---equivalently, if its entropy is
positive; a finite action set is nondegenerate if it has at least two
elements.  With logarithms in the fixed base used in the main text, put
\[
I_{qP}(A;X)
:=\Sh(m_{qP})-\sum_a q(a)\Sh(P(\cdot\mid a))
=\Sh(q)-\int_{\D(A)}\Sh(r)\,d\mu_{qP}(r).
\]

\begin{proposition}[Pure-trace lemma]\label{pt:prop:main}
Suppose the preference family satisfies main-text Axioms
\textup{(A1)}, \textup{(A2)}, and \textup{(A4)--(A8)}.  Then, for every record channel between
nonempty finite sets
$P:A\to\D(X)$ and every $q,q'\in\D(A)$,
\[
q\succeq_Pq'
\quad\Longleftrightarrow\quad
I_{qP}(A;X)\ge I_{q'P}(A;X).
\]
The same numerical scale represents every block-supported comparison: for
every nonempty finite family $\{P_i:A_i\to\D(X_i)\}_{i\in J}$ of channels
between nonempty finite sets, distinct $i,j\in J$,
and priors $q_i\in\D(A_i)$ and $q_j\in\D(A_j)$,
\[
q_i^i\succeq_{\bigsqcup_{k\in J}P_k}q_j^j
\quad\Longleftrightarrow\quad
I_{q_iP_i}(A_i;X_i)\ge I_{q_jP_j}(A_j;X_j).
\]
\end{proposition}

The rest of this subsection proves the proposition.  Assume its hypotheses;
every working property below is the corresponding Axiom~\textup{(A1)},
\textup{(A2)}, or \textup{(A4)--(A8)} applied to constant-payoff lifts.

\paragraph{Proof strategy.}
First, marked terminal laws give affine fixed-prior values, made exact under
relabelling and support transport.  Second, product and reached-branch
calibrations meet; their two-groupings argument eliminates the possible
product interaction and aligns every scale.  Third, the exact chain rule
produces a symmetric, continuous, strongly additive uncertainty functional,
identified as Shannon entropy by Faddeev's theorem.  The product-coefficient
and branch tangent-space lemmas may be read as technical interfaces.

\subsubsection{Fixed-prior affine values}

For a finite action set $A$, let $\mathsf X_A$ be the set of finitely
supported probability measures on $\Delta(A)$, and let
\[
 b(\mu):=\int_{\Delta(A)}s\,d\mu(s),
 \qquad
 \mathcal M_q:=\{\mu\in\mathsf X_A:b(\mu)=q\}.
\]
Thus $\mathcal M_q$ is the mixture set of finite Bayes-plausible posterior
laws with prior $q$, and $\mu_{qP}\in\mathcal M_q$ for every finite channel
$P$.  On the constant-payoff face, the marked terminal law from
\eqref{lf:eq:marked-law} is simply
\begin{equation}
 \eta_{K_P}=\delta_{o_\ast}\otimes\mu_{qP}.
 \label{pt:eq:pure-marked-law}
\end{equation}

Let $\delta_q\in\mathcal M_q$ be the posterior law of silence and let
\[
 \chi_q:=\sum_{a\in A}q(a)\delta_{e_a}\in\mathcal M_q
\]
be the posterior law of full revelation, where $e_a\in\D(A)$ is the point mass at $a$.

\begin{lemma}[Pure affine fibres]
\label{pt:lem:pure-affine}
\label{pt:lem:blockcoh}
\label{pt:lem:blackwell}\label{pt:lem:plsuff}
\label{pt:lem:publiccoin}
\label{pt:lem:postsep}
\label{pt:lem:convnorm}
\emph{Global order facts.}
Pure prior--channel pairs form a weak order with replacement, unaffected by
adding or deleting other blocks.  For every channel $P$ and priors $q,r$ on
its action set,
\begin{equation}
 q\succeq_P r
 \quad\Longleftrightarrow\quad
 (q,P)\succeq(r,P),
 \label{pt:eq:within-to-pair}
\end{equation}
and on fixed alphabets the pair-order graph is closed under entrywise
convergence of channels and priors.  \emph{Fixed full-support fibre.}
Fix a full-support $q\in\Delta(A)$.
Then the following statements hold.

\begin{enumerate}[label=\textup{(\alph*)},leftmargin=*]
\item If $\mu_{qP}=\mu_{qQ}$, record processors exist in both directions,
  so $(q,P)\sim(q,Q)$ and either channel may replace the other in a block
  comparison.  Conversely, every
  $\mu=\sum_{k=1}^n\lambda_k\delta_{s_k}\in\mathcal M_q$, written with
  $\lambda_k>0$, is implemented by
  \begin{equation}
   C_\mu(k\mid a):=\frac{\lambda_ks_k(a)}{q(a)},
   \qquad \mu_{qC_\mu}=\mu.
   \label{pt:eq:canonical-implementation}
  \end{equation}
  Hence the fixed-$q$ order descends to all of $\mathcal M_q$.

\item For publicly tagged mixtures,
  \begin{equation}
   \mu_{q,tP\oplus(1-t)R}
   =t\mu_{qP}+(1-t)\mu_{qR},
   \label{pt:eq:public-posterior-mixture}
  \end{equation}
  and, for $0<t<1$,
  \[
   (q,P)\succeq(q,Q)
   \quad\Longleftrightarrow\quad
   (q,tP\oplus(1-t)R)\succeq(q,tQ\oplus(1-t)R).
  \]

\item There is an affine $F_q:\mathcal M_q\to\mathbb R$ representing the
  fixed-prior order:
  \begin{equation}
   (q,P)\succeq(q,Q)
   \quad\Longleftrightarrow\quad
   F_q(\mu_{qP})\ge F_q(\mu_{qQ}).
   \label{pt:eq:fibre-representation}
  \end{equation}
  It may be normalised by
\begin{equation}
 F_q(\delta_q)=0,
 \qquad
 F_q(\mu)\ge0\quad(\mu\in\mathcal M_q),
 \label{pt:eq:zero-minimum}
\end{equation}
and, whenever $|A|\ge2$,
\begin{equation}
 F_q(\chi_q)=1.
 \label{pt:eq:full-revelation-anchor}
\end{equation}
These two anchors select a unique affine representative on every
nondegenerate full-support fibre.  On a singleton fibre we retain
$F_q\equiv0$.
\end{enumerate}
\end{lemma}

\begin{proof}
The order, replacement, and closed-graph assertions are respectively
Lemma~\ref{lf:lem:bookkeeping}\textup{(a)}, the duplication clause of
Axiom~\textup{(A5)}, and Axiom~\textup{(A2)}, all specialised to the lifts.

For part \textup{(a)}, equation~\eqref{pt:eq:pure-marked-law} turns equality
of posterior laws into equality of marked terminal laws.  The matching
processors in Lemma~\ref{lf:lem:terminal-affine}\textup{(a)} therefore give
record processors in both directions, including arbitrary completions on
zero-mass rows.  Replacement gives the quotient order.  Conversely, Bayes
plausibility gives $\sum_k\lambda_ks_k=q$, so the rows in
\eqref{pt:eq:canonical-implementation} sum to one and Bayes' formula gives
$\mu_{qC_\mu}=\mu$.

For part \textup{(b)}, Bayes' formula gives
\eqref{pt:eq:public-posterior-mixture}.  On constant-payoff lifts this is the
public-mixture identity~\eqref{lf:eq:public-mixture}, and the independence
argument in Lemma~\ref{lf:lem:terminal-affine}\textup{(b)} gives the
biconditional.  Tagged disjoint unions handle unequal record alphabets, with
part \textup{(a)} making the padding neutral.

For part \textup{(c)}, restrict the affine representative from
Lemma~\ref{lf:lem:terminal-affine}\textup{(b)} to laws
$\delta_{o_\ast}\otimes\mu$ and use parts \textup{(a)--(b)} and canonical
implementability.  If $|A|\ge2$, Lemma~\ref{lem:relevance-bridge}\textup{(b)}
makes the order nonconstant, so positive-affine uniqueness applies.  Deleting
the record garbles every channel to silence; hence $\delta_q$ is a minimum.
Translate its value to zero and divide by the positive value of $\chi_q$.
This gives the two anchors and nonnegativity.  If $|A|=1$, Bayes plausibility
makes $\mathcal M_q=\{\delta_q\}$.
\end{proof}

The only continuity used here is the finite-alphabet segment continuity
already verified in Lemma~\ref{lf:lem:terminal-affine}: implement the finitely
many laws on one tagged record alphabet and apply Axiom~\textup{(A2)} along
their mixtures.  No topology on laws with growing supports, and no pointwise
integrand on $\Delta(A)$, is being assumed.

\subsubsection{Exact transport across priors and support faces}

Action processing handles relabelling, undetectable distinctions, and
restriction to the positive support in one step.

For a bijection $\pi:A\to A'$, write $q\pi$ for the transported prior and
$\pi_\#\mu$ for the posterior law obtained by transporting every atom of
$\mu$.  If $B=\operatorname{supp}(q)$, write $q_B$ for the full-support
restriction of $q$ to $B$, $P_B$ for the restricted channel, and
$\pi_B\mu$ for the law obtained by restricting every posterior atom to $B$.

\begin{lemma}[Relabelling, undetectable distinctions, and support transport]
\label{pt:lem:transport}\label{pt:lem:actionbase}\label{pt:lem:neutrality}\label{pt:lem:supprestrict}
The canonically normalised fibre values have the following properties;
superscripts record the action set.

\begin{enumerate}[label=\textup{(\alph*)},leftmargin=*]
\item \emph{Exact relabelling.}
If $q$ has full support and $\pi:A\to A'$ is a bijection, then
\begin{equation}
 F_{q\pi}^{A'}(\pi_\#\mu)=F_q^A(\mu)
 \qquad(\mu\in\mathcal M_q).
 \label{pt:eq:exact-relabel}
\end{equation}

\item \emph{Undetectable distinctions.}
Let $S:A\to B$ be an onto deterministic map, let $q$ have full support, and
suppose that $P:A\to\Delta(X)$ is constant on every fibre of $S$.  Then
\begin{equation}
 (q,P)\sim(qS,S^qP).
 \label{pt:eq:undetectable}
\end{equation}
In particular, if $G_S$ reports only $S(a)$, then
$(q,G_S)\sim(qS,\Id_B)$.

\item \emph{Support restriction.}
For an arbitrary prior $q\in\Delta(A)$ with nonempty support
$B:=\operatorname{supp}(q)$ and every pair of finite channels
$P:A\to\Delta(X)$ and $Q:A\to\Delta(Y)$,
\begin{equation}
 (q,P)\sim(q_B,P_B).
 \label{pt:eq:support-indifference}
\end{equation}
Consequently comparisons at $q$ depend only on the rows indexed by $B$, and
\begin{equation}
 (q,P)\succeq(q,Q)
 \quad\Longleftrightarrow\quad
 (q_B,P_B)\succeq(q_B,Q_B).
 \label{pt:eq:support-comparison}
\end{equation}
Every atom of a law $\mu\in\mathcal M_q$ is supported on $B$, and the
definition
\begin{equation}
 F_q^A(\mu):=F_{q_B}^{B}(\pi_B\mu)
 \label{pt:eq:boundary-representative}
\end{equation}
represents all continuation comparisons at the boundary prior $q$.
\end{enumerate}
\end{lemma}

\begin{proof}
\emph{Part (a).}
Axiom~\textup{(A7)} applied to $\pi$ and $\pi^{-1}$ identifies the two
orders.  Hence $\mu\mapsto F_{q\pi}^{A'}(\pi_\#\mu)$ and $F_q^A$ are affine
representatives of the same order.  Relabelling preserves silence and full
revelation, so the two anchors and affine uniqueness give
\eqref{pt:eq:exact-relabel}; singleton values vanish.

\emph{Part (b).}
One direction of \eqref{pt:eq:undetectable} is Axiom~\textup{(A7)}.  For the
reverse, use the Bayesian refinement kernel
\[
 R(a\mid b):=
 \begin{cases}
 q(a)/(qS)(b),&S(a)=b,\\
 0,&S(a)\ne b.
 \end{cases}
\]
It satisfies $(qS)R=q$, and fibre constancy means that pushing $S^qP$ back
through $R$ recovers $P$.  A second application of Axiom~\textup{(A7)} gives
the reverse comparison.  Taking $P=G_S$ gives the final assertion.

\emph{Part (c).}
Lemma~\ref{lf:lem:bookkeeping}\textup{(b)} gives
\eqref{pt:eq:support-indifference}, and replacement gives
\eqref{pt:eq:support-comparison}.  If
$\mu=\sum_k\lambda_k\delta_{s_k}\in\mathcal M_q$ with $\lambda_k>0$, then
$0=q(a)=\sum_k\lambda_ks_k(a)$ off $B$; nonnegativity puts every $s_k$ on
$\Delta(B)$.  Thus $\pi_B\mu\in\mathcal M_{q_B}$, and the full-support
representation on $B$ proves~\eqref{pt:eq:boundary-representative}.
Compatibility with relabelling follows from part~\textup{(a)}.
\end{proof}

Henceforth action-set superscripts are suppressed, and boundary subscripts
mean transport from the positive support.

\paragraph{Normalisation ledger.}
$F_q$ has the silence/full-revelation normalisation; product calibration gives
$G_q=c_qF_q$ with a possible interaction $\kappa$; branch calibration gives
$V_q=G_q/a_q$.  Compatibility will make $a_q$ universal and $\kappa=0$.

\subsubsection{Scale coherence}

\paragraph{Product calibration and a common comparison scale.}

Write $P\succeq_qQ$ for $(q,P)\succeq(q,Q)$, and similarly for strict
preference.  For a channel $P:A\to\Delta(X)$, let $\widetilde P$ denote its
lift to $A\times B$, constant in the $B$-coordinate; for
$R:B\to\Delta(Y)$, let $\widetilde R$ denote the symmetric lift, constant in
the $A$-coordinate.
By Lemma~\ref{pt:lem:plsuff}, null records and record permutations are neutral.

\begin{lemma}[Cancellation of a common independent background]
\label{pt:lem:background}
Let $p\in\Delta(A)$ and $r\in\Delta(B)$ have full support.  For channels
$P,Q$ on $A$ and a channel $R$ on $B$,
\begin{equation}
 P\otimes R\succeq_{p\otimes r}Q\otimes R
 \quad\Longleftrightarrow\quad
 P\succeq_p Q.
\tag{BI}\label{pt:eq:backgroundinert}
\end{equation}
The symmetric assertion holds with the two coordinates interchanged.  The
equivalence, including its strict form, remains valid when a posterior reached
in the background coordinate lies on a proper support face.
\end{lemma}

\begin{proof}
Pad $P$ and $Q$ by null records so that they have the same record alphabet.
Regard the lift $\widetilde R$ of $R$ as a first-stage experiment on
$A\times B$, followed on every reached branch by $\widetilde P$, respectively
$\widetilde Q$.  Up to the harmless permutation of the two record
coordinates, the resulting compounds are $P\otimes R$ and $Q\otimes R$.

If $y$ is reached under $R$, its posterior on $A\times B$ is
\[
 p\otimes r_y,
 \qquad r_y:=r_y^{\,r,R}.
\]
The coordinate projection $A\times\operatorname{supp}(r_y)\to A$ and
Lemma~\ref{pt:lem:transport}\textup{(b)--(c)} give
\[
 (p\otimes r_y,\widetilde P)\sim(p,P),
 \qquad
 (p\otimes r_y,\widetilde Q)\sim(p,Q).
\tag{BI$'$}\label{pt:eq:bg-branch-transport}
\]
Thus every reached branch has the same weak comparison as $P$ versus $Q$ at
$p$.  If $P\succeq_pQ$, Lemma~\ref{pt:lem:finite-branch-extension}\textup{(a)} applied to
these branches proves the forward implication in~\eqref{pt:eq:backgroundinert}.

For the converse, suppose that $P\not\succeq_pQ$.  Completeness gives
$Q\succ_pP$.  By~\eqref{pt:eq:bg-branch-transport}, the reversed comparison is
strict on every reached branch; at least one such branch has positive mass.
The strict part of Lemma~\ref{pt:lem:finite-branch-extension}\textup{(a)} therefore gives
$Q\otimes R\succ_{p\otimes r}P\otimes R$, contradicting
$P\otimes R\succeq_{p\otimes r}Q\otimes R$.  This proves the converse and
also records why weak branch monotonicity alone would not suffice.  Interchange
the coordinates for the symmetric assertion.
\end{proof}

For finite posterior laws define
\[
 \left(\sum_i m_i\delta_{u_i}\right)\otimes
 \left(\sum_j n_j\delta_{v_j}\right)
 :=\sum_{i,j}m_i n_j\delta_{u_i\otimes v_j}.
\]
Bayes' rule gives
\begin{equation}
 \mu_{p\otimes r,P\otimes R}=\mu_{pP}\otimes\mu_{rR}.
\tag{PL}\label{pt:eq:productlaw}
\end{equation}

\begin{lemma}[Product gauge]
\label{pt:lem:coherentnorm}\label{pt:cor:permutationinvariance}
There are positive rescalings $G_q=c_qF_q$ of the zero-normalised affine
representatives and a single number $\kappa\in\mathbb R$ such that
\begin{equation}
 G_{p\otimes r}(\mu\otimes\nu)
 =G_p(\mu)+G_r(\nu)+\kappa G_p(\mu)G_r(\nu)
\tag{QA}\label{pt:eq:QA}
\end{equation}
for all full-support priors $p,r$ and all
$\mu\in\mathcal M_p$, $\nu\in\mathcal M_r$.  These rescalings can be chosen
exactly invariant under bijections.  If a factor is a singleton, its
representative is zero and the same formula holds.  Finally,
\begin{equation}
 1+\kappa G_r(\nu)>0
\tag{POS}\label{pt:eq:QAslope}
\end{equation}
whenever the $r$-fibre is nonconstant.
\end{lemma}

\begin{proof}
We give the algebra because both the common value of $\kappa$ and the strict
inequality~\eqref{pt:eq:QAslope} are used later.  Initially retain the notation
$F_q$.  For nondegenerate $p,r$, put
\[
 L_{p,r}(\mu,\nu):=F_{p\otimes r}(\mu\otimes\nu),
 \qquad x:=F_p(\mu),\quad y:=F_r(\nu).
\]
The map $L_{p,r}$ is separately affine, and
Lemma~\ref{pt:lem:background} says that each slice represents the appropriate
factor order.  We record the short uniqueness argument.  Fixing $\nu$ gives
$L_{p,r}=\alpha(\nu)x+\gamma(\nu)$ with $\alpha(\nu)>0$.  At
$\mu=\delta_p$, affine uniqueness gives
$\gamma(\nu)=B_{p,r}y$ for some $B_{p,r}>0$.  At any $\bar\mu$ with
$\bar x:=F_p(\bar\mu)\ne0$, the other slice has the form
$A_{p,r}\bar x+D_{p,r}y$ with $A_{p,r}>0$.  Comparing the two expressions
gives $\alpha(\nu)=A_{p,r}+C_{p,r}y$, where
$C_{p,r}:=(D_{p,r}-B_{p,r})/\bar x$.  Hence
\begin{equation}
 L_{p,r}(\mu,\nu)
 =A_{p,r}x+B_{p,r}y+C_{p,r}xy.
\label{pt:eq:pairbilinear}
\end{equation}
The symmetric calculation and strict slice cancellation give
\begin{equation}
 A_{p,r}+C_{p,r}y>0,
 \qquad B_{p,r}+C_{p,r}x>0
\label{pt:eq:pair-slopes}
\end{equation}
throughout the value ranges.

Let $s$ be a third nondegenerate prior, $\xi\in\mathcal M_s$, and
$z:=F_s(\xi)$.  Exact relabelling makes the evaluations under the two
associations of a triple product equal.  Comparing the coefficients of $x,y,z$ gives
\begin{align}
 A_{p\otimes r,s}A_{p,r}&=A_{p,r\otimes s},\label{pt:eq:Acoc1}\\
 A_{p\otimes r,s}B_{p,r}&=B_{p,r\otimes s}A_{r,s},\label{pt:eq:Acoc2}\\
 B_{p\otimes r,s}&=B_{p,r\otimes s}B_{r,s}.
\label{pt:eq:Acoc3}
\end{align}
The coordinate swap similarly gives
\[
 B_{r,p}=A_{p,r},\qquad A_{r,p}=B_{p,r},\qquad C_{r,p}=C_{p,r}.
\tag{SW}\label{pt:eq:coefficient-swap}
\]
This coefficient comparison is legitimate: the laws
$(1-t)\delta_p+t\chi_p$ make $x$ range over a nondegenerate interval, and
the same holds for $y$ and $z$.
Let $\rho(p,r):=B_{p,r}/A_{p,r}$.  Dividing
\eqref{pt:eq:Acoc2} by~\eqref{pt:eq:Acoc1} yields
\[
 \rho(p,r)=\rho(p,r\otimes s)A_{r,s}.
\]
Fix a nondegenerate reference prior $q_0$ and set
$g(r):=\rho(q_0,r)$.  Then
\begin{equation}
 A_{r,s}=\frac{g(r)}{g(r\otimes s)}.
\label{pt:eq:gaugeA}
\end{equation}
The uniqueness of the coefficients in~\eqref{pt:eq:pairbilinear}, applied
after product relabellings, shows that $g$ is invariant under bijections.
Consequently the positive rescaling $G_p:=c_pF_p$, with $c_p:=g(p)$,
preserves exact relabelling.  Equation~\eqref{pt:eq:gaugeA} makes the new first coefficient
equal to one, and~\eqref{pt:eq:coefficient-swap} makes the second coefficient
equal to one as well.  We henceforth use $G$ for these gauged representatives.

Write the remaining coefficient as $\kappa_{p,r}$.  Associativity of
\[
 x\oplus_{p,r}y:=x+y+\kappa_{p,r}xy
\]
and comparison of the coefficients of $xy,xz,yz$ give
\[
 \kappa_{p,r}=\kappa_{p,r\otimes s},\qquad
 \kappa_{p\otimes r,s}=\kappa_{p,r\otimes s},\qquad
 \kappa_{p\otimes r,s}=\kappa_{r,s}.
\tag{KA}\label{pt:eq:kappa-assoc}
\]
(The $xyz$ coefficient gives the product of the corresponding two sides and
adds no restriction.)  Taking $s=q_0$ in~\eqref{pt:eq:kappa-assoc} shows that
$\kappa_{p,r}=\kappa_{r,q_0}$, so it is independent of $p$; the swap in
\eqref{pt:eq:coefficient-swap} makes it symmetric, hence independent of $r$.
Call the common value $\kappa$.  After the gauge,
\eqref{pt:eq:pair-slopes} is exactly
$1+\kappa G_r(\nu)>0$ and its symmetric counterpart.

For a singleton prior set $G\equiv0$; the canonical bijections
$A\cong A\times\{*\}$ give~\eqref{pt:eq:QA}.  For a boundary prior, define
$G_r$ by exact transport from its positive-probability support, as in
Lemma~\ref{pt:lem:supprestrict}.  This convention is relabelling invariant.
\end{proof}

\begin{lemma}[Cross-prior block representation]
\label{pt:lem:blockbridge}
For arbitrary finite priors $q\in\Delta(A)$ and $r\in\Delta(B)$ and channels
$P,Q$ on the respective action sets,
\begin{equation}
 q^0\succeq_{P\sqcup Q}r^1
 \quad\Longleftrightarrow\quad
 G_q(\mu_{qP})\ge G_r(\mu_{rQ}).
\tag{BC}\label{pt:eq:blockcomp}
\end{equation}
\end{lemma}

\begin{proof}
First suppose that $q,r$ have full support.  By~\eqref{pt:eq:QA} and
$G_q(\delta_q)=G_r(\delta_r)=0$,
\[
 G_{q\otimes r}(\mu_{qP}\otimes\delta_r)=G_q(\mu_{qP}),
 \qquad
 G_{q\otimes r}(\delta_q\otimes\mu_{rQ})=G_r(\mu_{rQ}).
\]
The two laws on the left lie in the same fibre and are implemented by
$P\otimes U_B$ and $U_A\otimes Q$, so the affine fibre representation
compares them by the displayed numbers.

Lemma~\ref{lf:lem:bookkeeping}\textup{(c)} removes the independent dummy
coordinates, giving
$(q\otimes r,P\otimes U_B)\sim(q,P)$ and
$(q\otimes r,U_A\otimes Q)\sim(r,Q)$.  Replacement proves
\eqref{pt:eq:blockcomp}.  Arbitrary priors follow by support restriction;
singleton supports use the zero convention.
\end{proof}

\paragraph{Branch calibration and coherent face scales.}

Changing a continuation on one reached branch changes the posterior law by a
signed measure with zero total mass and zero barycentre.  The following space
is exactly the space of such feasible directions.
For a nonempty finite set $A$, let
\[
 W_A:=\left\{\eta:
 \begin{array}{l}
 \eta\text{ is a finitely supported signed measure on }\Delta(A),\\[-2pt]
 \eta(\Delta(A))=0,\quad \displaystyle\int p\,d\eta(p)=0
 \end{array}\right\}.
\]
If $\iota:B\hookrightarrow A$ is an injection, let
$i_\iota:W_B\to W_A$ push every posterior forward along $\iota$; for a
literal subset $B\subseteq A$, write $i_B$ for the inclusion case.  Let
$L_q^A$ be the linear part of the affine functional $G_q$ on the affine hull
of $\mathcal M_q$.

\begin{lemma}[Branch aggregation]
\label{pt:lem:branchagg}
Let $q$ have full support on $A$.  For every $r\in\Delta(A)$ with
$B:=\operatorname{supp}(r)$ and $|B|\ge2$, there is a unique
$\beta_A(q,r)>0$ satisfying
\begin{equation}
 L_q^A\circ i_B=\beta_A(q,r)L_{r_B}^B
 \quad\hbox{on }W_B.
\label{pt:eq:branch-linear}
\end{equation}
It depends only on $q,r$ and the face inclusion.  If $k\ge1$ and
$P_1:A\to\Delta(\{1,\ldots,k\})$, define
\begin{equation}
 m_i:=\sum_aq(a)P_1(i\mid a),\qquad
 r_i(a):=\frac{q(a)P_1(i\mid a)}{m_i}\quad(m_i>0).
\label{pt:eq:branch-notation}
\end{equation}
Here $r_i$ is the branch posterior, written $q_i$ in the main text.
Then every list of finite continuation channels $Q_1,\ldots,Q_k$ on $A$
satisfies
\begin{equation}
 G_q\bigl(\mu_{q,P_1\triangleright(Q_1,\ldots,Q_k)}\bigr)
 =G_q(\mu_{qP_1})+
 \sum_{i:m_i>0}m_i\beta_A(q,r_i)
 G_{r_i}(\mu_{r_i,Q_i}).
\tag{BA}\label{pt:eq:branchaggregation}
\end{equation}
For singleton $r_i$, the continuation value is zero and the corresponding
coefficient may be fixed by any positive convention.
\end{lemma}

\begin{proof}
We first record the common tangent space.  If $q$ has full support, then
\begin{equation}
 \operatorname{span}\{\mu-\nu:\mu,\nu\in\mathcal M_q\}=W_A.
\label{pt:eq:tangent-space}
\end{equation}
Only the reverse inclusion needs proof.  For $0\ne\eta\in W_A$, write its
Jordan decomposition as $\eta=\eta^+-\eta^-$, where both positive parts have
the same mass $c>0$ and, after division by $c$, the same barycentre $b$.
Choose $\varepsilon>0$ small enough that
$\tau=(q-\varepsilon b)/(1-\varepsilon)$ belongs to $\Delta(A)$.  Then
\[
 \mu^\pm:=\frac{\varepsilon}{c}\eta^\pm
          +(1-\varepsilon)\delta_\tau
 \in\mathcal M_q,
 \qquad
 \eta=\frac c\varepsilon(\mu^+-\mu^-),
\]
which proves~\eqref{pt:eq:tangent-space}.  The same argument intrinsically on
a face gives the space $W_B$.

Fix a nondegenerate $r$ supported on $B$.  It is reachable from $q$: choose
\[
 0<\varepsilon<\min\left\{1,
       \min_{b\in B}\frac{q(b)}{r(b)}\right\}
\]
and use a two-record experiment whose first record has conditional probability
$\varepsilon r(a)/q(a)$.  That record has mass $\varepsilon$ and posterior
$r$.

Implement $\nu,\nu'\in\mathcal M_{r_B}$ by channels on $B$, extend their
off-support rows arbitrarily to $A$, and pad their record alphabets if
necessary.  Vary only the continuation on the branch reaching $r$ from
$\nu$ to $\nu'$.  The aggregate posterior law changes by
$\varepsilon i_B(\nu-\nu')$.  Lemma~\ref{pt:lem:finite-branch-extension}\textup{(b)} and
support transport give
\[
 \operatorname{sgn}L_q^Ai_B(\nu-\nu')
 =\operatorname{sgn}L_{r_B}^B(\nu-\nu').
\]
Indeed, strict preference handles the positive case, swapping continuations
handles the negative case, and indifference gives both inequalities in the
zero case.  By~\eqref{pt:eq:tangent-space}, this sign identity extends by
positive scaling to $W_B$.  Positive trace orientation makes
\[
 L_{r_B}^B(\chi_{r_B}-\delta_{r_B})
 =G_{r_B}(\chi_{r_B})>0,
\]
so $L_{r_B}^B\ne0$.  Two nonzero linear functionals with the same positive
half-space have the same kernel and are positive scalar multiples; this gives
the unique coefficient in~\eqref{pt:eq:branch-linear}.  Since both linear
functionals in that equation are fixed before a reaching experiment is
chosen, the coefficient is path-independent.

For the aggregation formula, put $B_i:=\operatorname{supp}(r_i)$ and
\[
 \nu_i:=\pi_{B_i}\mu_{r_i,Q_i}\in\mathcal M_{(r_i)_{B_i}}.
\]
The compound law differs from the first-stage law by
\[
 \mu_{q,P_1\triangleright(Q_1,\ldots,Q_k)}-\mu_{qP_1}
 =\sum_{i:m_i>0}m_i\,i_{B_i}
   (\nu_i-\delta_{(r_i)_{B_i}}).
\]
Apply $L_q^A$, use~\eqref{pt:eq:branch-linear}, and use
$G_{r_i}(\delta_{r_i})=0$.  Singleton summands vanish.  This is
\eqref{pt:eq:branchaggregation}.
\end{proof}

For full-support $q\in\Delta(A)$, an injection
$\iota:B\hookrightarrow A$, and full-support $r\in\Delta(B)$, define
\[
 \beta_\iota(q,r):=\beta_A(q,\iota_\#r),
\]
where $\iota_\#r\in\D(A)$ is the image of $r$ under $\iota$.  The image posterior $\iota_\#r$ is reachable by the two-record construction
in the proof, and its support inclusion is precisely $i_\iota$.  Thus
\[
 L_q^A i_\iota=\beta_\iota(q,r)L_r^B
 \quad\hbox{on }W_B.
\]

\begin{lemma}[Cocycle, face coherence, and the chain gauge]
\label{pt:lem:chain}
There is a relabelling-invariant family of positive scales
\[
\{a_q^A: A\text{ finite},\ |A|\ge2,\ q\in\Delta(A)\text{ full support}\}
\]
for the coefficients of Lemma~\ref{pt:lem:branchagg}, such that
\begin{equation}
 \beta_\iota(q,r)=\frac{a_q^A}{a_r^B}
\tag{FS}\label{pt:eq:facescale}
\end{equation}
for every injection $\iota:B\hookrightarrow A$ with $|A|,|B|\ge2$ and every
full-support $q\in\Delta(A)$ and $r\in\Delta(B)$; here
$\beta_\iota(q,r)$ is the coefficient for reaching the posterior
$\iota_\#r$ on the face $\iota(B)$.  For each prior in this family, put
\[
 V_q:=\frac{G_q}{a_q}.
\]
These representatives satisfy the following exact chain rule.  For every
finite $A$ with $|A|\ge2$, every full-support $q\in\Delta(A)$, every $k\ge1$,
every $P_1:A\to\Delta(\{1,\ldots,k\})$, and every list of finite continuation
channels $Q_1,\ldots,Q_k$ on $A$, use the branch quantities in
\eqref{pt:eq:branch-notation}.  Then
\begin{equation}
 V_q\bigl(\mu_{q,P_1\triangleright(Q_1,\ldots,Q_k)}\bigr)
 =V_q(\mu_{qP_1})
  +\sum_{i:m_i>0}m_iV_{r_i}(\mu_{r_i,Q_i}).
\tag{CR}\label{pt:eq:chainrule}
\end{equation}
Singleton continuation terms are zero.
\end{lemma}

\begin{proof}
There are two normalisation tasks: first solve the branch cocycle within each
fixed action-set dimension, then remove the residual scalar attached to
embedding an $m$-point face in an $n$-point simplex.

We first prove the cocycle before choosing scales.  For nested injections
$C\xhookrightarrow{\jmath}B\xhookrightarrow{\iota}A$ and full-support
priors $q,r,s$ on $A,B,C$, respectively, the defining identities on $W_C$
are
\[
 L_q^A i_{\iota\circ\jmath}
 =\beta_\iota(q,r)L_r^B i_\jmath
 =\beta_\iota(q,r)\beta_\jmath(r,s)L_s^C.
\]
The direct coefficient from $q$ to $s$ is unique because $L_s^C\ne0$ by
Axiom~\textup{(A4)}.  Therefore
\begin{equation}
 \beta_{\iota\circ\jmath}(q,s)
 =\beta_\iota(q,r)\beta_\jmath(r,s).
\label{pt:eq:branch-cocycle}
\end{equation}
This single identity covers full-support paths, full-to-face paths, and
nested support faces.

For a fixed nondegenerate $A$, apply~\eqref{pt:eq:branch-cocycle} to identity
embeddings.  Fix a full-support basepoint $q_A$ and put
\[
 a_q^A:=\beta_{\mathrm{id}_A}(q,q_A).
\]
Then
\[
 \beta_{\mathrm{id}_A}(q,r)=\frac{a_q^A}{a_r^A}.
\tag{WS}\label{pt:eq:within-scale}
\]
Choosing $q_A$ to be uniform and using exact relabelling of $G$ makes these
initial scales invariant under bijections.  They remain free up to one
positive multiplier for each cardinality.

For an injection $\iota:B\hookrightarrow A$, define its residual defect by
\[
 \Gamma(\iota):=\beta_\iota(q,r)\frac{a_r^B}{a_q^A}.
\]
Equation~\eqref{pt:eq:within-scale} shows that this number is independent of
$q$ and $r$: changing either prior merely inserts and cancels the
corresponding within-set coefficient.  Equation~\eqref{pt:eq:branch-cocycle}
gives
\[
 \Gamma(\iota\circ\jmath)=\Gamma(\iota)\Gamma(\jmath).
\]
Exact relabelling implies that $\Gamma(\iota)$ depends only on
$n=|A|$ and $m=|B|$; write it $\Gamma_{n,m}$.  Thus
\[
 \Gamma_{n,\ell}=\Gamma_{n,m}\Gamma_{m,\ell},\qquad \Gamma_{n,n}=1.
\]
Set $t_n:=\Gamma_{n,2}$ and $t_2:=1$.  Then
$\Gamma_{n,m}=t_n/t_m$.  Replacing every $a_q^A$ on an $n$-point set by
$t_n a_q^A$ changes the defect to
$\Gamma_{n,m}t_m/t_n=1$, while preserving~\eqref{pt:eq:within-scale} and
relabel invariance.  This proves~\eqref{pt:eq:facescale}.

For a boundary prior, $a_r$ means the scale on its support; the preceding
construction says precisely that this agrees with the scale obtained by
reaching that support as a face.  Divide~\eqref{pt:eq:branchaggregation} by
$a_q$ and use~\eqref{pt:eq:facescale}.  Every nondegenerate branch coefficient
becomes
\[
 \frac1{a_q}\,\beta_A(q,r_i)G_{r_i}
 =\frac1{a_{r_i}}G_{r_i}=V_{r_i}.
\]
Singleton values vanish, so their scales and coefficients may be fixed
arbitrarily.  This proves~\eqref{pt:eq:chainrule}.
\end{proof}

\paragraph{Compatibility of the product and branch calibrations.}

The product and sequential evaluations of full revelation now meet.  Their
compatibility first removes the provisional product interaction and then
leaves a strongly additive uncertainty functional.

\begin{lemma}[Interaction collapse and universal chain scale]
\label{pt:lem:scalecoherence}
The coefficient $\kappa$ in~\eqref{pt:eq:QA} is zero.  Moreover there is a
single $C>0$ such that $a_q=C$ for every finite prior with support of size at
least two, and the same value may be assigned on singleton supports; for a
boundary prior, $a_q$ denotes the scale on its support face.
Consequently the final representatives $V_q=G_q/C$ satisfy, for all nonempty
finite sets $A,B$, all priors $p\in\Delta(A)$ and $r\in\Delta(B)$, and all
$\mu\in\mathcal M_p$ and $\nu\in\mathcal M_r$,
\begin{equation}
 V_{p\otimes r}(\mu\otimes\nu)
   =V_p(\mu)+V_r(\nu).
\tag{PA}\label{pt:eq:PArecovered}
\end{equation}
The chain rule~\eqref{pt:eq:chainrule} extends from full-support fibres to
all finite priors by exact support transport.
The cross-prior bridge of Lemma~\ref{pt:lem:blockbridge} is valid with $V$
in place of $G$.
\end{lemma}

\begin{proof}
For the rest of this subsection, put
\[
 H_G(q):=G_q(\chi_q),\qquad Z(q):=1+\kappa H_G(q).
\]
For nondegenerate $q$, Axiom~\textup{(A4)} and zero normalisation give
$H_G(q)>0$, while the strict slice-slope
\eqref{pt:eq:QAslope} gives
\begin{equation}
 Z(q)>0.
\label{pt:eq:Zpositive}
\end{equation}
For a singleton, $H_G=0$ and $Z=1$.

\emph{Product revelation versus sequential revelation.}
For nondegenerate full-support $p,r$, product quasi-additivity gives
\begin{equation}
 H_G(p\otimes r)
 =H_G(p)+H_G(r)+\kappa H_G(p)H_G(r).
\tag{QH}\label{pt:eq:QH}
\end{equation}
Compute the same full revelation by first revealing the $A$-coordinate and
then the $B$-coordinate.  The first-stage law is
$\chi_p\otimes\delta_r$ and has value $H_G(p)$ by~\eqref{pt:eq:QA}.
After action $a$ is revealed, the posterior is $e_a\otimes r$ on the face
$\{a\}\times B$; support restriction and relabelling identify the
continuation value with $H_G(r)$.  Every such branch has coefficient
$a_{p\otimes r}/a_r$ by~\eqref{pt:eq:facescale}.  Since the branch weights
sum to one, Lemma~\ref{pt:lem:branchagg} gives
\begin{equation}
 H_G(p\otimes r)=H_G(p)+\frac{a_{p\otimes r}}{a_r}H_G(r).
\tag{SA}\label{pt:eq:seqA}
\end{equation}
Comparing~\eqref{pt:eq:QH} and~\eqref{pt:eq:seqA}, and using
$H_G(r)>0$, yields
\[
 \frac{a_{p\otimes r}}{a_r}=Z(p).
\tag{R1}\label{pt:eq:ratio1}
\]
Revealing $B$ first gives symmetrically
\[
 \frac{a_{p\otimes r}}{a_p}=Z(r).
\tag{R2}\label{pt:eq:ratio2}
\]
Thus $a_rZ(p)=a_pZ(r)$ for every pair of nondegenerate priors, even on
different finite sets.  By~\eqref{pt:eq:Zpositive}, there is a universal
$C>0$ such that
\begin{equation}
 a_q=CZ(q).
\tag{AS}\label{pt:eq:ascaleZ}
\end{equation}
Assign $a=C$ on singleton supports.

\emph{The grouping weight.}
Let $q$ have full support and let
$\mathcal P=\{A_1,\ldots,A_m\}$ partition its action set into nonempty
cells.  Put
\[
 p_i=q(A_i),\qquad q^i=q(\,\cdot\mid A_i),\qquad
 p=(p_1,\ldots,p_m).
\]
The superscript in $q^i$ marks the conditioning cell, not a block tag.
If $C_{\mathcal P}$ reports only the cell, undetectable-distinctions
neutrality gives
\[
 (q,C_{\mathcal P})\sim(p,\Id_{\{1,\ldots,m\}}).
\]
The cross-prior bridge therefore identifies the coarse value as
\begin{equation}
 G_q(\mu_{q,C_{\mathcal P}})=H_G(p).
\tag{BG}\label{pt:eq:blockvalue}
\end{equation}
Append full revelation within each cell.  The final law is $\chi_q$, and the
branch formula together with~\eqref{pt:eq:ascaleZ} gives
\begin{equation}
 H_G(q)=H_G(p)+Z(q)\sum_i p_i\frac{H_G(q^i)}{Z(q^i)}.
\tag{GR}\label{pt:eq:groupingH}
\end{equation}
Here singleton cells contribute zero.

If $\kappa=0$, then $Z\equiv1$ and the conclusion follows.  Otherwise set
$w(q):=1/Z(q)>0$.  Because $\kappa\ne0$ in this case,
$\kappa H_G(x)=Z(x)-1$.  Substitution into
\eqref{pt:eq:groupingH} gives
\[
 Z(q)-1=Z(p)-1+
 Z(q)\sum_i p_i\left(1-\frac1{Z(q^i)}\right).
\]
Cancelling and inverting yields the positive grouping-weight equation
\begin{equation}
 w(q)=w(p)\sum_i p_iw(q^i).
\tag{W}\label{pt:eq:weight}
\end{equation}
All distributions in what follows are read on their positive-probability
supports.  For finite probability vectors $u$ and $v$, partition $u\otimes v$ by the
first coordinate: the coarse
distribution is $u$ and every conditional distribution is $v$.  Thus
\eqref{pt:eq:weight} gives
\begin{equation}
 w(u\otimes v)=w(u)w(v).
\tag{M}\label{pt:eq:wmult}
\end{equation}

\emph{Two groupings force the weight to be constant.}
Take arbitrary finite probability vectors $u,v$ and write
$x=w(u)$ and $y=w(v)$.  Let $\tau$ be the law of a triple
$(\ell,z_1,z_2)$: $\ell$ is uniform on $\{0,1\}$, and conditional on
$\ell=0$ the pair $(z_1,z_2)$ has law $u\otimes u$, while conditional on
$\ell=1$ it has law $v\otimes v$.  In tagged notation,
\[
 \tau:=\tfrac12(u\otimes u)^0\sqcup\tfrac12(v\otimes v)^1.
\]
Grouping first by $\ell$ and using~\eqref{pt:eq:wmult} gives
\begin{equation}
 w(\tau)=w(\tfrac12,\tfrac12)\frac{x^2+y^2}{2}.
\tag{E1}\label{pt:eq:twoeval1}
\end{equation}
Instead group first by $(\ell,z_1)$.  That marginal is
$S=\tfrac12u^0\sqcup\tfrac12v^1$, so
\[
 w(S)=w(\tfrac12,\tfrac12)\frac{x+y}{2}.
\]
The conditional law of $z_2$ is $u$ when $\ell=0$ and $v$ when $\ell=1$.
A second use of~\eqref{pt:eq:weight} therefore gives
\begin{equation}
 w(\tau)=w(S)\frac{x+y}{2}
 =w(\tfrac12,\tfrac12)\left(\frac{x+y}{2}\right)^2.
\tag{E2}\label{pt:eq:twoeval2}
\end{equation}
The common leading factor is positive.  Comparing
\eqref{pt:eq:twoeval1} and~\eqref{pt:eq:twoeval2} yields $(x-y)^2=0$.
Thus $w$ is constant on all finite probability vectors.  Its singleton value
is one, so $w\equiv1$.

It follows that $Z\equiv1$, hence $\kappa H_G(q)=0$ for every $q$.
Axiom~\textup{(A4)} supplies a nondegenerate $q$ with $H_G(q)>0$, so
$\kappa=0$, contradicting the temporary alternative.  Equation
\eqref{pt:eq:ascaleZ} now gives $a_q=C$ universally.
Dividing by $C$ proves \eqref{pt:eq:PArecovered} and
\eqref{pt:eq:chainrule} at full support; exact support transport extends
them to boundary priors.  The common positive division also preserves the
cross-prior bridge.
\end{proof}

\subsubsection{Entropy identification}

\begin{lemma}[Entropy reduction and Faddeev identification]
\label{pt:lem:faddeevsketch}
For every nonempty finite $A$ and every $q\in\Delta(A)$, define the value of
full revelation by
\[
 \mathcal H(q):=V_q(\chi_q),
\]
with boundary priors read on their support.
For every nonempty finite $X$ and every channel $P:A\to\Delta(X)$, put
$m_x:=\sum_aq(a)P(x\mid a)$ and, for each reached record $x$,
$r_x(a):=q(a)P(x\mid a)/m_x$.  Then
\begin{equation}
 V_q(\mu_{qP})
 =\mathcal H(q)-\sum_{x:m_x>0}m_x\mathcal H(r_x).
\tag{ER}\label{pt:eq:entropy-reduction}
\end{equation}
Moreover, there is a single $\alpha>0$, independent of all finite alphabets
and priors, such that
\[
 \mathcal H(q)=\alpha\Sh(q)
\]
and consequently
\[
 V_q(\mu_{qP})=\alpha I_{qP}(A;X).
\tag{MI}\label{pt:eq:value-MI}
\]
\end{lemma}

\begin{proof}
\emph{Entropy reduction.}
Assume first that $q$ has full support.  Append full revelation after $P$.
The final posterior law is $\chi_q$, while the continuation after record $x$
has full-revelation law $\chi_{r_x}$ on the support of $r_x$.  The exact
chain rule~\eqref{pt:eq:chainrule} therefore gives
\[
 \mathcal H(q)=V_q(\mu_{qP})
 +\sum_{x:m_x>0}m_x\mathcal H(r_x),
\]
which is~\eqref{pt:eq:entropy-reduction}.  A boundary prior is reduced to its
support; null action rows and null records affect neither side.

\emph{Strong additivity.}
Let $\mathcal P=\{A_1,\ldots,A_m\}$ be a partition of the support of $q$,
and put
\[
 p_i=q(A_i),\qquad q^i=q(\,\cdot\mid A_i),\qquad
 p=(p_1,\ldots,p_m).
\]
In~\eqref{pt:eq:groupingH}, Lemma~\ref{pt:lem:scalecoherence} gives
$\kappa=0$ and $V_q=G_q/C$, hence $Z\equiv1$ and $H_G=C\mathcal H$.  Dividing that identity by $C$ yields
\begin{equation}
 \mathcal H(q)=\mathcal H(p)+\sum_i p_i\mathcal H(q^i).
\tag{FA}\label{pt:eq:faddeev-recursion}
\end{equation}
Relabelling invariance makes $\mathcal H$ symmetric, support restriction makes
it expansive, and $\mathcal H$ vanishes on point masses.  Record processing
gives $\mathcal H\ge0$, while Axiom~\textup{(A4)} gives strict positivity at
nondegenerate full-support priors.  Equation~\eqref{pt:eq:faddeev-recursion}
also allows zero outer weights by deleting and reinserting null cells.

\emph{Continuity: the binary case.}
Write $h(t):=\mathcal H(t,1-t)$.  We derive, rather than assume, its
continuity from Axiom~\textup{(A2)}.  Fix $c\ge0$.  Choose a full-support
binary prior $\theta$ with $\mathcal H(\theta)>0$, as supplied by
Axiom~\textup{(A4)}.  Choose $n\ge1$ and $\zeta\in[0,1]$ so that
\[
 c=\zeta n\mathcal H(\theta).
\]
For $\theta_n:=\theta^{\otimes n}$, product additivity
\eqref{pt:eq:PArecovered} gives
$\mathcal H(\theta_n)=n\mathcal H(\theta)$.  Let $E_{n,\zeta}$ reveal the
$n$-bit action fully with probability $\zeta$ and report a common erasure
symbol otherwise.  Its posterior law is
$\zeta\chi_{\theta_n}+(1-\zeta)\delta_{\theta_n}$, so affinity gives
\[
 V_{\theta_n}(\mu_{\theta_n,E_{n,\zeta}})=c.
\]

Now fix the block environment
\[
 K_c:=\Id_{\{0,1\}}\sqcup E_{n,\zeta}.
\]
Let $q_t^0$ put the binary prior $(t,1-t)$ on its first block and let
$\theta_n^1$ put $\theta_n$ on its second block.  The cross-prior bridge says
\begin{align*}
 h(t)\ge c&\quad\Longleftrightarrow\quad q_t^0\succeq_{K_c}\theta_n^1,\\
 h(t)\le c&\quad\Longleftrightarrow\quad \theta_n^1\succeq_{K_c}q_t^0.
\end{align*}
The environment is fixed and $t\mapsto q_t^0$ is continuous, so
Axiom~\textup{(A2)} makes both level sets closed for every $c\ge0$.  For a
threshold $c<0$, the superlevel set is all of $[0,1]$ and the sublevel set is
empty because $h\ge0$.
Thus every upper and lower level set of $h$ is closed, and $h$ is continuous.

\emph{Continuity on every finite simplex.}
Proceed by induction on the number of coordinates.  For
$q=(q_1,(1-q_1)\bar q)$ with $q_1<1$, recursion
\eqref{pt:eq:faddeev-recursion} for the partition
$\{\{1\},\{2,\ldots,n\}\}$ gives
\begin{equation}
 \mathcal H(q)=h(q_1)+(1-q_1)\mathcal H(\bar q).
\label{pt:eq:continuity-induction}
\end{equation}
The right side is continuous away from $q_1=1$ by binary continuity and the
induction hypothesis.  At $q_1=0$, support restriction and $h(0)=0$ give
the same formula on the face.  At $q_1=1$, the induction hypothesis makes
$\mathcal H$ bounded on the compact $(n-1)$-simplex, so the second term in
\eqref{pt:eq:continuity-induction} tends to zero and the first tends to
$h(1)=0$.  Relabelling handles the other faces.  Hence $\mathcal H$ is
continuous on each closed finite simplex.

We have now verified all hypotheses of the strong-additivity form of
Faddeev's theorem.  Axiom~\textup{(A2)} supplied continuity;
Lemma~\ref{pt:lem:transport} supplied relabelling invariance and
expansibility; $\mathcal H$ vanishes on point masses; and
Lemmas~\ref{pt:lem:chain} and~\ref{pt:lem:scalecoherence} supplied strong
additivity~\eqref{pt:eq:faddeev-recursion}.  Record processing has already
established $\mathcal H\ge0$.  The finite entropy characterisation
\cite[Theorem~6]{BaezFritzLeinster2011} therefore gives
\[
 \mathcal H(q)=\alpha\Sh(q)
\]
for some $\alpha\ge0$.  Axiom~\textup{(A4)} forces $\alpha>0$.
Substitution in~\eqref{pt:eq:entropy-reduction} gives
\[
 V_q(\mu_{qP})
 =\alpha\left(\Sh(q)-\sum_x m_x\Sh(r_x)\right)
 =\alpha I_{qP}(A;X),
\]
where the sum is over positive-mass records.  This proves
\eqref{pt:eq:value-MI}.
\end{proof}

\begin{proof}[Proof of Proposition~\ref{pt:prop:main}]
The cross-prior bridge and Lemma~\ref{pt:lem:faddeevsketch} give, for arbitrary
pairs $(q,P)$ and $(r,Q)$,
\[
 q^0\succeq_{P\sqcup Q}r^1
 \quad\Longleftrightarrow\quad
 I_{qP}\ge I_{rQ},
\tag{BMI}\label{pt:eq:final-block-MI}
\]
including boundary priors by support restriction.  With $Q=P$, the
duplication clause of Axiom~\textup{(A5)} gives the proposition's fixed-channel
claim.  The positive multiplier $\alpha$ cancels from these comparisons; its
cardinal calibration against payoff utility is the next step.  Thus the
conclusion here is ordinal: it fixes the common ranking of pure-trace pairs,
not their exchange rate against material utility.  For distinct
$i,j\in J$, the irrelevant-block clause of Axiom~\textup{(A5)} reduces an
arbitrary block union to the canonical two-block comparison:
\[
 q_i^i\succeq_{\bigsqcup_{k\in J}P_k}q_j^j
 \quad\Longleftrightarrow\quad
 q_i^0\succeq_{P_i\sqcup P_j}q_j^1.
\]
On the information side, the block tag is constant under a block-supported
prior, so
\[
 I_{q_i^i,\,\bigsqcup_{k\in J}P_k}=I_{q_iP_i},
\]
and similarly for block $j$.  Applying~\eqref{pt:eq:final-block-MI} to the
right-hand side proves the block claim.
\end{proof}

\subsection{The two common scales}

For $\ell\in\D(O)$, write $\ell$ also for the channel that pays $\ell$
independently of the action and writes a constant record; the
\emph{payoff-lottery order} is the order induced on $\D(O)$ by this
identification.

\begin{lemma}[Material utility and a common affine scale]\label{lf:lem:material-scale}
Under Axioms \textup{(A1)--(A3)} and
\textup{(A5)--(A8)}, there are outcomes $o^+,o^-$ and a
nonconstant function $u:O\to\R$ such that
\[
u(o^+)=1,
\qquad
u(o^-)=0.
\]
For every payoff lottery $\ell\in\D(O)$, the payoff-lottery order is represented
by $\sum_o\ell(o)u(o)$.  For every full-support prior $q$, the representatives
in Lemma~\ref{lf:lem:terminal-affine} may be normalised so that
\begin{equation}
W_q(\ell)=\sum_o\ell(o)u(o).
\label{lf:eq:payoff-normalisation}
\end{equation}
With this normalisation, for every full-support $q\in\D(A)$, every nonempty
finite $B$, every full-support $s\in\D(B)$, and every finite experiment $E$ at
$q$,
\[
W_{q\otimes s}(E^\uparrow)=W_q(E),
\]
where $E^\uparrow$ ignores the dummy coordinate.  Moreover, for every pair of nonempty
finite action sets, every two full-support priors $q,q'$ on those sets, and
every two finite experiments $E,G$ on the respective sets,
\begin{equation}
(q,E)\succeq(q',G)
\quad\Longleftrightarrow\quad
W_q(E)\ge W_{q'}(G).
\label{lf:eq:common-scale}
\end{equation}
\end{lemma}

\begin{proof}
On a singleton action set, terminal-law sufficiency identifies an experiment
with its payoff lottery.  Public-mixture independence and continuity, proved as
in Lemma~\ref{lf:lem:terminal-affine}, give an affine utility on $\D(O)$.
By Lemma~\ref{lem:relevance-bridge}, Axiom~\textup{(A3)} supplies
$o^+\succ o^-$.  Normalise their values to one
and zero and let $u(o)$ be the value of the degenerate lottery at $o$.
Affinity on the finite simplex gives expected utility.

An action processor may ignore its input and draw an arbitrary target prior.
Axiom~\textup{(A7)}, applied in both directions, therefore makes this payoff
lottery order independent of the prior and action alphabet.  Normalise each
$W_q$ at the same two deterministic payoffs.  The payoff-lottery embedding is
affine and order-reflecting, so positive-affine uniqueness gives
\eqref{lf:eq:payoff-normalisation}.

Dummy lifting between full-support fibres is likewise an affine order
embedding.  Positive-affine uniqueness initially permits
\[
W_{q\otimes s}(E^\uparrow)=cW_q(E)+d,
\qquad c>0.
\]
The two payoff anchors have values zero and one on both sides; hence $c=1$ and
$d=0$.  To compare $E$ at $q$ with $G$ at $q'$, lift each to the common product
prior $q\otimes q'$, using the other prior as the dummy.  Equation
\eqref{lf:eq:dummy-ordinal}, the replacement part of
Lemma~\ref{lf:lem:bookkeeping}, and the common product representative give
\eqref{lf:eq:common-scale}.
\end{proof}

Retain the outcome $o_\ast$ fixed in the preceding pure-trace subsection to
witness Axiom~\textup{(A4)}.  Proposition~\ref{pt:prop:main} has established
that mutual information represents every within-channel and block-supported
comparison on this constant-payoff face.
The normalisation $u(o^+)=1$ and $u(o^-)=0$ fixes only the material origin
and unit.  The trace anchor $o_\ast$ need not be either material anchor, and
its intercept $u(o_\ast)$ is retained.  It remains to identify the cardinal
coefficient of mutual information relative to the material unit and to show
that this coefficient is common to all priors and finite alphabets.  Call a
prior \emph{nontrivial} if its entropy is positive.

\begin{lemma}[The trace scale]\label{lf:lem:trace-scale}
Under Axioms \textup{(A1)--(A8)}, there is one number $\lambda>0$, independent
of the prior and all finite alphabets, such that, for every nonempty finite
$A$ with $\lvert A\rvert\ge2$, every full-support $q\in\D(A)$, every nonempty
finite record set $R$, and every experiment
$E:A\to\D(O\times R)$ that pays $o_\ast$ surely,
\begin{equation}
W_q(E)=u(o_\ast)+\lambda I_q(E).
\label{lf:eq:sure-anchor}
\end{equation}
\end{lemma}

\begin{proof}
Fix a nontrivial full-support $q$.  On the constant-$o_\ast$ face, $W_q$ is
affine by Lemma~\ref{lf:lem:terminal-affine}.  For any experiment $E$ on the constant-$o_\ast$ face, the
finite entropy identity
\begin{equation}
I_q(E)=\Sh(q)-\sum_{(o,\pi)\in\supp\eta_E}
\eta_E(o,\pi)\Sh(\pi).
\label{lf:eq:posterior-entropy}
\end{equation}
shows that mutual information depends only on the marked terminal law and is
affine under public mixtures.  Proposition~\ref{pt:prop:main} says that $W_q$ and
$I_q$ represent the same nonconstant order on the constant-$o_\ast$ face.
Positive-affine uniqueness therefore gives
\[
W_q(E)=\lambda_q I_q(E)+b_q,
\qquad \lambda_q>0.
\]
The uninformative experiment has mutual information zero and, by
\eqref{lf:eq:payoff-normalisation}, value $u(o_\ast)$.  Thus $b_q=u(o_\ast)$.

The coefficient does not depend on $q$.  Let $q$ and $q'$ be nontrivial
full-support priors.  Lift full revelation at $q$ to $q\otimes q'$ by adding an
independent dummy coordinate.  Normalised value and mutual information are both
unchanged.  Since the common information value is $\Sh(q)>0$, cancellation
gives $\lambda_{q\otimes q'}=\lambda_q$.  Lifting full revelation from the
other coordinate gives $\lambda_{q\otimes q'}=\lambda_{q'}$.  Hence all such
coefficients agree.  Equivalently, fix the uniform prior on a two-point dummy
alphabet and call its coefficient $\lambda$.  Axiom~\textup{(A4)} fixes the
strict orientation, so $\lambda>0$.

Finally, any experiment that pays $o_\ast$ surely is a pure record experiment
after the constant payoff coordinate is deleted; restoring that coordinate
changes neither its marked terminal law nor its mutual information.  This
proves \eqref{lf:eq:sure-anchor}.
\end{proof}

\subsection{Exact branch weights}

Let $A$ and $Y$ be arbitrary nonempty finite sets, let $q\in\D(A)$ have full
support with $\lvert A\rvert\ge2$, and let
$P:A\to\D(Y)$ be a first-stage record channel.  For every $y\in Y$, put
\[
m_y:=\sum_{a\in A}q(a)P(y\mid a).
\]
Whenever $m_y>0$, define the posterior on $A$ and its support by
\[
\widehat r_y(a):=\frac{q(a)P(y\mid a)}{m_y},
\qquad
S_y:=\supp\widehat r_y,
\]
and write $r_y\in\D(S_y)$ for the full-support restriction of
$\widehat r_y$ to $S_y$.

We need notation for changing a continuation without fixing what happens on
the other branches.  For a finite continuation experiment
$E:S_y\to\D(O\times R_E)$, choose a retraction
$\rho_y:A\to S_y$ whose restriction to $S_y$ is the identity and set
$E^{\rho_y}(\cdot\mid a):=E(\cdot\mid\rho_y(a))$.  Given an arbitrary
background $\mathbf H=(H_z)_{z\ne y}$ of continuation channels on $A$, let
\[
K_z^E:=
\begin{cases}
E^{\rho_y},&z=y,\\
H_z,&z\ne y,
\end{cases}
\qquad
J_{y,\mathbf H}(E):=P\triangleright(K_z^E)_{z\in Y}.
\]
Here and below, fix any ordering of $Y$ when writing the compound.
Before taking the compound, all continuation record alphabets are embedded as
tagged copies in one finite disjoint union.  Payoff-preserving record
isomorphisms make the choice of tags immaterial.  Also,
$a\notin S_y$ implies $P(y\mid a)=0$, because $q$ has full support; hence the
rows of $E^{\rho_y}$ outside $S_y$ are never used on branch $y$.
The compound channel itself is therefore independent of the chosen
retraction.

\begin{lemma}[Exact reached-branch weight]\label{lf:lem:branch-weight}
Under Axioms \textup{(A1)--(A8)}, for every nonempty finite $A,Y$, every
full-support $q\in\D(A)$ with $\lvert A\rvert\ge2$, every
$P:A\to\D(Y)$, every reached $y\in Y$, every two finite continuation
experiments $E,G$ on $S_y$, and every finite background
$\mathbf H=(H_z)_{z\ne y}$,
\begin{equation}
W_q\bigl(J_{y,\mathbf H}(E)\bigr)
-W_q\bigl(J_{y,\mathbf H}(G)\bigr)
=m_y\bigl[W_{r_y}(E)-W_{r_y}(G)\bigr].
\label{lf:eq:branch-weight}
\end{equation}
\end{lemma}

\begin{proof}
At the branch posterior $\widehat r_y$, support restriction identifies
$(\widehat r_y,E^{\rho_y})$ with $(r_y,E)$, and likewise for $G$.
Lemma~\ref{pt:lem:finite-branch-extension}\textup{(b)} therefore gives, for every $E$ and $G$,
\[
E\succeq_{r_y}G
\quad\Longleftrightarrow\quad
J_{y,\mathbf H}(E)\succeq_qJ_{y,\mathbf H}(G),
\]
with the same equivalence for strict preference.  This statement includes
continuations with different record alphabets: the tagged disjoint union puts
them in the common comparison environment required by the lemma.

To see the affine structure explicitly, let
$\iota_y:\D(S_y)\to\D(A)$ pad a distribution by zeros outside $S_y$.
There is a fixed subprobability measure $\beta_{y,\mathbf H}$, of mass
$1-m_y$, containing the terminal atoms from all branches other than $y$, such
that
\[
\eta_{J_{y,\mathbf H}(E)}
=\beta_{y,\mathbf H}
 +m_y(\mathrm{id}_O,\iota_y)_\#\eta_E.
\]
Thus insertion is affine on marked terminal laws.  If $S_y$ is nontrivial,
the local order is nonconstant: at payoff $o_\ast$, full revelation and
silence have, by Lemma~\ref{lf:lem:trace-scale}, value difference
$\lambda\Sh(r_y)>0$.  Affine
uniqueness consequently gives, for every fixed background $\mathbf H$,
constants $c_{y,\mathbf H}>0$ and $d_{y,\mathbf H}\in\R$ such that
\begin{equation}
W_q\bigl(J_{y,\mathbf H}(E)\bigr)
=c_{y,\mathbf H}W_{r_y}(E)+d_{y,\mathbf H}.
\label{lf:eq:insertion-affine}
\end{equation}

It remains to identify the coefficient and to show that the background cannot
alter it.  Let $\mathbf H^\ast$ pay $o_\ast$ for sure and reveal only a
singleton record after every $z\ne y$.  For every continuation $E$, the
mutual-information chain rule gives
\begin{equation}
I_q\bigl(J_{y,\mathbf H^\ast}(E)\bigr)
=I_q(A;Y)+m_y I_{r_y}(E).
\label{lf:eq:insertion-chain}
\end{equation}
Take for $E$ full revelation, $K^{\mathrm{id}}_{o_\ast}$, and for $G$
silence, $K^0_{o_\ast}$, on $S_y$.  By
Lemma~\ref{lf:lem:trace-scale}, their local value difference is
$\lambda\Sh(r_y)$, while \eqref{lf:eq:insertion-chain} makes their aggregate
value difference $\lambda m_y\Sh(r_y)$.  Comparing these differences in
\eqref{lf:eq:insertion-affine}, and using
$\lambda\Sh(r_y)>0$, yields
$c_{y,\mathbf H^\ast}=m_y$.

Now fix arbitrary $E,G$, and an arbitrary background $\mathbf H$.  The crossed
public half-mixtures
\[
\tfrac12J_{y,\mathbf H}(E)
\oplus\tfrac12J_{y,\mathbf H^\ast}(G)
\quad\text{and}\quad
\tfrac12J_{y,\mathbf H}(G)
\oplus\tfrac12J_{y,\mathbf H^\ast}(E)
\]
have the same marked terminal law.  Indeed, each contains one half of each
background contribution and, on branch $y$, one half of each of
$m_y\eta_E$ and $m_y\eta_G$ (with posterior supports padded by $\iota_y$).
Lemma~\ref{lf:lem:terminal-affine}\textup{(a)} makes the mixtures indifferent.
Affinity of $W_q$ and cancellation therefore give
\[
W_q(J_{y,\mathbf H}(E))-W_q(J_{y,\mathbf H}(G))
=W_q(J_{y,\mathbf H^\ast}(E))-W_q(J_{y,\mathbf H^\ast}(G)).
\]
Applying \eqref{lf:eq:insertion-affine} to the right-hand side with
$c_{y,\mathbf H^\ast}=m_y$ proves \eqref{lf:eq:branch-weight} whenever
$\lvert S_y\rvert\ge2$.

If $S_y$ is a singleton, let $s$ be a full-support prior on a two-point set
$B$.  Adjoin $B$ independently to the global prior, lift the first-stage
channel by $P^\uparrow(y\mid a,b)=P(y\mid a)$, and lift every local and
background continuation so that it ignores $b$.  The reached posterior is
$r_y\otimes s$, its support is nontrivial, and its branch mass remains $m_y$.
The preceding argument applies to this common dummy extension.
Lemma~\ref{lf:lem:material-scale} says that adjoining or deleting the
independent dummy preserves both aggregate values and both local values.
Removing it gives \eqref{lf:eq:branch-weight} for the singleton case.
\end{proof}

For a deterministic payoff profile $g:Y\to O$, let $C(P,g)$ denote the
compound that first runs $P$ and then pays $g(y)$ with no further record.
For $o\in O$, write $C(P,o)$ for $C(P,g)$ with $g\equiv o$.
Write $g_y^o$ for the profile that pays $o$ at $y$ and $o_\ast$ elsewhere.

\begin{lemma}[Deterministic branch assembly]
\label{lf:lem:deterministic-assembly}\label{lf:lem:branch-increment}
Under Axioms \textup{(A1)--(A8)}, for every nonempty finite $A,Y$, every
full-support $q\in\D(A)$ with $\lvert A\rvert\ge2$, every
$P:A\to\D(Y)$, and every deterministic payoff profile $g:Y\to O$,
\begin{equation}
W_q(C(P,g))
=\lambda I_q(A;Y)+\sum_{y\in Y}m_yu(g(y)).
\label{lf:eq:deterministic-assembly}
\end{equation}
\end{lemma}

\begin{proof}
Fix $y\in Y$ and $o\in O$.
If $m_y=0$, changing the payoff after $y$ leaves the marked terminal law
unchanged.  If $m_y>0$, apply Lemma~\ref{lf:lem:branch-weight} with
deterministic local payoffs $o$ and $o_\ast$ and with every other branch paying
$o_\ast$.  Equation~\eqref{lf:eq:payoff-normalisation} gives the local value
difference, and hence
\begin{equation}
W_q\bigl(C(P,g_y^o)\bigr)
-W_q\bigl(C(P,o_\ast)\bigr)
=m_y\bigl[u(o)-u(o_\ast)\bigr].
\label{lf:eq:branch-increment}
\end{equation}
The arbitrary-background clause of Lemma~\ref{lf:lem:branch-weight} gives the
same increment when the payoffs on the other branches have already been
changed.  Enumerate the reached branches and change their payoffs from
$o_\ast$ to $g(y)$ one at a time; unreached branches have no effect.  The
resulting telescope is
\begin{equation}
W_q(C(P,g))-W_q(C(P,o_\ast))
=\sum_{y\in Y}m_y\bigl[u(g(y))-u(o_\ast)\bigr].
\label{lf:eq:telescope}
\end{equation}
The baseline compound pays $o_\ast$ surely and its complete record is $Y$.
Lemma~\ref{lf:lem:trace-scale} therefore gives
$W_q(C(P,o_\ast))=u(o_\ast)+\lambda I_q(A;Y)$.
Substitution in \eqref{lf:eq:telescope}, followed by
$\sum_{y\in Y}m_y=1$, proves the formula.
\end{proof}

\subsection{Proof of the theorem}

\begin{proof}[Proof of Theorem~\ref{thm:main}]
Assume first Axioms \textup{(A1)--(A8)}.  Let $q$ have full support on a
nontrivial action set, and let $K:A\to\D(O\times R)$ be arbitrary.  Put
$Y=O\times R$ and take as first stage
\[
P_K(y\mid a):=K(y\mid a),
\qquad y=(o,r).
\]
After $y=(o,r)$, pay $o$ for sure and reveal nothing further.  Call the
resulting compound $K^{\mathrm{seq}}$.  Conditional on its final record
$((o,r),\ast)$, the posterior over actions is exactly the posterior conditional
on $(o,r)$ under $K$, and both records carry payoff $o$.  Thus $K$ and
$K^{\mathrm{seq}}$ have the same marked terminal law.  By
Lemma~\ref{lf:lem:terminal-affine}\textup{(a)},
$(q,K)\sim(q,K^{\mathrm{seq}})$.

Apply Lemma~\ref{lf:lem:deterministic-assembly} with $g(o,r)=o$.  The two
identities needed are
\[
I_q(A;Y)=I_{qK}(A;O,R)
\]
and
\[
\sum_{o,r}p_{qK}(o,r)u(o)=\E_{qK}[u(O)],
\]
respectively.  Together with sequentialisation, they give
\begin{equation}
W_q(K)=\E_{qK}[u(O)]+\lambda I_{qK}(A;O,R).
\label{lf:eq:full-support-value}
\end{equation}

If $A$ is a singleton, mutual information is zero and the marked terminal law
is that of the induced payoff lottery.  Equation
\eqref{lf:eq:payoff-normalisation} gives
\eqref{lf:eq:full-support-value} directly.  If $q$ is on the boundary, put
$S=\supp q$ and write $(q_S,K_S)$ for the restriction.  Then
\begin{equation}
(q,K)\sim(q_S,K_S),
\qquad
V(q,K)=V(q_S,K_S).
\label{lf:eq:support-value}
\end{equation}
The first identity is Lemma~\ref{lf:lem:bookkeeping}\textup{(b)}; the second
follows directly from the finite sums defining expected utility and mutual
information.  Thus no representative $W_q$ is assigned to a boundary fibre.

For arbitrary finite experiments $(q,K)$ and $(p,L)$, restrict both priors to
their supports, apply
\eqref{lf:eq:common-scale} and \eqref{lf:eq:full-support-value} there, and use
\eqref{lf:eq:support-value} to transport the comparison back.  This gives
\begin{equation}
(q,K)\succeq(p,L)
\quad\Longleftrightarrow\quad
V(q,K)\ge V(p,L),
\label{lf:eq:pair-value}
\end{equation}
with the same $u$ and $\lambda$ for every pair of finite alphabets.  The
duplication clause of Axiom~\textup{(A5)} turns this pair formula into the
within-channel representation.  Its irrelevant-block clause applies the same
formula to any two block-supported priors in any finite block environment.
This is exactly the theorem's same-scale moreover clause.  The construction
made $u$ nonconstant and $\lambda>0$, proving the forward implication.

Conversely, suppose that a nonconstant $u:O\to\R$ and $\lambda>0$ represent
every finite environment by $V$.  Completeness and transitivity of the order
on $\R$ give Axiom~\textup{(A1)}.  On every fixed finite
$(A,O,R)$, expected utility and mutual information are jointly continuous in
$(K,q)$; for mutual information, use its finite entropy formula with
$0\log0=0$.  Therefore
\[
\{(K,q,p):V(q,K)\ge V(p,K)\}
\]
is closed, which is Axiom~\textup{(A2)}.

For a duplicated channel, a prior supported on either tagged copy has the same
payoff law and information as its untagged prior.  This proves both directions
of the duplication clause in Axiom~\textup{(A5)}.  More generally, under a
prior supported on block $i$ of $\bigsqcup_{k\in I}K_k$, the block label is
constant and
\[
V\left(q_i^i,\bigsqcup_{k\in I}K_k\right)=V(q_i,K_i).
\]
The same identity for block $j$ proves both directions of the
irrelevant-block clause, for every finite $I$ and every ordered pair
$i\ne j$.  Thus Axiom~\textup{(A5)} holds.

A payoff-preserving record processor leaves the payoff marginal unchanged and,
by the Markov chain
$A-(O,R)-(O,R')$, weakly lowers mutual information.  This is
Axiom~\textup{(A6)}.  For every action processor $S:A\to\D(B)$ and every
Bayesian pushforward in \eqref{eq:action-processing}, the payoff--record
marginal is again unchanged and the undivided joint law forms the Markov chain
\[
(O,R)-A-B.
\]
Hence $I(B;O,R)\le I(A;O,R)$, proving Axiom~\textup{(A7)}.  A completion on a
processed action with $(qS)(b)=0$ has zero joint mass, so this argument covers
every completion allowed by the axiom.

For Axiom~\textup{(A8)}, fix every binary first-stage channel and continuation
triple covered by that axiom, and put
\[
\alpha:=\sum_{a\in A}q(a)P(1\mid a)>0.
\]
The expected-utility and mutual-information chain rules give the exact
difference
\begin{align*}
&V\bigl(q,P\triangleright(K,M)\bigr)
 -V\bigl(q,P\triangleright(L,M)\bigr)\\
&\qquad=\alpha\bigl[V(q_1,K)-V(q_1,L)\bigr].
\end{align*}
The first-stage information term and the common continuation after record $2$
cancel.  Since $\alpha>0$, the aggregate difference is nonnegative exactly
when the conditional difference is nonnegative, and it is positive exactly
when the conditional difference is positive.  This proves the weak
biconditional in Axiom~\textup{(A8)} and the strict biconditional.

Finally, nonconstancy of $u$ supplies distinct $o^+,o^-\in O$ with
$u(o^+)>u(o^-)$.  In the two-action channel specified by
Axiom~\textup{(A3)}, the two pure lotteries have those values and zero mutual
information, proving material relevance.  For any $o\in O$, the two lotteries
in the four-action channel specified by Axiom~\textup{(A4)} both pay $o$
surely, while their information values are $\log2$ and zero.  Because
$\lambda>0$, the required strict comparison holds (in fact at every $o$).
This verifies Axioms~\textup{(A3)} and~\textup{(A4)} and completes the
converse.  The displayed block identity also proves the theorem's moreover
clause with the same $u$ and $\lambda$.
\end{proof}
\clearpage
\section{Further results and choice implications}\label{app:auxiliary-results}

This appendix follows the order of Sections~\ref{sec:theorem}
and~\ref{sec:implications}.  It first proves the representation's immediate
consequences and the independence result, then turns to the trace test,
deliberate randomisation, record value, and measurement.  None of these
results is used to prove Theorem~\ref{thm:main}; Appendix~\ref{app:main-proof}
contains the complete representation argument.

For $q\in\D(A)$ and $K:A\to\D(O\times R)$, write
$E_{qK}:=\E_{qK}[u(O)]$ and, whenever $p_{qK}(o,r)>0$,
\[
 \pi^{qK}_{or}(a):=\frac{q(a)K(o,r\mid a)}{p_{qK}(o,r)},
 \qquad
 \mu_{qK}:=\sum_{o,r}p_{qK}(o,r)\,\delta_{\pi^{qK}_{or}},
\]
where zero-probability terms are omitted.

\subsection{Consequences of the representation}
\label{app:representation-consequences}

\begin{corollary}[No payoff--trace complementarity]\label{cor:matching}
Suppose Axioms \textup{(A1)--(A8)} hold.  Fix $A$ and $q\in\D(A)$.  If channels $K$ and $L$ satisfy
$E_{qK}=E_{qL}$ and $\mu_{qK}=\mu_{qL}$, then $(q,K)\sim(q,L)$.
\end{corollary}

\begin{proof}
The result follows from Theorem~\ref{thm:main} and
$\I_{qK}(A;O,R)=\Sh(q)-\int\Sh(\pi)d\mu_{qK}(\pi)$.
\end{proof}

\begin{proof}[Proof of Proposition~\ref{prop:uniqueness}]
\emph{Range.}
Choose $o_*\in O$ with $u(o_*)=\underline u$.  The range of $V$ is
$[\underline u,\infty)$.  The lower bound is immediate.  Let a constant-$o_*$
channel reveal the realised action with probability $t$ and otherwise report $*$:
\[
 G_t(o_*,a'\mid a)=t\1_{a'=a},
 \qquad
 G_t(o_*,*\mid a)=1-t.
\]
At the uniform lottery on $n$ actions,
$V(q,G_t)=\underline u+\lambda t\log n$ fills
$[\underline u,\underline u+\lambda\log n]$; these intervals are unbounded.

\emph{Ordinal uniqueness.}
Let $W$ be common-scale and define
$\varphi(v):=W(q,K)$ for any pair with $V(q,K)=v$.  This is well defined: equality of the two $V$-values gives
indifference in their two-block environment by Theorem~\ref{thm:main}, and the common-scale property gives equality
of their $W$-values; strict inequalities are preserved in the same way.  Thus
$W=\varphi\circ V$ for a strictly increasing $\varphi$.  The converse follows from
Theorem~\ref{thm:main}, proving part~\textup{(i)}.

\emph{Affine uniqueness.}
The two chain rules give, for every
$P:A\to\D(\{1,\ldots,k\})$ and every
$K_1,\ldots,K_k,L_1,\ldots,L_k:A\to\D(O\times R)$, writing
$m_i:=\sum_aq(a)P(i\mid a)$,
\[
 V\bigl(q,P\triangleright(K_1,\ldots,K_k)\bigr)
 -V\bigl(q,P\triangleright(L_1,\ldots,L_k)\bigr)
 =\sum_{i:m_i>0}m_i
 \bigl[V(q_i,K_i)-V(q_i,L_i)\bigr].
\]
Hence every $\alpha V+\beta$, $\alpha>0$, is common-scale and branch-additive.

Conversely, write a branch-additive common-scale representation as
$W=\varphi\circ V$ and put
\[
 \psi(x):=\varphi(\underline u+x)-\varphi(\underline u),
 \qquad x\ge0.
\]
For $x\in[0,\lambda\log n]$, use the range construction after the first
branch of a public coin with probabilities $\theta$ and $1-\theta$, and the
uninformative constant-$o_*$ channel elsewhere.  Comparing this compound with
the uninformative compound in \eqref{eq:branch-additive-index} gives
\[
 \psi(\theta x)=\theta\psi(x)
 \qquad(0\le\theta\le1).
\]
Taking $x=\lambda\log n$ shows that $\psi$ is linear on $[0,\lambda\log n]$.  These intervals are nested and
unbounded, so $\psi(x)=\alpha x$ on $[0,\infty)$ for one $\alpha>0$.  Therefore
$\varphi(v)=\alpha v+\beta$, proving part~\textup{(ii)}.

\emph{Parameters.}
Any second trace-tempered index $\widetilde V$ is common-scale by block
invariance and branch-additive by the same chain rules.  Part~\textup{(ii)},
restricted first to deterministic payoffs and then to constant-payoff record channels, gives
$\widetilde u=\alpha u+\beta$ and $\widetilde\lambda=\alpha\lambda$.
\end{proof}

\begin{proof}[Proof of Corollary~\ref{cor:signduality}]
Part \textup{(i)} is Theorem~\ref{thm:main}.  For part \textup{(ii)}, reverse
the primitive relation in every environment.  The reversed family satisfies
Axioms~\textup{(A1)}, \textup{(A2)}, \textup{(A5)}, and \textup{(A8)}; Axiom~\textup{(A3)} holds with
its two witnessing outcomes interchanged.  It also satisfies the forward
versions of Axioms~\textup{(A4)}, \textup{(A6)}, and
\textup{(A7)}.  Theorem~\ref{thm:main} therefore represents the reversed
family by
\(
  \E[\widehat u(O)]+\widehat\lambda\I(A;O,R)
\)
with \(\widehat\lambda>0\).  Negating that index represents the original
family, so it has the asserted form with
\(u=-\widehat u\) and \(\lambda=-\widehat\lambda<0\).

For part \textup{(iii)}, collapse the action alphabet of any pair $(q,K)$ to
a singleton.  Axiom~$(\mathrm{A}7^0)$ makes the pair indifferent to the
resulting one-action channel, whose output law is $p_{qK}$.  Axiom
$(\mathrm{A}6^0)$ then erases its record without changing value.  Thus every
pair is indifferent to the one-action channel that delivers precisely its
payoff marginal.  Axiom~\textup{(A5)} makes these reductions coherent across
blocks.  On payoff lotteries, Axiom~\textup{(A8)} applied to an
action-independent public coin gives mixture independence, Axiom
\textup{(A2)} gives continuity, and Axiom~\textup{(A3)} gives
non-triviality.  The mixture-space theorem therefore yields a nonconstant
expected-utility index.  Conversely, expected utility is unaffected by record
or action processing and satisfies the three indifference clauses.  The
positive and negative formulae satisfy their respective clauses by data
processing and the two chain rules.
\end{proof}

\subsection{Independence of the axioms}\label{app:axiom-independence}

\begin{proof}[Proof of Proposition~\ref{prop:independence}]
Fix nonconstant $u$ and write $U(q,K):=\E_{qK}[u(O)]$, $Z:=(O,R)$, and
$I(q,K):=\I_{qK}(A;Z)$.  All scores below use the same formula in block and
compound environments.  Both coordinates are block invariant and processing
monotone; under a common-branch replacement their differences are multiplied
by the reached-branch probability.  Moreover, $U$ distinguishes the material
test and $I$ takes values $\log2$ and zero in the trace test.  For witnesses
built only from $(U,I)$, these facts verify every unmentioned axiom below.

\emph{Failure of Axiom \textup{(A1)}.}
Use the Pareto relation on the two coordinates $(U,I)$:
\[
 (q,K)\succeq(p,L)
 \quad\Longleftrightarrow\quad
 U(q,K)\ge U(p,L)\ \text{and}\ I(q,K)\ge I(p,L).
\]
It is transitive with closed graph, but incomplete whenever payoff and
information rank two alternatives oppositely.

\emph{Failure of Axiom \textup{(A2)}.}
Order alternatives lexicographically by $(U,I)$, with payoff first.  This is
a weak order but is not continuous even within one fixed channel.  Let a
three-action channel reveal the action fully and let only action $2$ pay the high outcome.  Put
\[
 p=(1/4,1/2,1/4),\qquad q=(0.49,0.50,0.01),\qquad
 q_n=(0.49-\varepsilon_n,0.50+\varepsilon_n,0.01),
\]
where $\varepsilon_n\downarrow0$.  The lotteries $p$ and $q$ have the same expected utility but
$\Sh(q)<\Sh(p)$, so $p\succ_Kq$.  Every $q_n$ has strictly higher expected utility than $p$, so
$q_n\succ_Kp$.  Thus the graph in Axiom~\textup{(A2)} is not closed.

\emph{Failure of Axioms \textup{(A3)} and \textup{(A4)}.}
For Axiom~\textup{(A3)}, use $V:=I$: pure lotteries carry no information,
although the trace test has values $\log2$ and zero.  For
Axiom~\textup{(A4)}, use $V:=U$: suitable outcomes have different utility,
but the two trace-test lotteries have the same payoff.

\emph{Failure of Axiom \textup{(A5)}.}
The idea is to add a linear-in-the-prior bonus that vanishes in every canonical
two-block comparison but becomes active in a three-block environment.
Define the weights directly from the literal labels of each finite action alphabet $A$.  If every element of
$A$ is literally a pair $(x,i)$ and the set of outermost second coordinates is
$\{0,\ldots,m-1\}$ for some $m\ge3$, put
\[
 w_A(x,i):=\1_{\{i=0\}}.
\]
Otherwise put $w_A=0$.  In a nested tagged sum this convention reads only the literal outermost coordinate, so
it is a single-valued function of the actual action set and not of its construction history.  For
$\varepsilon>0$, use within each environment
\[
 V_\varepsilon(q,K):=U(q,K)+I(q,K)+\varepsilon\sum_aq(a)w_A(a).
\]
This is a continuous weak order.  The material- and trace-relevance tests can use action sets on which
$w_A=0$, so Axioms~\textup{(A3)} and~\textup{(A4)} retain their usual witnesses.
Every comparison in Axioms~\textup{(A6)--(A8)} is evaluated in a canonical outer two-block environment;
there the outermost tag set is $\{0,1\}$, even when a component alphabet is itself tagged.  Hence $w_A=0$
throughout those comparisons and the score reduces to $U+I$.  For three identical one-action,
constant-payoff, silent channels, however, blocks $0$ and $1$ are indifferent
in their canonical two-block comparison but differ by $\varepsilon$ in the
three-block environment.  Thus Axiom~\textup{(A5)} fails.

\emph{Failure of Axiom \textup{(A6)}.}
For $\varepsilon>0$, use
\begin{equation}
 V(q,K):=U(q,K)+I(q,K)-\varepsilon\Sh_{qK}(Z\mid A),
 \label{eq:no-record-dp}
\end{equation}
where $\Sh_{qK}(Z\mid A)$ is the conditional entropy of $Z$ given $A$ under the
joint law $q(a)K(z\mid a)$; the subscript is dropped when the law is clear.
Conditional entropy is continuous and block invariant, and mutual information
and conditional entropy both obey the branch chain rule.  Under action processing $A\to B$, the
original-minus-processed difference is
\[
 \bigl[\I(A;Z)-\I(B;Z)\bigr]
 +\varepsilon\bigl[\Sh(Z\mid B)-\Sh(Z\mid A)\bigr]
 =(1+\varepsilon)\I(A;Z\mid B)\ge0,
\]
so Axiom~\textup{(A7)} holds.  In the trace test the two pairs
$(I(q,K),\Sh(Z\mid A))$ are $(\log2,0)$ and $(0,\log2)$.  But with constant
payoff and a fair record independent of $A$, the score is
$-\varepsilon\log2$; deleting that noise raises it to zero, violating
Axiom~\textup{(A6)}.

\emph{Failure of Axiom \textup{(A7)}.}
Let $G(q):=1-\sum_aq(a)^2$ be Gini entropy and define its posterior reduction
\[
 J_G(q,K):=G(q)-\int G(\pi)\,d\mu_{qK}(\pi),
 \qquad V:=U+J_G.
\]
The score $J_G$ is continuous and block invariant.  Concavity of $G$ gives
record processing, and the potential-difference form gives the branch chain
rule; the trace-test values are $1/2$ and zero.
Action processing can nevertheless increase $J_G$.  Take
$q=(1/8,1/8,3/4)$, a constant payoff, and a binary record with
\[
 \Pr(0\mid a_1)=\Pr(0\mid a_2)=0,
 \qquad \Pr(0\mid a_3)=1/4,
\]
and merge $a_1,a_2$.  Direct calculation gives $J_G=9/416$ before the merge
and $J_G=3/104$ afterwards, violating Axiom~\textup{(A7)}.

\emph{Failure of Axiom \textup{(A8)}.}
The idea is that a cap preserves processing monotonicity but makes a branch's
value depend on how much information has already accumulated elsewhere.
Use $V(q,K):=U(q,K)+\min\{I(q,K),1\}$ in bits.  The cap preserves continuity, block
coherence, processing, and both relevance tests.  Let $q$ be uniform on four
actions and let a first-stage record reveal which of two pairs contains the
action.  Keep payoffs constant and one continuation silent; after the other
record, compare a silent continuation with one that reveals the action within its pair.  The
latter is strictly preferred at the branch posterior, but the silent and revealing compounds have
one and three-halves bits, respectively; both are capped at one and are therefore indifferent.
Taking $K$ silent and $L$ revealing in
\eqref{eq:recordwise-sure-thing}, the compound comparison is true by
indifference whereas $(q_1,K)\not\succeq(q_1,L)$; the biconditional fails.
For an arbitrary cap $c>0$, the same construction serves.  If $c\le1$, both
compound values already exceed the cap, while the branch comparison,
$\min\{1,c\}=c$ against $0$, remains strict.  If $c>1$, take $q$ uniform on
$2^n$ actions with $n-1\ge c$ and let the first stage reveal which of the
$2^{n-1}$ pairs contains the action: the first stage alone reaches the cap, so
both compounds are capped at $c$, and the within-pair branch comparison of one
bit against zero is unchanged.
\end{proof}

\subsection{Choice implications}\label{app:choice-proofs}

The next four blocks mirror Section~\ref{sec:implications}.  Write
$Z:=O\times R$; when $K:A\to\D(Z)$ is fixed, put
$K_a:=K(\cdot\mid a)$ and
$\bar u(a):=\E_{K_a}[u(O)]$; its output law under $q$ is $p_{qK}$.

\subsubsection{The trace test}\label{app:trace-test-proof}

\begin{proof}[Proof of Proposition~\ref{prop:signature}]
(i) Under either lottery, each visible pair $(o_0,r_j)$, $j=1,2$, has probability $\tfrac12$: under $q'$ because
each $a_j$ produces $(o_0,r_j)$ for sure, under $q''$ because each of $a_3,a_4$ produces either pair with equal
probability.  Hence $p_{q'K}=p_{q''K}$.  On $\supp(q')$ the rows are degenerate and distinct, so the visible pair
reveals the action and $\I_{q'K}=\Sh(q')=\log2$; on $\supp(q'')$ the rows are identical, so the visible pair is
independent of the action and $\I_{q''K}=0$.  The payoff terms agree because the payoff marginals do, giving the
displayed gap.
(ii) All rows of $K^\circ$ are equal, so $\I_{qK^\circ}=0$ for every $q$; since the payoff is constantly $o_0$,
the representation gives indifference.
\end{proof}

\subsubsection{Deliberate randomisation}
\label{app:randomisation-proofs}

\begin{proposition}[Optimal mixing]\label{prop:binary}
Let $\lambda>0$ and $A=\{a,b\}$ with $K_a\ne K_b$ and
$\Delta:=\bar u(a)-\bar u(b)\ge0$.  For $x\in[0,1]$ let $q_x$ put mass $x$ on $b$, so that
$p_{q_xK}=(1-x)K_a+xK_b=:p_x$, and for $x\in(0,1)$ define
\[
 g(x):=\KL\!\left(K_b\,\middle\|\,p_x\right)-\KL\!\left(K_a\,\middle\|\,p_x\right).
\]
\begin{enumerate}[label=\textup{(\roman*)},leftmargin=*]
\item $g$ is continuous and strictly decreasing, with
$g(0^+)=\KL(K_b\,\|\,K_a)$ and
$g(1^-)=-\KL(K_a\,\|\,K_b)$ as extended-real limits.  The pure lottery $\delta_a$ is optimal if and only if
$\Delta\ge\lambda\KL(K_b\,\|\,K_a)$.  Below this threshold the optimum is unique and interior,
$q^*=q_{x^*}$ with $x^*$ the unique solution of $\Delta=\lambda g(x)$.
\item The optimal mass $x^*(\Delta,\lambda)$ on the inferior action is nonincreasing in $\Delta$ and
nondecreasing in $\lambda$.  While the optimum is interior, it is strictly decreasing in $\Delta$ and,
when $\Delta>0$, strictly increasing in $\lambda$.  At $\Delta=0$ it is the capacity-achieving input
of the two-row channel, independently of $\lambda$.
\end{enumerate}
\end{proposition}

\begin{proof}
Write
\[
 f(x):=V(q_x,K)=\bar u(a)-x\Delta+\lambda I(x),
 \qquad
 I(x):=\Sh(p_x)-(1-x)\Sh(K_a)-x\Sh(K_b).
\]
For $x\in(0,1)$, differentiation gives $I'(x)=g(x)$ and
\[
 g'(x)
 =-c_{\log}\sum_{z\in Z}\frac{(K_b(z)-K_a(z))^2}{p_x(z)}<0,
\]
where $c_{\log}=1/\ln b_0>0$ for the fixed logarithm base $b_0$
and zero summands are omitted.  Thus $f'(x)=-\Delta+\lambda g(x)$ is strictly
decreasing.  As $x\downarrow0$, finite summands converge to those in
$\KL(K_b\|K_a)$, while any $z$ with $K_b(z)>0=K_a(z)$ makes the limit
$+\infty$; exchanging $a$ and $b$ gives the limit at one.  Consequently
$x=0$ is optimal exactly when
$\Delta\ge\lambda\KL(K_b\|K_a)$.  Otherwise $f'(0^+)>0$ and
$f'(1^-)=-\Delta-\lambda\KL(K_a\|K_b)<0$, so the unique optimum is the
root $\Delta=\lambda g(x)$.  At an interior optimum,
$x^*=g^{-1}(\Delta/\lambda)$; since $g^{-1}$ is strictly decreasing, this
gives strict decrease in $\Delta$ and, when $\Delta>0$, strict increase in
$\lambda$, together with the stated weak comparative statics at the boundary.
At $\Delta=0$, $g(x)=0$ is the usual capacity condition and is independent
of $\lambda$.
\end{proof}

\paragraph{General optimality.}
For every $q$, the entropy identity gives
\begin{equation}\label{eq:choice-value}
 \begin{aligned}
 V(q,K)
 &=\lambda\Sh(p_{qK})
   +\sum_aq(a)\bigl[\bar u(a)-\lambda\Sh(K_a)\bigr] \\
 &=\sum_{a\in\supp(q)}q(a)
   \left[\bar u(a)+\lambda
   \KL\!\left(K_a\,\middle\|\,p_{qK}\right)\right].
 \end{aligned}
\end{equation}
Define
\[
 \Phi_a(p):=\bar u(a)+\lambda\KL\!\left(K_a\,\middle\|\,p\right).
\]
For $q_t:=(1-t)q+t\delta_a$, direct differentiation in the extended-real
sense gives
\begin{equation}\label{eq:choice-direction}
 \left.\frac{d}{dt}V(q_t,K)\right|_{t=0^+}
 =\Phi_a(p_{qK})-V(q,K).
\end{equation}

\begin{proof}[Proof of Proposition~\ref{prop:optimal-choice}]
Expected utility is linear in $q$, while mutual information is concave in
the input distribution \cite{CoverThomas2006}; continuity on the compact
simplex gives existence.  At an optimum $q^*$, every vertex direction in
\eqref{eq:choice-direction} is nonpositive, so
$\Phi_a(p_{q^*K})\le V(q^*,K)$ for all $a$.  Equation
\eqref{eq:choice-value} says their $q^*$-weighted average is $V(q^*,K)$;
hence equality holds on $\supp(q^*)$.

Conversely, the proposition's equalities and inequalities, averaged under
$q^*$, give $c=V(q^*,K)$.  The derivative toward any lottery $r$ is
\[
 \sum_a r(a)\bigl[\Phi_a(p_{q^*K})-c\bigr]\le0.
\]
Concavity makes this condition sufficient for global optimality.  At
constant $u$ it is the classical capacity condition.
\end{proof}

\begin{proposition}[Support and pure choice]\label{cor:support-pure}
Let $\lambda>0$ and let
$Z_K:=\{z\in Z:K_a(z)>0\text{ for some }a\in A\}$.
\begin{enumerate}[label=\textup{(\roman*)},leftmargin=*]
\item Every optimal lottery $q^*$ satisfies $p_{q^*K}(z)>0$ for every
$z\in Z_K$.  Consequently, if every action $a$ has a private
consequence---some $z_a$ with $K_a(z_a)>0$ and $K_b(z_a)=0$ for all
$b\ne a$---then every optimal lottery has full support.
\item The pure lottery $\delta_a$ is optimal if and only if
\[
 \bar u(b)+\lambda\KL\!\left(K_b\,\middle\|\,K_a\right)
 \le \bar u(a)
 \qquad\text{for every }b\in A,
\]
and uniquely optimal if these inequalities are strict for $b\ne a$.  In
particular, if some rival reaches a consequence outside the support of
$K_a$, then $\delta_a$ is not optimal, however large its payoff advantage.
\end{enumerate}
\end{proposition}

\begin{proof}
(i) Put $p^*:=p_{q^*K}$.  If $p^*(z)=0<K_a(z)$, then
$\KL(K_a\|p^*)=+\infty$, contradicting either the on-support equality or
the off-support inequality in Proposition~\ref{prop:optimal-choice}.
If action $a$ has a private consequence but $q^*(a)=0$, that consequence
has zero probability under $p^*$, contradicting the first claim.

(ii) At $q=\delta_a$ the output law is $K_a$ and the on-support value is
$\bar u(a)$.  Proposition~\ref{prop:optimal-choice} therefore gives exactly
the displayed inequalities.  If they are strict off $a$, every nonzero
direction from $\delta_a$ has strictly negative directional derivative, so
concavity makes $\delta_a$ the unique maximiser.
If $K_b$ is not supported by $K_a$, then $\KL(K_b\|K_a)=+\infty$, so the
displayed condition fails.
\end{proof}

\begin{proof}[Proof of Corollary~\ref{rem:logit}]
Pairwise-disjoint row supports make the visible consequence identify the
action, so $\I_{qK}(A;Z)=\Sh(q)$; they also give every action a private
consequence, so Proposition~\ref{cor:support-pure}\textup{(i)} gives full
support at the optimum.  On $\supp K_a$,
$p_{qK}(z)=q(a)K_a(z)$.  Hence
$\KL(K_a\|p_{qK})=-\log q(a)$.  The equalities in
Proposition~\ref{prop:optimal-choice} give
$\bar u(a)-\lambda\log q^*(a)=c$.  In natural-log units, solving and
normalising yields the displayed logit formula.  Strict concavity of
$\sum_aq(a)\bar u(a)+\lambda\Sh(q)$ gives uniqueness; changing the
logarithm base merely rescales $\lambda$.
\end{proof}

\begin{proposition}[Optimal marginal and multiplicity]\label{prop:optimal-marginal}
Let $\lambda>0$ and write $C^*(K):=\max_{q\in\D(A)}V(q,K)$.  There is a unique $p^*\in\D(Z)$ such that
$p_{q^*K}=p^*$ for every optimal lottery $q^*$.  With
\[
 \mathcal A^*
 :=\bigl\{a\in A:
 \bar u(a)+\lambda
 \KL\!\left(K_a\,\middle\|\,p^*\right)=C^*(K)
 \bigr\},
\]
the set of optimal lotteries is exactly
\[
 \bigl\{q\in\D(A):
 p_{qK}=p^*,\ \supp(q)\subseteq\mathcal A^*
 \bigr\}.
\]
The optimum is unique whenever the rows $\{K_a:a\in\mathcal A^*\}$ are affinely independent.  If actions are
partitioned into classes with identical rows, $V(q,K)$ depends on $q$ within each class only through the class's
total mass.  The optimal lotteries are exactly the lifts of the optimal lotteries of the collapsed channel, and
each class's mass may be distributed arbitrarily among its members.
\end{proposition}

\begin{proof}[Proof of Proposition~\ref{prop:optimal-marginal}]
If optimal $q_0,q_1$ induced different output laws, the first line of
\eqref{eq:choice-value} would be strictly concave along their segment while
its remaining term is linear, producing a value above $C^*(K)$.  Thus every
optimum induces the same $p^*$.  Proposition~\ref{prop:optimal-choice} and
\eqref{eq:choice-value} then show that every active score equals $C^*(K)$,
so $\supp(q^*)\subseteq\mathcal A^*$.  Conversely, if $p_{qK}=p^*$ and
$\supp(q)\subseteq\mathcal A^*$, then
$V(q,K)=\sum_aq(a)\Phi_a(p^*)=C^*(K)$.
Affine independence makes the representation $p^*=\sum_aq(a)K_a$ with $\supp(q)\subseteq\mathcal A^*$ unique.
Identical rows share $\bar u(a)$ and $\Sh(K_a)$, so both $p_{qK}$ and
\eqref{eq:choice-value} are
unchanged by redistribution within a class; collapsing the classes and lifting aggregate masses proves the last
claim.
\end{proof}

\begin{remark}[Overlapping traces and IIA]\label{rem:overlap-iia-example}
Let three actions have equal expected utility and rows
\[
 K_1=(1,0,0),\qquad K_2=(0,1,0),\qquad K_3=(1/2,0,1/2).
\]
With only the first two actions, the optimal lottery is $(1/2,1/2)$.  Adding the third gives
\[
 q^*=(1/3,4/9,2/9),\qquad p_{q^*K}=(4/9,4/9,1/9).
\]
Each active row has divergence $\log(9/4)$ from $p_{q^*K}$, so
Proposition~\ref{prop:optimal-choice} verifies optimality.  The ratio of the first two choice probabilities falls
from $1$ to $3/4$.  Thus overlapping rows can generate the red-bus/blue-bus pattern even when expected utilities
are equal.
\end{remark}

\subsubsection{The value of being on record}
\label{app:record-value-proofs}

\begin{proposition}[Cost of a record garbling]\label{prop:garbling-price}
Let $K:A\to\D(O\times R)$, let $T$ be a record processor, and write
$L:=KT$.  Fix $q\in\D(A)$ and let $(A,O,R,R')$ carry the joint law $q(a)K(o,r\mid a)T(o,r'\mid o,r)$.  Then
\[
 V(q,K)-V(q,L)=\lambda\,\I\!\left(A;R\mid O,R'\right)\ge0,
\]
with equality if and only if, after $(O,R')$ has been observed, the original record $R$ contains no further
information about the action.
\end{proposition}

\begin{proof}[Proof of Proposition~\ref{prop:garbling-price}]
The processor leaves the payoff coordinate unchanged, so the expected-utility terms coincide and the difference is
$\lambda[\I_{qK}(A;O,R)-\I_{qL}(A;O,R')]$.  Within this proof, write $Z$ for the random pair $(O,R)$ and $Z'$ for
$(O,R')$; posteriors evaluated at the realised pair are written
$\pi^{qK}_{Z}$ and $\pi^{qL}_{Z'}$.  By construction $A\to Z\to Z'$ is
a Markov chain, so expanding $\I(A;Z,Z')$ in both orders gives
\[
 \I(A;Z)-\I(A;Z')=\I(A;Z\mid Z')=\I(A;R\mid O,R'),
\]
the last equality because $O$ is part of $Z'$.  Conditional mutual information is nonnegative and vanishes exactly
when $A$ and $(O,R)$ are conditionally independent given $(O,R')$ \cite{CoverThomas2006}.  This is the zero-price
condition stated in the proposition.  Under a garbling, it is also equivalent to equality of the posterior laws,
$\mu_{qK}=\mu_{qL}$.  Indeed, the Markov property gives
$\Pr(A=\cdot\mid Z,Z')=\pi^{qK}_{Z}$ almost surely, while conditional independence gives
$\Pr(A=\cdot\mid Z,Z')=\pi^{qL}_{Z'}$; the two random posteriors then coincide almost surely.  Conversely, by the
identity in the proof of Corollary~\ref{cor:matching}, equality of the posterior laws gives equality of the
information terms, and the payoff terms already agree.
\end{proof}

\begin{corollary}[The recorded coin]\label{cor:recorded-coin}
Let $K,L:A\to\D(O\times R)$ and $t\in(0,1)$.  Let $tK\oplus(1-t)L$ be the public mixture that draws an
action-independent coin $Y$ with $\Pr(Y=0)=t$, applies $K$ if $Y=0$ and $L$ otherwise, and includes the coin in
the record, so the full public-mixture record is $(Y,R)$; let $M_t:=tK+(1-t)L$ be the rowwise mixture, which hides it.  Then, for every $q\in\D(A)$,
\begin{enumerate}[label=\textup{(\roman*)},leftmargin=*]
\item $V\bigl(q,\,tK\oplus(1-t)L\bigr)=t\,V(q,K)+(1-t)\,V(q,L)$;
\item $V\bigl(q,\,tK\oplus(1-t)L\bigr)-V(q,M_t)=\lambda\,\I(A;Y\mid O,R)\ge0$, with equality if and only if $A$
and $Y$ are conditionally independent given $(O,R)$.
\end{enumerate}
\end{corollary}

\begin{proof}[Proof of Corollary~\ref{cor:recorded-coin}]
Part \textup{(i)} is the chain rule with an action-independent first stage:
$\I(A;Y)=0$, and each reached branch reproduces the joint law of the
corresponding constituent.  Part \textup{(ii)} follows from
Proposition~\ref{prop:garbling-price}: deleting the coin coordinate is a
deterministic payoff-preserving record processor carrying the public mixture
to $M_t$.
\end{proof}

\begin{proposition}[Uniform dominance is more-capable, not garbling]
\label{prop:blackwell}\label{def:channel-orders}
Let $K:A\to\D(O\times R)$ and $L:A\to\D(O\times R')$ have identical
payoff kernels.  Write $K\succeq_gL$ if $L=KT$ for some record processor
$T$, and write $K\succeq_cL$ if
$\I_{qK}(A;O,R)\ge\I_{qL}(A;O,R')$ for every $q\in\D(A)$.  The latter is
the \emph{more-capable} order \cite{KornerMarton1977}.
\begin{enumerate}[label=\textup{(\roman*)},leftmargin=*]
\item $K\succeq_cL$ if and only if $V(q,K)\ge V(q,L)$ for every $q\in\D(A)$.
\item $K\succeq_gL\Rightarrow K\succeq_cL$, but the converse fails.  There are channels for which
$K\succeq_cL$ although $L\ne KT$ for every record processor $T$.
\end{enumerate}
\end{proposition}

\begin{proof}[Proof of Proposition~\ref{prop:blackwell}]
(i) Because the payoff kernels coincide,
$V(q,K)-V(q,L)=\lambda[\I_{qK}-\I_{qL}]$ with $\lambda>0$.  Uniform value dominance and
$K\succeq_cL$ are therefore the same statement.

(ii) If $L=KT$, Proposition~\ref{prop:garbling-price} gives $\I_{qK}\ge\I_{qL}$ at every $q$, and hence
$K\succeq_gL\Rightarrow K\succeq_cL$.  For failure of the converse, let $A=\{0,1\}$, fix $\bar o\in O$, and let
both channels pay $\bar o$ surely.  Let $K$ carry the binary
erasure record with erasure probability $\varepsilon=\tfrac12$ and $L$ the binary symmetric record with crossover
$p=\tfrac15$:
\[
 K(\bar o,a\mid a)=1-\varepsilon,\quad K(\bar o,\mathrm e\mid a)=\varepsilon,
 \qquad
 L(\bar o,a\mid a)=1-p,\quad L(\bar o,1-a\mid a)=p.
\]
The payoff kernels agree row by row, so it suffices to compare informations, and the comparison is invariant to the
base; compute in bits with $h$ the binary entropy.  For $q=(1-t,t)$, grouping records gives
\[
 \I_{qK}=(1-\varepsilon)\,h(t),
 \qquad
 \I_{qL}=h(p\ast t)-h(p),
 \qquad p\ast t:=p(1-t)+(1-p)t.
\]
By Mrs Gerber's Lemma \cite{WynerZiv1973}, $v\mapsto h\bigl(p\ast h^{-1}(v)\bigr)$ is convex on $[0,1]$, with
$h^{-1}$ valued in $[0,\tfrac12]$; since $h(p\ast t)$ is invariant under $t\mapsto1-t$, take $t\le\tfrac12$ and
$v=h(t)$.  The chord between $v=0$ and $v=1$ gives
\[
 h(p\ast t)\le(1-v)\,h(p)+v,
 \qquad\text{hence}\qquad
 \I_{qL}\le\bigl(1-h(p)\bigr)h(t)\le(1-\varepsilon)\,h(t)=\I_{qK},
\]
the last step because $\varepsilon=\tfrac12\le h(\tfrac15)$.  Thus $K\succeq_cL$.  Now suppose $L=KT$ for a
processor $T$, and write $t_r:=T(\bar o,0\mid \bar o,r)$ for $r\in\{0,1,\mathrm e\}$.  Matching the probability of the
processed record $0$ in the two rows,
\[
 (1-\varepsilon)t_0+\varepsilon t_{\mathrm e}=1-p,
 \qquad
 (1-\varepsilon)t_1+\varepsilon t_{\mathrm e}=p,
\]
and subtracting gives $t_0-t_1=(1-2p)/(1-\varepsilon)=\tfrac65>1$, impossible.  Compositions of record processors
are record processors, so no chain of garblings produces $L$ from $K$ either.
\end{proof}

\subsubsection{Measuring the trace weight}
\label{app:measurement-proofs}

\begin{proposition}[Single information price]
\label{prop:single-price}\label{lem:info-payoff-plane}
If $(q,K)\sim(p,L)$ and $\I_{qK}\ne\I_{pL}$, then
\[
 \frac{E_{pL}-E_{qK}}{\I_{qK}-\I_{pL}}=\lambda,
\]
and $\lambda$ is the only number satisfying this identity at every such
indifference.  More generally, writing
\[
 \phi(q,K):=\bigl(\I_{qK}(A;O,R),E_{qK}\bigr),\qquad
 \underline u:=\min_o u(o),\quad \overline u:=\max_o u(o),
\]
the image of $\phi$ over all finite pairs is
$[0,\infty)\times[\underline u,\overline u]$.  In these coordinates the
block-comparison indifference classes are the nonempty intersections of the
lines $E=c-\lambda\I$ with that strip.
\end{proposition}

\begin{proof}
The ratio identity rearranges
$E_{qK}+\lambda\I_{qK}=E_{pL}+\lambda\I_{pL}$.
For the image claim, only the reverse inclusion needs proof.  Given
$(t,w)$ in the strip, choose $n$ with $\log n\ge t$ and, by continuity of
entropy between a point mass and the uniform lottery, choose $q_t$ on
$\{1,\ldots,n\}$ with $\Sh(q_t)=t$.  Choose $o^+,o^-\in O$ with
$u(o^+)=\overline u$ and $u(o^-)=\underline u$.  Let a channel reveal the
action in its record and, independently of the action, pay $o^+$ with
probability $\beta$ and $o^-$ otherwise, where
$\beta u(o^+)+(1-\beta)u(o^-)=w$.  It realises
$\phi(q_t,K)=(t,w)$, and Theorem~\ref{thm:main} gives the stated lines.

For uniqueness, if $\lambda'\ne\lambda$, put
$t_0:=\tfrac12\min\{1,(\overline u-\underline u)/\lambda\}>0$.
The image claim supplies pairs with coordinates
$(t_0,\overline u-\lambda t_0)$ and $(0,\overline u)$.  They are indifferent,
whereas their $E+\lambda'\I$ values differ by
$(\lambda'-\lambda)t_0\ne0$.
\end{proof}

\begin{proposition}[Willingness to pay for trace]
\label{prop:wtp}\label{def:dilution}
Let $o^+,o^-$ attain the maximum and minimum of $u$, and put
\[
 \Delta u:=u(o^+)-u(o^-)>0,
 \qquad
 w(\beta):=\beta u(o^+)+(1-\beta)u(o^-).
\]
For a channel $N:C\to\D(O\times Y)$, zero-pad it to
$Y^\ast:=Y\sqcup\{\ast\}$, giving the added record probability zero under $N$.  For $\varepsilon\in[0,1)$ and
$\beta\in[0,1]$, let $P_\varepsilon:C\to\D(\{1,2\})$ satisfy
$P_\varepsilon(2\mid c)=\varepsilon$ for every $c$, and let $J_\beta$ have
constant rows
\[
 J_\beta(o^+,\ast\mid c)=\beta,
 \qquad J_\beta(o^-,\ast\mid c)=1-\beta.
\]
Write
$D_{\varepsilon,\beta}N:=P_\varepsilon\triangleright(N,J_\beta)$.
Fix $q\in\D(A)$, $p\in\D(B)$, and channels $K:A\to\D(O\times R)$ and $L:B\to\D(O\times S)$.
\begin{enumerate}[label=\textup{(\roman*)},leftmargin=*]
\item For every $\varepsilon\in[0,1)$ and $\beta\in[0,1]$,
\[
 V\bigl(q,\,D_{\varepsilon,\beta}K\bigr)=(1-\varepsilon)\,V(q,K)+\varepsilon\,w(\beta).
\]
\item If $\varepsilon\in(0,1)$ and $\beta,\beta'\in[0,1]$ satisfy
$\bigl(q,D_{\varepsilon,\beta}K\bigr)\sim\bigl(p,D_{\varepsilon,\beta'}L\bigr)$, then
\[
 \frac{\varepsilon}{1-\varepsilon}\bigl[w(\beta)-w(\beta')\bigr]
 =V(p,L)-V(q,K).
\]
When the material parts agree, $E_{qK}=E_{pL}$, the left-hand side equals
$\lambda\bigl[\I_{pL}(B;O,S)-\I_{qK}(A;O,R)\bigr]$: information is priced linearly at rate $\lambda$.
\item Such a pair $(\beta,\beta')$ exists if and only if
$\lvert V(p,L)-V(q,K)\rvert\le\frac{\varepsilon}{1-\varepsilon}\Delta u$; a boundary response therefore brackets
the value difference instead of measuring it.
\end{enumerate}
\end{proposition}

\begin{proof}
(i) The material part of the compound is $(1-\varepsilon)\E_{qK}[u(O)]+\varepsilon w(\beta)$.  For the trace part,
the coin is action-independent, so both branch posteriors equal $q$ and the chain rule for the compound gives an
information term of $(1-\varepsilon)\I_{qK}(A;O,R)$, the side-bet branch contributing zero.
(ii) By the moreover clause of Theorem~\ref{thm:main}, the indifference holds if and only if
$(1-\varepsilon)V(q,K)+\varepsilon w(\beta)=(1-\varepsilon)V(p,L)+\varepsilon w(\beta')$; rearrange, and
substitute the representation when the material parts cancel.
(iii) $w(\beta)-w(\beta')$ is continuous on $[0,1]^2$ with range $[-\Delta u,\Delta u]$; apply the intermediate
value theorem to the equation in \textup{(ii)}.
\end{proof}

\begin{proof}[Proof of Proposition~\ref{prop:elicit}]
Write $K^{\mathrm{id}}$ and $K^{0}$ for the revealing and the silent channel at the constant payoff, so that
$V(q,K^{\mathrm{id}})-V(q,K^{0})=\lambda\Sh(q)$.  By Proposition~\ref{prop:wtp}\textup{(i)}, the two diluted
values are $(1-\varepsilon)V(q,K^{\mathrm{id}})+\varepsilon w(0)$ and
$(1-\varepsilon)V(q,K^{0})+\varepsilon w(\beta)$, and under the normalisation $w(\beta)=\beta$.  Indifference
therefore holds if and only if $\varepsilon\beta=(1-\varepsilon)\lambda\Sh(q)$.  The left side is continuous and
strictly increasing in $\beta$ with range $[0,\varepsilon]$, so a solution exists, and is then unique, exactly
when $\lambda\Sh(q)\le\varepsilon/(1-\varepsilon)$; when this fails, the revealing side's value exceeds the
silent side's at every $\beta$.
\end{proof}

\begin{remark}[Implementation and controls]\label{rem:elicitation}
Normalise $u(o^+)=1$ and $u(o^-)=0$.  All comparisons below are primitive
block comparisons.  For each $o\in O$, elicit $\beta_o$ from
$(\delta_\ast,K_o)\sim(\delta_\ast,J_{\beta_o})$, where $K_o$ pays $o$
surely; both one-action channels carry zero information, so
$u(o)=\beta_o$.  Next take the uniform $q$ on $n\ge2$ actions and the
constant-payoff revealing and silent channels $K_o^{\mathrm{id}}$ and
$K_o^0$, and elicit $\beta$ from
\[
 (q,D_{\varepsilon,\beta}K_o^0)
 \sim(q,D_{\varepsilon,0}K_o^{\mathrm{id}}),
 \qquad
 \lambda=\frac{\varepsilon}{1-\varepsilon}\frac{\beta}{\log n}.
\]
If even $\beta=1$ does not produce indifference, the response gives the lower bound
$\lambda>\varepsilon/((1-\varepsilon)\log n)$, and $\varepsilon$ is raised.

Three controls are essential.  \emph{Orientation:} revelation must beat
silence at every constant payoff.  \emph{Garbling:} in a separate comparison of the undiluted,
equal-payoff channels $K_o^{\mathrm{id}}$ and $K_o^0$, total record
erasure must make the two blocks indifferent; otherwise some display feature
still traces the action.  \emph{Transport:} changing the prior, payoff, or
action-set size must return the same $\lambda$ by
Proposition~\ref{prop:single-price}; a decline with $\log n$ instead signals
the capped-information criterion.
\end{remark}

\clearpage
\begin{thebibliography}{99}
\raggedright

\bibitem{AgranovOrtoleva2017}
Marina Agranov and Pietro Ortoleva.
\newblock Stochastic choice and preferences for randomization.
\newblock {\em Journal of Political Economy}, 125(1):40--68, 2017.

\bibitem{AndreoniBernheim2009}
James Andreoni and B. Douglas Bernheim.
\newblock Social image and the 50--50 norm: A theoretical and experimental analysis of audience effects.
\newblock {\em Econometrica}, 77(5):1607--1636, 2009.

\bibitem{AzrieliLehrer2008}
Yaron Azrieli and Ehud Lehrer.
\newblock The value of a stochastic information structure.
\newblock {\em Games and Economic Behavior}, 63(2):679--693, 2008.

\bibitem{BartlingFehrHerz2014}
Bj\"orn Bartling, Ernst Fehr, and Holger Herz.
\newblock The intrinsic value of decision rights.
\newblock {\em Econometrica}, 82(6):2005--2039, 2014.

\bibitem{BaezFritzLeinster2011}
John C. Baez, Tobias Fritz, and Tom Leinster.
\newblock A characterization of entropy in terms of information loss.
\newblock {\em Entropy}, 13(11):1945--1957, 2011.

\bibitem{BenabouTirole2002}
Roland B\'enabou and Jean Tirole.
\newblock Self-confidence and personal motivation.
\newblock {\em Quarterly Journal of Economics}, 117(3):871--915, 2002.

\bibitem{BenabouTirole2006}
Roland B\'enabou and Jean Tirole.
\newblock Incentives and prosocial behavior.
\newblock {\em American Economic Review}, 96(5):1652--1678, 2006.

\bibitem{Bernheim1994}
B. Douglas Bernheim.
\newblock A theory of conformity.
\newblock {\em Journal of Political Economy}, 102(5):841--877, 1994.

\bibitem{Blackwell1953}
David Blackwell.
\newblock Equivalent comparisons of experiments.
\newblock {\em Annals of Mathematical Statistics}, 24(2):265--272, 1953.


\bibitem{CabralesGossnerSerrano2013}
Antonio Cabrales, Olivier Gossner, and Roberto Serrano.
\newblock Entropy and the value of information for investors.
\newblock {\em American Economic Review}, 103(1):360--377, 2013.

\bibitem{CaplinDeanLeahy2022}
Andrew Caplin, Mark Dean, and John Leahy.
\newblock Rationally inattentive behavior: Characterizing and generalizing {S}hannon entropy.
\newblock {\em Journal of Political Economy}, 130(6):1676--1715, 2022.

\bibitem{CaplinLeahy2001}
Andrew Caplin and John Leahy.
\newblock Psychological expected utility theory and anticipatory feelings.
\newblock {\em Quarterly Journal of Economics}, 116(1):55--79, 2001.

\bibitem{CerreiaVioglioDillenbergerOrtolevaRiella2019}
Simone Cerreia-Vioglio, David Dillenberger, Pietro Ortoleva, and Gil Riella.
\newblock Deliberately stochastic.
\newblock {\em American Economic Review}, 109(7):2425--2445, 2019.

\bibitem{CorraoFudenbergMudekerezaLevine2026}
Roberto Corrao, Drew Fudenberg, Florian Mudekereza, and David Levine.
\newblock Randomization, surprise, and adversarial forecasters.
\newblock Working paper, August 2026.

\bibitem{CoverThomas2006}
Thomas M. Cover and Joy A. Thomas.
\newblock {\em Elements of Information Theory}.
\newblock Wiley-Interscience, Hoboken, second edition, 2006.

\bibitem{Denti2022}
Tommaso Denti.
\newblock Posterior separable cost of information.
\newblock {\em American Economic Review}, 112(10):3215--3259, 2022.

\bibitem{Dostoevsky1864}
Fyodor Dostoevsky.
\newblock {\em Notes from the Underground}, Part I.
\newblock 1864.  Translated by Constance Garnett, 1918.

\bibitem{ElyFrankelKamenica2015}
Jeffrey Ely, Alexander Frankel, and Emir Kamenica.
\newblock Suspense and surprise.
\newblock {\em Journal of Political Economy}, 123(1):215--260, 2015.

\bibitem{Faddeev1956}
D. K. Faddeev.
\newblock On the concept of entropy of a finite probabilistic scheme.
\newblock {\em Uspekhi Matematicheskikh Nauk}, 11(1):227--231, 1956.
\newblock In Russian.

\bibitem{FrankelKamenica2019}
Alexander Frankel and Emir Kamenica.
\newblock Quantifying information and uncertainty.
\newblock {\em American Economic Review}, 109(10):3650--3680, 2019.

\bibitem{GilboaLehrer1991}
Itzhak Gilboa and Ehud Lehrer.
\newblock The value of information---an axiomatic approach.
\newblock {\em Journal of Mathematical Economics}, 20(5):443--459, 1991.

\bibitem{GrantKajiiPolak1998}
Simon Grant, Atsushi Kajii, and Ben Polak.
\newblock Intrinsic preference for information.
\newblock {\em Journal of Economic Theory}, 83(2):233--259, 1998.

\bibitem{HersteinMilnor1953}
I. N. Herstein and John Milnor.
\newblock An axiomatic approach to measurable utility.
\newblock {\em Econometrica}, 21(2):291--297, 1953.

\bibitem{Holmstrom1979}
Bengt Holmstr\"om.
\newblock Moral hazard and observability.
\newblock {\em Bell Journal of Economics}, 10(1):74--91, 1979.

\bibitem{Holmstrom1982}
Bengt Holmstr\"om.
\newblock Moral hazard in teams.
\newblock {\em Bell Journal of Economics}, 13(2):324--340, 1982.


\bibitem{KlyubinPolaniNehaniv2005}
Alexander S. Klyubin, Daniel Polani, and Chrystopher L. Nehaniv.
\newblock Empowerment: A universal agent-centric measure of control.
\newblock In {\em Proceedings of the 2005 IEEE Congress on Evolutionary Computation}, volume~1, pages 128--135. IEEE,
  2005.

\bibitem{KornerMarton1977}
J\'anos K\"orner and Katalin Marton.
\newblock Comparison of two noisy channels.
\newblock In Imre Csisz\'ar and Peter Elias, editors, {\em Topics in Information Theory}, Colloquia Mathematica
  Societatis J\'anos Bolyai, volume~16, pages 411--423. North-Holland, Amsterdam, 1977.


\bibitem{MatejkaMcKay2015}
Filip Mat\v{e}jka and Alisdair McKay.
\newblock Rational inattention to discrete choices: A new foundation for the multinomial logit model.
\newblock {\em American Economic Review}, 105(1):272--298, 2015.

\bibitem{MistryLiljeholm2016}
Prachi Mistry and Mimi Liljeholm.
\newblock Instrumental divergence and the value of control.
\newblock {\em Scientific Reports}, 6:36295, 2016.

\bibitem{NortonLiljeholm2020}
Kaitlyn G. Norton and Mimi Liljeholm.
\newblock The rostrolateral prefrontal cortex mediates a preference for high-agency environments.
\newblock {\em Journal of Neuroscience}, 40(22):4401--4409, 2020.


\bibitem{PomattoStrackTamuz2023}
Luciano Pomatto, Philipp Strack, and Omer Tamuz.
\newblock The cost of information: The case of constant marginal costs.
\newblock {\em American Economic Review}, 113(5):1360--1393, 2023.

\bibitem{SalgeGlackinPolani2014}
Christoph Salge, Cornelius Glackin, and Daniel Polani.
\newblock Empowerment---an introduction.
\newblock In Mikhail Prokopenko, editor, {\em Guided Self-Organization: Inception}, pages 67--114. Springer, Berlin,
  2014.

\bibitem{Savage1954}
Leonard J. Savage.
\newblock {\em The Foundations of Statistics}.
\newblock Wiley, New York, 1954.

\bibitem{Shannon1948}
Claude E. Shannon.
\newblock A mathematical theory of communication.
\newblock {\em Bell System Technical Journal}, 27:379--423, 623--656, 1948.

\bibitem{Sims2003}
Christopher A. Sims.
\newblock Implications of rational inattention.
\newblock {\em Journal of Monetary Economics}, 50(3):665--690, 2003.

\bibitem{WynerZiv1973}
Aaron D. Wyner and Jacob Ziv.
\newblock A theorem on the entropy of certain binary sequences and applications: Part {I}.
\newblock {\em IEEE Transactions on Information Theory}, 19(6):769--772, 1973.

\end{thebibliography}


\end{document}
